%% OVERFORE-Formulario-T-14.tex — generado desde OVERFORE-Formulario-T-14.docx
%% Formato: clase overfore, modo formulario. El formato lo pone la clase; aquí sólo va el contenido.
\documentclass[formulario]{../../comun/overfore}
\numsubdoc{7}
\formulario{T-14}{Plan de trabajo, EDT y carta Gantt}
\informe{Informe Preparatorio N° 2}
\fecha{12 de octubre de 2026}
\externos{A1,S1,S3,S4,S5,T9}
\begin{document}

\apartado{Objeto y estructura del formulario}
Este formulario presenta la estructura de descomposición del trabajo (EDT) de SITCO Aranda, su carta Gantt de 56 meses y su diccionario. El proyecto comprende 8 ramas, 55 entregables y 175 paquetes, con los mismos códigos en diagramas, cronograma y fichas. La Tabla \ref{tab:campos-diccionario-edt} define los campos de cada ficha.

\begin{tabla}{I{0.286}I{0.714}}{Campos del diccionario de la EDT.}{tab:campos-diccionario-edt}
  {\cab{Campo} & \cab{Qué se escribe en él}}
  {Fuente: elaboración propia a partir de la Guía del PMBOK (PMI, 2017, p. 162).}
Código & El identificador del código de cuentas. \\
Nombre & El nombre del componente. \\
Descripción del trabajo & Qué trabajo comprende el paquete y qué queda fuera de él. \\
Entregable & El producto o resultado concreto que el paquete produce. \\
Criterio de aceptación & La condición verificable con que se dará por recibido el entregable. \\
Responsable & La persona o la organización que responde por el paquete. \\
Hitos asociados & Los hitos del cronograma que dependen de este paquete. \\
Esfuerzo estimado & Las horas-persona estimadas, desagregadas por perfil. \\
Supuestos & Lo que se dio por cierto al estimar este paquete. \\
Referencias & Los requisitos, contratos o documentos técnicos aplicables. \\
\end{tabla}
Cada ficha comienza con el código y el nombre. La del entregable resume el trabajo, el responsable, los meses, los hitos y las HH. La del paquete añade los demás campos; incluye Supuestos cuando existe uno específico.

Las actividades elementales de trabajo se estiman individualmente en un rango de 8 a 80 HH. Para aquellos paquetes de trabajo complejos o de ejecución continua cuya entrega o unidad indivisible de gestión requiere un esfuerzo superior a este límite (por ejemplo, los 17 paquetes identificados con N = 1), la cifra declarada representa el esfuerzo consolidado del paquete completo, el cual se desagrega operativamente en tareas elementales de 8 a 80 HH dentro de la nivelación de recursos del plan de trabajo. Los paquetes continuos reciben el saldo de horas de su rol. Las ramas 1.1 a 1.7 suman 60.394 HH y la operación, 122.672 HH, sujetas a ajustes mediante planificación gradual (PMI, 2017).

Las HH se estiman con PERT de tres valores: optimista (O), más probable (M) y pesimista (P), con $E = \frac{(O + 4M + P) }{6}$ y $\sigma = \frac{(P - O)}{6}$ (PMI, 2017). Cada ficha indica la clase de incertidumbre de la Tabla \ref{tab:clases-incertidumbre-estimacion-pert-resultado}. O y P se obtienen de M con los factores de esa tabla, tomados del cono de incertidumbre (McConnell, 2006), y M se despeja para que E sea la cifra de la ficha.

\begin{tabla}{I{0.097}I{0.382}I{0.185}D{0.131}D{0.114}D{0.091}}{Clases de incertidumbre de la estimación PERT y su resultado.}{tab:clases-incertidumbre-estimacion-pert-resultado}
  {\cab{Clase} & \cab{Trabajo que incluye} & \cab{O y P, en veces M (aproximado)} & \cab{Paquetes} & \cab{E (HH)} & \cab{σ (HH)}}
  {Fuente: elaboración propia, con los rangos de McConnell (2006) y la estimación de tres valores de la Guía del PMBOK (PMI, 2017, pp. 201–202).}
Baja & Trabajo recurrente o de contenido fijado: informes, actas, actualizaciones, pruebas periódicas, jornadas de capacitación y digitalización por lotes & 0,80 y 1,25 & 50 & 12.368 & 272 \\
Media & Trabajo definido por el catálogo de requerimientos: diseño, construcción, pruebas e instalación & 0,67 y 1,50 & 83 & 22.224 & 514 \\
Alta & Trabajo que depende de terceros o de información que no existe al ofertar: fabricantes, sindicato, proveedores, inventario de los diez sistemas y migración de datos & 0,50 y 2,00 & 21 & 6.568 & 507 \\
Sin PERT & Saldo de la dedicación de los roles clave, cobertura de turnos y capacidad mensual comprometida & No aplica & 21 & 141.906 & No aplica \\
Total & — & — & 175 & 183.066 & 772 \\
\end{tabla}
Los paquetes de incertidumbre alta reúnen el 16 \% de las HH estimadas con PERT y el 43 \% de su varianza. Con 80 \% de confianza, las 41.160 HH estimadas con PERT suben a lo más 648 HH; con 90 \%, 988 HH. Esas diferencias se cubren con las 1.280 HH de corrección de 1.4.3.7 y 1.4.4.7. Si todas las desviaciones van en el mismo sentido, las 5.372 HH se cubren con la reserva de gestión, fuera de la EDT, que libera el Comité Ejecutivo.

Los 21 paquetes sin PERT suman 141.906 HH (77,5 \%). Son turnos 24x7 (Art. 78.2), acompañamiento, estabilización, comités, capacidad mensual de mantención y saldo de los roles clave.

La faena trabaja en turnos 7x7 (siete días de trabajo y siete de descanso) de doce horas. H1 a H12 son los hitos del Formulario E-25 (sección E), y “M” seguido de un número es un mes del contrato. CIO es el Centro Integrado de Operaciones del cliente; NOC y SOC, los centros de operaciones de red y de seguridad; QA, el ambiente de pruebas; Preproducción, el ambiente igual a producción donde se ensaya cada paso; ERP, el sistema de gestión empresarial; SD-WAN, la red de área amplia definida por software; TLS, el cifrado del transporte; WCAG, las pautas de accesibilidad web; OWASP SAMM y ASVS, los modelos de madurez y de verificación de seguridad del software; SLSA, SPDX y CycloneDX, las normas de procedencia y de inventario de componentes, y NIST, el instituto de normas de Estados Unidos. Un sprint es un ciclo de desarrollo de dos semanas, y el backlog, la lista ordenada del trabajo pendiente. Los códigos RP, RF, RNF, TRA y SUP son los del Formulario T-12 y del Subdocumento 2.

La Contraparte Técnica revisa cada versión en diez días hábiles y Overfore corrige las observaciones en otros diez (Art. 18.3). Los paquetes sin hito no generan pagos; los de la rama 1.8 acreditan el pago mensual de operación.

\begin{tabla}{I{0.121}I{0.351}C{0.391}D{0.137}}{Cálculo PERT.}{tab:calculo-pert}
  {\cab{Código EDT} & \cab{Clase y cantidad (N)} & \cab{Cálculo (HH)} & \cab{σ (HH)\newline = N(P−O)/6}}
  {Fuente: elaboración propia a partir de las fichas y la Tabla \ref{tab:clases-incertidumbre-estimacion-pert-resultado} del Formulario T-14 y de la Guía del PMBOK (PMI, 2017).}
1.1.1.1 & PERT baja, E = 24 HH.\newline N = 1, paquete completo. & $M = 6 \times 24 / (4 + 0{,}80 + 1{,}25) \approx 24$\newline $O = 0{,}80 \times 24 \approx 19$\newline $P = 6 \times 24 - 19 - 4 \times 24 = 29$ & 1,667 \\
1.1.1.2 & PERT media, E = 80 HH.\newline N = 1, paquete completo. & $M = 6 \times 80 / (4 + 0{,}67 + 1{,}50) \approx 78$\newline $O = 0{,}67 \times 78 \approx 52$\newline $P = 6 \times 80 - 52 - 4 \times 78 = 116$ & 10,667 \\
1.1.1.3 & PERT baja, E = 16 HH.\newline N = 1, paquete completo. & $M = 6 \times 16 / (4 + 0{,}80 + 1{,}25) \approx 16$\newline $O = 0{,}80 \times 16 \approx 13$\newline $P = 6 \times 16 - 13 - 4 \times 16 = 19$ & 1,000 \\
1.1.1.4 & \multicolumn{3}{I{0.896}}{Sin PERT. Saldo del JP en M1–M21. 160 HH al mes por 15 meses, 152 por 5 meses y 88 en M21, menos 1.904 HH de sus otros paquetes.\newline $H = 160 \times 15 + 152 \times 5 + 88 - 1.904 = 1.344 \mathrm{HH}$} \\
1.1.2.1 & PERT baja, E = 24 HH.\newline N = 1, paquete completo. & $M = 6 \times 24 / (4 + 0{,}80 + 1{,}25) \approx 24$\newline $O = 0{,}80 \times 24 \approx 19$\newline $P = 6 \times 24 - 19 - 4 \times 24 = 29$ & 1,667 \\
1.1.2.2 & PERT baja, E = 16 HH.\newline N = 20 informes.\newline H = 16 × 20 = 320 HH. & $M = 6 \times 16 / (4 + 0{,}80 + 1{,}25) \approx 16$\newline $O = 0{,}80 \times 16 \approx 13$\newline $P = 6 \times 16 - 13 - 4 \times 16 = 19$ & 20,000 \\
1.1.2.3 & PERT baja, E = 12 HH.\newline N = 20 solicitudes estimadas, dos en el mes 2 y una por mes entre los meses 3 y 20.\newline H = 12 × 20 = 240 HH. & $M = 6 \times 12 / (4 + 0{,}80 + 1{,}25) \approx 12$\newline $O = 0{,}80 \times 12 \approx 10$\newline $P = 6 \times 12 - 10 - 4 \times 12 = 14$ & 13,333 \\
1.1.3.1 & PERT baja, E = 160 HH.\newline N = 1, paquete completo. & $M = 6 \times 160 / (4 + 0{,}80 + 1{,}25) \approx 159$\newline $O = 0{,}80 \times 159 \approx 127$\newline $P = 6 \times 160 - 127 - 4 \times 159 = 197$ & 11,667 \\
1.1.3.2 & PERT baja, E = 8 HH.\newline N = 19 actualizaciones.\newline H = 8 × 19 = 152 HH. & $M = 6 \times 8 / (4 + 0{,}80 + 1{,}25) \approx 8$\newline $O = 0{,}80 \times 8 \approx 6$\newline $P = 6 \times 8 - 6 - 4 \times 8 = 10$ & 12,667 \\
1.1.4.1 & PERT media, E = 56 HH.\newline N = 1, paquete completo. & $M = 6 \times 56 / (4 + 0{,}67 + 1{,}50) \approx 54$\newline $O = 0{,}67 \times 54 \approx 36$\newline $P = 6 \times 56 - 36 - 4 \times 54 = 84$ & 8,000 \\
1.1.4.2 & PERT media, E = 64 HH.\newline N = 1, paquete completo. & $M = 6 \times 64 / (4 + 0{,}67 + 1{,}50) \approx 62$\newline $O = 0{,}67 \times 62 \approx 42$\newline $P = 6 \times 64 - 42 - 4 \times 62 = 94$ & 8,667 \\
1.1.4.3 & PERT baja, E = 32 HH.\newline N = 1, paquete completo. & $M = 6 \times 32 / (4 + 0{,}80 + 1{,}25) \approx 32$\newline $O = 0{,}80 \times 32 \approx 26$\newline $P = 6 \times 32 - 26 - 4 \times 32 = 38$ & 2,000 \\
1.1.4.4 & PERT baja, E = 24 HH.\newline N = 3 evaluaciones.\newline H = 24 × 3 = 72 HH. & $M = 6 \times 24 / (4 + 0{,}80 + 1{,}25) \approx 24$\newline $O = 0{,}80 \times 24 \approx 19$\newline $P = 6 \times 24 - 19 - 4 \times 24 = 29$ & 5,000 \\
1.1.5.1 & \multicolumn{3}{I{0.896}}{Sin PERT. Comités y coordinación, 40 HH al mes por 20 meses. LDA y LI suman 8 HH al mes en el comité de arquitectura por 20 meses, y el SRE 4 HH al mes en el de operación por 8 meses.\newline $H = 40 \times 20 + 8 \times 20 + 4 \times 8 = 992 \mathrm{HH}$} \\
1.1.5.2 & PERT alta, E = 16 HH.\newline N = 13 sesiones.\newline H = 16 × 13 = 208 HH. & $M = 6 \times 16 / (4 + 0{,}50 + 2{,}00) \approx 15$\newline $O = 0{,}50 \times 15 \approx 8 (\mathrm{al} \mathrm{par})$\newline $P = 6 \times 16 - 8 - 4 \times 15 = 28$ & 43,333 \\
1.1.5.3 & PERT baja, E = 16 HH.\newline N = 7 versiones (inicial en el mes 1 y una por trimestre).\newline H = 16 × 7 = 112 HH. & $M = 6 \times 16 / (4 + 0{,}80 + 1{,}25) \approx 16$\newline $O = 0{,}80 \times 16 \approx 13$\newline $P = 6 \times 16 - 13 - 4 \times 16 = 19$ & 7,000 \\
1.1.6.1 & PERT media, E = 40 HH.\newline N = 1, paquete completo. & $M = 6 \times 40 / (4 + 0{,}67 + 1{,}50) \approx 39$\newline $O = 0{,}67 \times 39 \approx 26$\newline $P = 6 \times 40 - 26 - 4 \times 39 = 58$ & 5,333 \\
1.1.6.2 & PERT alta, E = 32 HH.\newline N = 7 contratos.\newline H = 32 × 7 = 224 HH. & $M = 6 \times 32 / (4 + 0{,}50 + 2{,}00) \approx 30$\newline $O = 0{,}50 \times 30 \approx 15$\newline $P = 6 \times 32 - 15 - 4 \times 30 = 57$ & 49,000 \\
1.1.6.3 & PERT baja, E = 24 HH.\newline N = 1, paquete completo. & $M = 6 \times 24 / (4 + 0{,}80 + 1{,}25) \approx 24$\newline $O = 0{,}80 \times 24 \approx 19$\newline $P = 6 \times 24 - 19 - 4 \times 24 = 29$ & 1,667 \\
1.1.6.4 & PERT baja, E = 32 HH.\newline N = 1, paquete completo. & $M = 6 \times 32 / (4 + 0{,}80 + 1{,}25) \approx 32$\newline $O = 0{,}80 \times 32 \approx 26$\newline $P = 6 \times 32 - 26 - 4 \times 32 = 38$ & 2,000 \\
1.1.7.1 & PERT baja, E = 336 HH.\newline N = 1, paquete completo. & $M = 6 \times 336 / (4 + 0{,}80 + 1{,}25) \approx 333$\newline $O = 0{,}80 \times 333 \approx 266$\newline $P = 6 \times 336 - 266 - 4 \times 333 = 418$ & 25,333 \\
1.1.7.2 & PERT baja, E = 8 HH.\newline N = 20 meses.\newline H = 8 × 20 = 160 HH. & $M = 6 \times 8 / (4 + 0{,}80 + 1{,}25) \approx 8$\newline $O = 0{,}80 \times 8 \approx 6$\newline $P = 6 \times 8 - 6 - 4 \times 8 = 10$ & 13,333 \\
1.1.8.1 & PERT baja, E = 24 HH.\newline N = 1, paquete completo. & $M = 6 \times 24 / (4 + 0{,}80 + 1{,}25) \approx 24$\newline $O = 0{,}80 \times 24 \approx 19$\newline $P = 6 \times 24 - 19 - 4 \times 24 = 29$ & 1,667 \\
1.1.8.2 & PERT baja, E = 24 HH.\newline N = 1, paquete completo. & $M = 6 \times 24 / (4 + 0{,}80 + 1{,}25) \approx 24$\newline $O = 0{,}80 \times 24 \approx 19$\newline $P = 6 \times 24 - 19 - 4 \times 24 = 29$ & 1,667 \\
1.1.8.3 & PERT baja, E = 16 HH.\newline N = 2 etapas.\newline H = 16 × 2 = 32 HH. & $M = 6 \times 16 / (4 + 0{,}80 + 1{,}25) \approx 16$\newline $O = 0{,}80 \times 16 \approx 13$\newline $P = 6 \times 16 - 13 - 4 \times 16 = 19$ & 2,000 \\
1.1.9.1 & PERT media, E = 48 HH.\newline N = 1, paquete completo. & $M = 6 \times 48 / (4 + 0{,}67 + 1{,}50) \approx 47$\newline $O = 0{,}67 \times 47 \approx 31$\newline $P = 6 \times 48 - 31 - 4 \times 47 = 69$ & 6,333 \\
1.1.9.2 & PERT baja, E = 16 HH.\newline N = 4 actualizaciones.\newline H = 16 × 4 = 64 HH. & $M = 6 \times 16 / (4 + 0{,}80 + 1{,}25) \approx 16$\newline $O = 0{,}80 \times 16 \approx 13$\newline $P = 6 \times 16 - 13 - 4 \times 16 = 19$ & 4,000 \\
1.1.10.1 & PERT baja, E = 16 HH.\newline N = 1, paquete completo. & $M = 6 \times 16 / (4 + 0{,}80 + 1{,}25) \approx 16$\newline $O = 0{,}80 \times 16 \approx 13$\newline $P = 6 \times 16 - 13 - 4 \times 16 = 19$ & 1,000 \\
1.1.10.2 & PERT media, E = 24 HH.\newline N = 7 documentos.\newline H = 24 × 7 = 168 HH. & $M = 6 \times 24 / (4 + 0{,}67 + 1{,}50) \approx 23$\newline $O = 0{,}67 \times 23 \approx 15$\newline $P = 6 \times 24 - 15 - 4 \times 23 = 37$ & 25,667 \\
1.2.1.1 & PERT alta, E = 16 HH.\newline N = 10 fichas.\newline H = 16 × 10 = 160 HH. & $M = 6 \times 16 / (4 + 0{,}50 + 2{,}00) \approx 15$\newline $O = 0{,}50 \times 15 \approx 8 (\mathrm{al} \mathrm{par})$\newline $P = 6 \times 16 - 8 - 4 \times 15 = 28$ & 33,333 \\
1.2.1.2 & PERT alta, E = 48 HH.\newline N = 1, paquete completo. & $M = 6 \times 48 / (4 + 0{,}50 + 2{,}00) \approx 44$\newline $O = 0{,}50 \times 44 \approx 22$\newline $P = 6 \times 48 - 22 - 4 \times 44 = 90$ & 11,333 \\
1.2.2.1 & PERT media, E = 40 HH.\newline N = 6 módulos de la Etapa 1.\newline H = 40 × 6 = 240 HH. & $M = 6 \times 40 / (4 + 0{,}67 + 1{,}50) \approx 39$\newline $O = 0{,}67 \times 39 \approx 26$\newline $P = 6 \times 40 - 26 - 4 \times 39 = 58$ & 32,000 \\
1.2.2.2 & PERT media, E = 56 HH.\newline N = 1, paquete completo. & $M = 6 \times 56 / (4 + 0{,}67 + 1{,}50) \approx 54$\newline $O = 0{,}67 \times 54 \approx 36$\newline $P = 6 \times 56 - 36 - 4 \times 54 = 84$ & 8,000 \\
1.2.2.3 & PERT media, E = 16 HH.\newline N = 6 módulos de la Etapa 1.\newline H = 16 × 6 = 96 HH. & $M = 6 \times 16 / (4 + 0{,}67 + 1{,}50) \approx 16$\newline $O = 0{,}67 \times 16 \approx 11$\newline $P = 6 \times 16 - 11 - 4 \times 16 = 21$ & 10,000 \\
1.2.2.4 & \multicolumn{3}{I{0.896}}{Sin PERT. Saldo del LF en M1–M12. 160 HH al mes por 12 meses, menos 576 HH de sus otros paquetes.\newline $H = 160 \times 12 - 576 = 1.344 \mathrm{HH}$} \\
1.2.2.5 & PERT media, E = 16 HH.\newline N = 4 talleres (mina y planta, en los dos grupos del 7x7).\newline H = 16 × 4 = 64 HH. & $M = 6 \times 16 / (4 + 0{,}67 + 1{,}50) \approx 16$\newline $O = 0{,}67 \times 16 \approx 11$\newline $P = 6 \times 16 - 11 - 4 \times 16 = 21$ & 6,667 \\
1.2.2.6 & PERT media, E = 40 HH.\newline N = 1, paquete completo. & $M = 6 \times 40 / (4 + 0{,}67 + 1{,}50) \approx 39$\newline $O = 0{,}67 \times 39 \approx 26$\newline $P = 6 \times 40 - 26 - 4 \times 39 = 58$ & 5,333 \\
1.2.3.1 & PERT media, E = 12 HH.\newline N = 14 decisiones.\newline H = 12 × 14 = 168 HH. & $M = 6 \times 12 / (4 + 0{,}67 + 1{,}50) \approx 12$\newline $O = 0{,}67 \times 12 \approx 8$\newline $P = 6 \times 12 - 8 - 4 \times 12 = 16$ & 18,667 \\
1.2.3.2 & PERT media, E = 16 HH.\newline N = 1, paquete completo. & $M = 6 \times 16 / (4 + 0{,}67 + 1{,}50) \approx 16$\newline $O = 0{,}67 \times 16 \approx 11$\newline $P = 6 \times 16 - 11 - 4 \times 16 = 21$ & 1,667 \\
1.2.4.1 & PERT alta, E = 72 HH.\newline N = 1, paquete completo. & $M = 6 \times 72 / (4 + 0{,}50 + 2{,}00) \approx 66$\newline $O = 0{,}50 \times 66 \approx 33$\newline $P = 6 \times 72 - 33 - 4 \times 66 = 135$ & 17,000 \\
1.2.4.2 & PERT alta, E = 40 HH.\newline N = 3 fabricantes.\newline H = 40 × 3 = 120 HH. & $M = 6 \times 40 / (4 + 0{,}50 + 2{,}00) \approx 37$\newline $O = 0{,}50 \times 37 \approx 18 (\mathrm{al} \mathrm{par})$\newline $P = 6 \times 40 - 18 - 4 \times 37 = 74$ & 28,000 \\
1.2.5.1 & PERT media, E = 16 HH.\newline N = 6 componentes de la Etapa 2 (tres módulos nuevos y tres ampliaciones).\newline H = 16 × 6 = 96 HH. & $M = 6 \times 16 / (4 + 0{,}67 + 1{,}50) \approx 16$\newline $O = 0{,}67 \times 16 \approx 11$\newline $P = 6 \times 16 - 11 - 4 \times 16 = 21$ & 10,000 \\
1.2.5.2 & PERT media, E = 32 HH.\newline N = 1, paquete completo. & $M = 6 \times 32 / (4 + 0{,}67 + 1{,}50) \approx 31$\newline $O = 0{,}67 \times 31 \approx 21$\newline $P = 6 \times 32 - 21 - 4 \times 31 = 47$ & 4,333 \\
1.2.5.3 & PERT media, E = 40 HH.\newline N = 1, paquete completo. & $M = 6 \times 40 / (4 + 0{,}67 + 1{,}50) \approx 39$\newline $O = 0{,}67 \times 39 \approx 26$\newline $P = 6 \times 40 - 26 - 4 \times 39 = 58$ & 5,333 \\
1.2.5.4 & \multicolumn{3}{I{0.896}}{Sin PERT. Saldo del LF en M13–M21. 1.392 HH disponibles, menos 420 HH de sus otros paquetes.\newline $H = 1.392 - 420 = 972 \mathrm{HH}$} \\
1.2.5.5 & PERT alta, E = 56 HH.\newline N = 1, paquete completo. & $M = 6 \times 56 / (4 + 0{,}50 + 2{,}00) \approx 52$\newline $O = 0{,}50 \times 52 \approx 26$\newline $P = 6 \times 56 - 26 - 4 \times 52 = 102$ & 12,667 \\
1.2.6.1 & PERT alta, E = 40 HH.\newline N = 1, paquete completo. & $M = 6 \times 40 / (4 + 0{,}50 + 2{,}00) \approx 37$\newline $O = 0{,}50 \times 37 \approx 18 (\mathrm{al} \mathrm{par})$\newline $P = 6 \times 40 - 18 - 4 \times 37 = 74$ & 9,333 \\
1.2.6.2 & PERT alta, E = 24 HH.\newline N = 1, paquete completo. & $M = 6 \times 24 / (4 + 0{,}50 + 2{,}00) \approx 22$\newline $O = 0{,}50 \times 22 \approx 11$\newline $P = 6 \times 24 - 11 - 4 \times 22 = 45$ & 5,667 \\
1.3.1.1 & PERT media, E = 80 HH.\newline N = 1, paquete completo. & $M = 6 \times 80 / (4 + 0{,}67 + 1{,}50) \approx 78$\newline $O = 0{,}67 \times 78 \approx 52$\newline $P = 6 \times 80 - 52 - 4 \times 78 = 116$ & 10,667 \\
1.3.1.2 & PERT media, E = 80 HH.\newline N = 1, paquete completo. & $M = 6 \times 80 / (4 + 0{,}67 + 1{,}50) \approx 78$\newline $O = 0{,}67 \times 78 \approx 52$\newline $P = 6 \times 80 - 52 - 4 \times 78 = 116$ & 10,667 \\
1.3.1.3 & PERT media, E = 72 HH.\newline N = 1, paquete completo. & $M = 6 \times 72 / (4 + 0{,}67 + 1{,}50) \approx 70$\newline $O = 0{,}67 \times 70 \approx 47$\newline $P = 6 \times 72 - 47 - 4 \times 70 = 105$ & 9,667 \\
1.3.1.4 & PERT media, E = 40 HH.\newline N = 2 planes.\newline H = 40 × 2 = 80 HH. & $M = 6 \times 40 / (4 + 0{,}67 + 1{,}50) \approx 39$\newline $O = 0{,}67 \times 39 \approx 26$\newline $P = 6 \times 40 - 26 - 4 \times 39 = 58$ & 10,667 \\
1.3.1.5 & PERT alta, E = 152 HH.\newline N = 1, paquete completo. & $M = 6 \times 152 / (4 + 0{,}50 + 2{,}00) \approx 140$\newline $O = 0{,}50 \times 140 \approx 70$\newline $P = 6 \times 152 - 70 - 4 \times 140 = 282$ & 35,333 \\
1.3.2.1 & PERT media, E = 16 HH.\newline N = 9 dominios.\newline H = 16 × 9 = 144 HH. & $M = 6 \times 16 / (4 + 0{,}67 + 1{,}50) \approx 16$\newline $O = 0{,}67 \times 16 \approx 11$\newline $P = 6 \times 16 - 11 - 4 \times 16 = 21$ & 15,000 \\
1.3.2.2 & PERT media, E = 8 HH.\newline N = 9 dominios.\newline H = 8 × 9 = 72 HH. & $M = 6 \times 8 / (4 + 0{,}67 + 1{,}50) \approx 8$\newline $O = 0{,}67 \times 8 \approx 5$\newline $P = 6 \times 8 - 5 - 4 \times 8 = 11$ & 9,000 \\
1.3.2.3 & PERT media, E = 32 HH.\newline N = 1, paquete completo. & $M = 6 \times 32 / (4 + 0{,}67 + 1{,}50) \approx 31$\newline $O = 0{,}67 \times 31 \approx 21$\newline $P = 6 \times 32 - 21 - 4 \times 31 = 47$ & 4,333 \\
1.3.3.1 & PERT media, E = 72 HH.\newline N = 1, paquete completo. & $M = 6 \times 72 / (4 + 0{,}67 + 1{,}50) \approx 70$\newline $O = 0{,}67 \times 70 \approx 47$\newline $P = 6 \times 72 - 47 - 4 \times 70 = 105$ & 9,667 \\
1.3.3.2 & PERT media, E = 72 HH.\newline N = 1, paquete completo. & $M = 6 \times 72 / (4 + 0{,}67 + 1{,}50) \approx 70$\newline $O = 0{,}67 \times 70 \approx 47$\newline $P = 6 \times 72 - 47 - 4 \times 70 = 105$ & 9,667 \\
1.3.3.3 & PERT media, E = 56 HH.\newline N = 1, paquete completo. & $M = 6 \times 56 / (4 + 0{,}67 + 1{,}50) \approx 54$\newline $O = 0{,}67 \times 54 \approx 36$\newline $P = 6 \times 56 - 36 - 4 \times 54 = 84$ & 8,000 \\
1.3.3.4 & PERT media, E = 40 HH.\newline N = 1, paquete completo. & $M = 6 \times 40 / (4 + 0{,}67 + 1{,}50) \approx 39$\newline $O = 0{,}67 \times 39 \approx 26$\newline $P = 6 \times 40 - 26 - 4 \times 39 = 58$ & 5,333 \\
1.3.4.1 & PERT media, E = 40 HH.\newline N = 2 grupos del 7x7.\newline H = 40 × 2 = 80 HH. & $M = 6 \times 40 / (4 + 0{,}67 + 1{,}50) \approx 39$\newline $O = 0{,}67 \times 39 \approx 26$\newline $P = 6 \times 40 - 26 - 4 \times 39 = 58$ & 10,667 \\
1.3.4.2 & PERT media, E = 72 HH.\newline N = 1, paquete completo. & $M = 6 \times 72 / (4 + 0{,}67 + 1{,}50) \approx 70$\newline $O = 0{,}67 \times 70 \approx 47$\newline $P = 6 \times 72 - 47 - 4 \times 70 = 105$ & 9,667 \\
1.3.4.3 & PERT media, E = 48 HH.\newline N = 6 canales.\newline H = 48 × 6 = 288 HH. & $M = 6 \times 48 / (4 + 0{,}67 + 1{,}50) \approx 47$\newline $O = 0{,}67 \times 47 \approx 31$\newline $P = 6 \times 48 - 31 - 4 \times 47 = 69$ & 38,000 \\
1.3.5.1 & PERT media, E = 32 HH.\newline N = 6 componentes de la Etapa 2 (tres módulos nuevos y tres ampliaciones).\newline H = 32 × 6 = 192 HH. & $M = 6 \times 32 / (4 + 0{,}67 + 1{,}50) \approx 31$\newline $O = 0{,}67 \times 31 \approx 21$\newline $P = 6 \times 32 - 21 - 4 \times 31 = 47$ & 26,000 \\
1.3.5.2 & PERT alta, E = 24 HH.\newline N = 4 conectores.\newline H = 24 × 4 = 96 HH. & $M = 6 \times 24 / (4 + 0{,}50 + 2{,}00) \approx 22$\newline $O = 0{,}50 \times 22 \approx 11$\newline $P = 6 \times 24 - 11 - 4 \times 22 = 45$ & 22,667 \\
1.3.6.1 & PERT media, E = 56 HH.\newline N = 1, paquete completo. & $M = 6 \times 56 / (4 + 0{,}67 + 1{,}50) \approx 54$\newline $O = 0{,}67 \times 54 \approx 36$\newline $P = 6 \times 56 - 36 - 4 \times 54 = 84$ & 8,000 \\
1.3.6.2 & PERT media, E = 64 HH.\newline N = 1, paquete completo. & $M = 6 \times 64 / (4 + 0{,}67 + 1{,}50) \approx 62$\newline $O = 0{,}67 \times 62 \approx 42$\newline $P = 6 \times 64 - 42 - 4 \times 62 = 94$ & 8,667 \\
1.4.1.1 & PERT media, E = 48 HH.\newline N = 1, paquete completo. & $M = 6 \times 48 / (4 + 0{,}67 + 1{,}50) \approx 47$\newline $O = 0{,}67 \times 47 \approx 31$\newline $P = 6 \times 48 - 31 - 4 \times 47 = 69$ & 6,333 \\
1.4.1.2 & PERT media, E = 64 HH.\newline N = 1, paquete completo. & $M = 6 \times 64 / (4 + 0{,}67 + 1{,}50) \approx 62$\newline $O = 0{,}67 \times 62 \approx 42$\newline $P = 6 \times 64 - 42 - 4 \times 62 = 94$ & 8,667 \\
1.4.1.3 & PERT media, E = 40 HH.\newline N = 1, paquete completo. & $M = 6 \times 40 / (4 + 0{,}67 + 1{,}50) \approx 39$\newline $O = 0{,}67 \times 39 \approx 26$\newline $P = 6 \times 40 - 26 - 4 \times 39 = 58$ & 5,333 \\
1.4.1.4 & \multicolumn{3}{I{0.896}}{Sin PERT. Saldo del LD en M1–M21. 160 HH al mes por 15 meses y 152 por 6, menos 1.754 HH de sus otros paquetes (1.732 de la implementación y 22 de la operación en M21).\newline $H = 160 \times 15 + 152 \times 6 - 1.754 = 1.558 \mathrm{HH}$} \\
1.4.2.1 & PERT alta, E = 80 HH.\newline N = 19 sprints.\newline H = 80 × 19 = 1.520 HH. & $M = 6 \times 80 / (4 + 0{,}50 + 2{,}00) \approx 74$\newline $O = 0{,}50 \times 74 \approx 37$\newline $P = 6 \times 80 - 37 - 4 \times 74 = 147$ & 348,333 \\
1.4.2.2 & PERT media, E = 64 HH.\newline N = 19 sprints.\newline H = 64 × 19 = 1.216 HH. & $M = 6 \times 64 / (4 + 0{,}67 + 1{,}50) \approx 62$\newline $O = 0{,}67 \times 62 \approx 42$\newline $P = 6 \times 64 - 42 - 4 \times 62 = 94$ & 164,667 \\
1.4.2.3 & PERT alta, E = 80 HH.\newline N = 10 sprints.\newline H = 80 × 10 = 800 HH. & $M = 6 \times 80 / (4 + 0{,}50 + 2{,}00) \approx 74$\newline $O = 0{,}50 \times 74 \approx 37$\newline $P = 6 \times 80 - 37 - 4 \times 74 = 147$ & 183,333 \\
1.4.3.1 & PERT media, E = 64 HH.\newline N = 12 sprints.\newline H = 64 × 12 = 768 HH. & $M = 6 \times 64 / (4 + 0{,}67 + 1{,}50) \approx 62$\newline $O = 0{,}67 \times 62 \approx 42$\newline $P = 6 \times 64 - 42 - 4 \times 62 = 94$ & 104,000 \\
1.4.3.2 & PERT media, E = 80 HH.\newline N = 16 sprints.\newline H = 80 × 16 = 1.280 HH. & $M = 6 \times 80 / (4 + 0{,}67 + 1{,}50) \approx 78$\newline $O = 0{,}67 \times 78 \approx 52$\newline $P = 6 \times 80 - 52 - 4 \times 78 = 116$ & 170,667 \\
1.4.3.3 & PERT media, E = 76 HH.\newline N = 14 sprints.\newline H = 76 × 14 = 1.064 HH. & $M = 6 \times 76 / (4 + 0{,}67 + 1{,}50) \approx 74$\newline $O = 0{,}67 \times 74 \approx 50$\newline $P = 6 \times 76 - 50 - 4 \times 74 = 110$ & 140,000 \\
1.4.3.4 & PERT media, E = 80 HH.\newline N = 12 sprints.\newline H = 80 × 12 = 960 HH. & $M = 6 \times 80 / (4 + 0{,}67 + 1{,}50) \approx 78$\newline $O = 0{,}67 \times 78 \approx 52$\newline $P = 6 \times 80 - 52 - 4 \times 78 = 116$ & 128,000 \\
1.4.3.5 & PERT media, E = 64 HH.\newline N = 14 sprints.\newline H = 64 × 14 = 896 HH. & $M = 6 \times 64 / (4 + 0{,}67 + 1{,}50) \approx 62$\newline $O = 0{,}67 \times 62 \approx 42$\newline $P = 6 \times 64 - 42 - 4 \times 62 = 94$ & 121,333 \\
1.4.3.6 & PERT media, E = 48 HH.\newline N = 8 sprints.\newline H = 48 × 8 = 384 HH. & $M = 6 \times 48 / (4 + 0{,}67 + 1{,}50) \approx 47$\newline $O = 0{,}67 \times 47 \approx 31$\newline $P = 6 \times 48 - 31 - 4 \times 47 = 69$ & 50,667 \\
1.4.3.7 & PERT media, E = 64 HH.\newline N = 12 sprints.\newline H = 64 × 12 = 768 HH. & $M = 6 \times 64 / (4 + 0{,}67 + 1{,}50) \approx 62$\newline $O = 0{,}67 \times 62 \approx 42$\newline $P = 6 \times 64 - 42 - 4 \times 62 = 94$ & 104,000 \\
1.4.4.1 & PERT alta, E = 80 HH.\newline N = 10 sprints.\newline H = 80 × 10 = 800 HH. & $M = 6 \times 80 / (4 + 0{,}50 + 2{,}00) \approx 74$\newline $O = 0{,}50 \times 74 \approx 37$\newline $P = 6 \times 80 - 37 - 4 \times 74 = 147$ & 183,333 \\
1.4.4.2 & PERT media, E = 64 HH.\newline N = 8 sprints.\newline H = 64 × 8 = 512 HH. & $M = 6 \times 64 / (4 + 0{,}67 + 1{,}50) \approx 62$\newline $O = 0{,}67 \times 62 \approx 42$\newline $P = 6 \times 64 - 42 - 4 \times 62 = 94$ & 69,333 \\
1.4.4.3 & PERT media, E = 72 HH.\newline N = 8 sprints.\newline H = 72 × 8 = 576 HH. & $M = 6 \times 72 / (4 + 0{,}67 + 1{,}50) \approx 70$\newline $O = 0{,}67 \times 70 \approx 47$\newline $P = 6 \times 72 - 47 - 4 \times 70 = 105$ & 77,333 \\
1.4.4.4 & PERT alta, E = 64 HH.\newline N = 6 sprints.\newline H = 64 × 6 = 384 HH. & $M = 6 \times 64 / (4 + 0{,}50 + 2{,}00) \approx 59$\newline $O = 0{,}50 \times 59 \approx 30 (\mathrm{al} \mathrm{par})$\newline $P = 6 \times 64 - 30 - 4 \times 59 = 118$ & 88,000 \\
1.4.4.5 & PERT media, E = 64 HH.\newline N = 6 sprints.\newline H = 64 × 6 = 384 HH. & $M = 6 \times 64 / (4 + 0{,}67 + 1{,}50) \approx 62$\newline $O = 0{,}67 \times 62 \approx 42$\newline $P = 6 \times 64 - 42 - 4 \times 62 = 94$ & 52,000 \\
1.4.4.6 & PERT media, E = 64 HH.\newline N = 6 sprints.\newline H = 64 × 6 = 384 HH. & $M = 6 \times 64 / (4 + 0{,}67 + 1{,}50) \approx 62$\newline $O = 0{,}67 \times 62 \approx 42$\newline $P = 6 \times 64 - 42 - 4 \times 62 = 94$ & 52,000 \\
1.4.4.7 & PERT media, E = 64 HH.\newline N = 8 ciclos.\newline H = 64 × 8 = 512 HH. & $M = 6 \times 64 / (4 + 0{,}67 + 1{,}50) \approx 62$\newline $O = 0{,}67 \times 62 \approx 42$\newline $P = 6 \times 64 - 42 - 4 \times 62 = 94$ & 69,333 \\
1.4.5.1 & PERT media, E = 64 HH.\newline N = 19 sprints.\newline H = 64 × 19 = 1.216 HH. & $M = 6 \times 64 / (4 + 0{,}67 + 1{,}50) \approx 62$\newline $O = 0{,}67 \times 62 \approx 42$\newline $P = 6 \times 64 - 42 - 4 \times 62 = 94$ & 164,667 \\
1.4.5.2 & PERT media, E = 64 HH.\newline N = 16 sprints.\newline H = 64 × 16 = 1.024 HH. & $M = 6 \times 64 / (4 + 0{,}67 + 1{,}50) \approx 62$\newline $O = 0{,}67 \times 62 \approx 42$\newline $P = 6 \times 64 - 42 - 4 \times 62 = 94$ & 138,667 \\
1.4.6.1 & PERT alta, E = 24 HH.\newline N = 7 orígenes de la migración (Subdocumento 5, Tabla \refext{S5}{tab:volumen-migracion-origen}).\newline H = 24 × 7 = 168 HH. & $M = 6 \times 24 / (4 + 0{,}50 + 2{,}00) \approx 22$\newline $O = 0{,}50 \times 22 \approx 11$\newline $P = 6 \times 24 - 11 - 4 \times 22 = 45$ & 39,667 \\
1.4.6.2 & PERT alta, E = 80 HH.\newline N = 3 cargas (dos ensayos y la carga final).\newline H = 80 × 3 = 240 HH. & $M = 6 \times 80 / (4 + 0{,}50 + 2{,}00) \approx 74$\newline $O = 0{,}50 \times 74 \approx 37$\newline $P = 6 \times 80 - 37 - 4 \times 74 = 147$ & 55,000 \\
1.4.6.3 & PERT baja, E = 3.184 HH.\newline N = 1, paquete completo. & $M = 6 \times 3.184 / (4 + 0{,}80 + 1{,}25) \approx 3.158$\newline $O = 0{,}80 \times 3.158 \approx 2.526$\newline $P = 6 \times 3.184 - 2.526 - 4 \times 3.158 = 3.946$ & 236,667 \\
1.4.6.4 & PERT alta, E = 64 HH.\newline N = 3 cargas (dos ensayos y la carga final).\newline H = 64 × 3 = 192 HH. & $M = 6 \times 64 / (4 + 0{,}50 + 2{,}00) \approx 59$\newline $O = 0{,}50 \times 59 \approx 30 (\mathrm{al} \mathrm{par})$\newline $P = 6 \times 64 - 30 - 4 \times 59 = 118$ & 44,000 \\
1.4.6.5 & PERT media, E = 120 HH.\newline N = 1, paquete completo. & $M = 6 \times 120 / (4 + 0{,}67 + 1{,}50) \approx 117$\newline $O = 0{,}67 \times 117 \approx 78$\newline $P = 6 \times 120 - 78 - 4 \times 117 = 174$ & 16,000 \\
1.5.1.1 & PERT media, E = 24 HH.\newline N = 32 quincenas de los meses 3 a 18.\newline H = 24 × 32 = 768 HH. & $M = 6 \times 24 / (4 + 0{,}67 + 1{,}50) \approx 23$\newline $O = 0{,}67 \times 23 \approx 15$\newline $P = 6 \times 24 - 15 - 4 \times 23 = 37$ & 117,333 \\
1.5.1.2 & PERT media, E = 32 HH.\newline N = 32 quincenas de los meses 3 a 18.\newline H = 32 × 32 = 1.024 HH. & $M = 6 \times 32 / (4 + 0{,}67 + 1{,}50) \approx 31$\newline $O = 0{,}67 \times 31 \approx 21$\newline $P = 6 \times 32 - 21 - 4 \times 31 = 47$ & 138,667 \\
1.5.1.3 & PERT baja, E = 8 HH.\newline N = 32 quincenas de los meses 3 a 18.\newline H = 8 × 32 = 256 HH. & $M = 6 \times 8 / (4 + 0{,}80 + 1{,}25) \approx 8$\newline $O = 0{,}80 \times 8 \approx 6$\newline $P = 6 \times 8 - 6 - 4 \times 8 = 10$ & 21,333 \\
1.5.2.1 & PERT media, E = 48 HH.\newline N = 4 canales.\newline H = 48 × 4 = 192 HH. & $M = 6 \times 48 / (4 + 0{,}67 + 1{,}50) \approx 47$\newline $O = 0{,}67 \times 47 \approx 31$\newline $P = 6 \times 48 - 31 - 4 \times 47 = 69$ & 25,333 \\
1.5.2.2 & PERT media, E = 48 HH.\newline N = 3 canales.\newline H = 48 × 3 = 144 HH. & $M = 6 \times 48 / (4 + 0{,}67 + 1{,}50) \approx 47$\newline $O = 0{,}67 \times 47 \approx 31$\newline $P = 6 \times 48 - 31 - 4 \times 47 = 69$ & 19,000 \\
1.5.3.1 & PERT media, E = 40 HH.\newline N = 3 pruebas (carga sostenida, estrés y resistencia).\newline H = 40 × 3 = 120 HH. & $M = 6 \times 40 / (4 + 0{,}67 + 1{,}50) \approx 39$\newline $O = 0{,}67 \times 39 \approx 26$\newline $P = 6 \times 40 - 26 - 4 \times 39 = 58$ & 16,000 \\
1.5.3.2 & PERT media, E = 16 HH.\newline N = 5 experimentos de falla.\newline H = 16 × 5 = 80 HH. & $M = 6 \times 16 / (4 + 0{,}67 + 1{,}50) \approx 16$\newline $O = 0{,}67 \times 16 \approx 11$\newline $P = 6 \times 16 - 11 - 4 \times 16 = 21$ & 8,333 \\
1.5.3.3 & PERT media, E = 40 HH.\newline N = 6 módulos de la Etapa 1.\newline H = 40 × 6 = 240 HH. & $M = 6 \times 40 / (4 + 0{,}67 + 1{,}50) \approx 39$\newline $O = 0{,}67 \times 39 \approx 26$\newline $P = 6 \times 40 - 26 - 4 \times 39 = 58$ & 32,000 \\
1.5.4.1 & PERT media, E = 336 HH.\newline N = 1, paquete completo. & $M = 6 \times 336 / (4 + 0{,}67 + 1{,}50) \approx 327$\newline $O = 0{,}67 \times 327 \approx 219$\newline $P = 6 \times 336 - 219 - 4 \times 327 = 489$ & 45,000 \\
1.5.4.2 & PERT media, E = 32 HH.\newline N = 1, paquete completo. & $M = 6 \times 32 / (4 + 0{,}67 + 1{,}50) \approx 31$\newline $O = 0{,}67 \times 31 \approx 21$\newline $P = 6 \times 32 - 21 - 4 \times 31 = 47$ & 4,333 \\
1.5.4.3 & PERT baja, E = 8 HH.\newline N = 18 meses.\newline H = 8 × 18 = 144 HH. & $M = 6 \times 8 / (4 + 0{,}80 + 1{,}25) \approx 8$\newline $O = 0{,}80 \times 8 \approx 6$\newline $P = 6 \times 8 - 6 - 4 \times 8 = 10$ & 12,000 \\
1.5.5.1 & PERT media, E = 40 HH.\newline N = 2 ejercicios.\newline H = 40 × 2 = 80 HH. & $M = 6 \times 40 / (4 + 0{,}67 + 1{,}50) \approx 39$\newline $O = 0{,}67 \times 39 \approx 26$\newline $P = 6 \times 40 - 26 - 4 \times 39 = 58$ & 10,667 \\
1.5.5.2 & PERT media, E = 80 HH.\newline N = 2 cortes.\newline H = 80 × 2 = 160 HH. & $M = 6 \times 80 / (4 + 0{,}67 + 1{,}50) \approx 78$\newline $O = 0{,}67 \times 78 \approx 52$\newline $P = 6 \times 80 - 52 - 4 \times 78 = 116$ & 21,333 \\
1.5.6.1 & PERT media, E = 40 HH.\newline N = 3 pruebas.\newline H = 40 × 3 = 120 HH. & $M = 6 \times 40 / (4 + 0{,}67 + 1{,}50) \approx 39$\newline $O = 0{,}67 \times 39 \approx 26$\newline $P = 6 \times 40 - 26 - 4 \times 39 = 58$ & 16,000 \\
1.5.6.2 & PERT media, E = 16 HH.\newline N = 5 experimentos de falla.\newline H = 16 × 5 = 80 HH. & $M = 6 \times 16 / (4 + 0{,}67 + 1{,}50) \approx 16$\newline $O = 0{,}67 \times 16 \approx 11$\newline $P = 6 \times 16 - 11 - 4 \times 16 = 21$ & 8,333 \\
1.5.6.3 & PERT media, E = 64 HH.\newline N = 1, paquete completo. & $M = 6 \times 64 / (4 + 0{,}67 + 1{,}50) \approx 62$\newline $O = 0{,}67 \times 62 \approx 42$\newline $P = 6 \times 64 - 42 - 4 \times 62 = 94$ & 8,667 \\
1.6.1.1 & PERT media, E = 248 HH.\newline N = 1, paquete completo. & $M = 6 \times 248 / (4 + 0{,}67 + 1{,}50) \approx 241$\newline $O = 0{,}67 \times 241 \approx 161$\newline $P = 6 \times 248 - 161 - 4 \times 241 = 363$ & 33,667 \\
1.6.1.2 & PERT media, E = 224 HH.\newline N = 1, paquete completo. & $M = 6 \times 224 / (4 + 0{,}67 + 1{,}50) \approx 218$\newline $O = 0{,}67 \times 218 \approx 146$\newline $P = 6 \times 224 - 146 - 4 \times 218 = 326$ & 30,000 \\
1.6.1.3 & PERT alta, E = 48 HH.\newline N = 3 enlaces (segunda fibra, satélite y diodo).\newline H = 48 × 3 = 144 HH. & $M = 6 \times 48 / (4 + 0{,}50 + 2{,}00) \approx 44$\newline $O = 0{,}50 \times 44 \approx 22$\newline $P = 6 \times 48 - 22 - 4 \times 44 = 90$ & 34,000 \\
1.6.2.1 & PERT media, E = 336 HH.\newline N = 1, paquete completo. & $M = 6 \times 336 / (4 + 0{,}67 + 1{,}50) \approx 327$\newline $O = 0{,}67 \times 327 \approx 219$\newline $P = 6 \times 336 - 219 - 4 \times 327 = 489$ & 45,000 \\
1.6.2.2 & PERT media, E = 256 HH.\newline N = 1, paquete completo. & $M = 6 \times 256 / (4 + 0{,}67 + 1{,}50) \approx 249$\newline $O = 0{,}67 \times 249 \approx 167$\newline $P = 6 \times 256 - 167 - 4 \times 249 = 373$ & 34,333 \\
1.6.2.3 & PERT media, E = 64 HH.\newline N = 1, paquete completo. & $M = 6 \times 64 / (4 + 0{,}67 + 1{,}50) \approx 62$\newline $O = 0{,}67 \times 62 \approx 42$\newline $P = 6 \times 64 - 42 - 4 \times 62 = 94$ & 8,667 \\
1.6.3.1 & \multicolumn{3}{I{0.896}}{Sin PERT. Marcha blanca de la Etapa 1, 40 HH a la semana durante 13 semanas (M13–M15).\newline $H = 40 \times 13 = 520 \mathrm{HH}$} \\
1.6.3.2 & \multicolumn{3}{I{0.896}}{Sin PERT. Marcha blanca de la Etapa 2, 40 HH a la semana durante 9 semanas, que cubren los 62 días de julio y agosto de 2028.\newline $H = 40 \times 9 = 360 \mathrm{HH}$} \\
1.6.3.3 & \multicolumn{3}{I{0.896}}{Sin PERT. Acompañamiento en 560 turnos de una persona, de 12 HH cada uno, más 8 HH a la semana de coordinación de LF y LIG durante 13 semanas.\newline $H = 12 \times 560 + 8 \times 13 = 6.824 \mathrm{HH}$} \\
1.6.3.4 & \multicolumn{3}{I{0.896}}{Sin PERT. Acompañamiento en 330 turnos de una persona, de 12 HH cada uno, más 32 HH a la semana de implantación y atención funcional durante 9 semanas.\newline $H = 12 \times 330 + 32 \times 9 = 4.248 \mathrm{HH}$} \\
1.6.3.5 & PERT media, E = 48 HH.\newline N = 1, paquete completo. & $M = 6 \times 48 / (4 + 0{,}67 + 1{,}50) \approx 47$\newline $O = 0{,}67 \times 47 \approx 31$\newline $P = 6 \times 48 - 31 - 4 \times 47 = 69$ & 6,333 \\
1.6.4.1 & \multicolumn{3}{I{0.896}}{Sin PERT. Estabilización en 166 turnos de una persona, de 12 HH cada uno, más 72 HH del SRE para la ventana de corte y la coordinación.\newline $H = 12 \times 166 + 72 = 2.064 \mathrm{HH}$} \\
1.6.4.2 & \multicolumn{3}{I{0.896}}{Sin PERT. Estabilización en 140 turnos de una persona, de 12 HH cada uno, más 72 HH del SRE para la ventana de corte y la coordinación.\newline $H = 12 \times 140 + 72 = 1.752 \mathrm{HH}$} \\
1.6.4.3 & \multicolumn{3}{I{0.896}}{Sin PERT. Atención en 184 turnos de 12 HH, más 24 HH al mes de coordinación del SRE durante 3 meses.\newline $H = 12 \times 184 + 24 \times 3 = 2.280 \mathrm{HH}$} \\
1.6.5.1 & PERT baja, E = 50 HH.\newline N = 24 semanas de ola (tres olas de cuatro semanas por etapa).\newline H = 50 × 24 = 1.200 HH. & $M = 6 \times 50 / (4 + 0{,}80 + 1{,}25) \approx 50$\newline $O = 0{,}80 \times 50 \approx 40$\newline $P = 6 \times 50 - 40 - 4 \times 50 = 60$ & 80,000 \\
1.6.5.2 & PERT baja, E = 32 HH.\newline N = 9 meses.\newline H = 32 × 9 = 288 HH. & $M = 6 \times 32 / (4 + 0{,}80 + 1{,}25) \approx 32$\newline $O = 0{,}80 \times 32 \approx 26$\newline $P = 6 \times 32 - 26 - 4 \times 32 = 38$ & 18,000 \\
1.6.5.3 & PERT media, E = 40 HH.\newline N = 6 meses.\newline H = 40 × 6 = 240 HH. & $M = 6 \times 40 / (4 + 0{,}67 + 1{,}50) \approx 39$\newline $O = 0{,}67 \times 39 \approx 26$\newline $P = 6 \times 40 - 26 - 4 \times 39 = 58$ & 32,000 \\
1.6.6.1 & PERT media, E = 64 HH.\newline N = 1, paquete completo. & $M = 6 \times 64 / (4 + 0{,}67 + 1{,}50) \approx 62$\newline $O = 0{,}67 \times 62 \approx 42$\newline $P = 6 \times 64 - 42 - 4 \times 62 = 94$ & 8,667 \\
1.6.6.2 & PERT baja, E = 8 HH.\newline N = 16 meses.\newline H = 8 × 16 = 128 HH. & $M = 6 \times 8 / (4 + 0{,}80 + 1{,}25) \approx 8$\newline $O = 0{,}80 \times 8 \approx 6$\newline $P = 6 \times 8 - 6 - 4 \times 8 = 10$ & 10,667 \\
1.6.6.3 & PERT media, E = 24 HH.\newline N = 11 meses.\newline H = 24 × 11 = 264 HH. & $M = 6 \times 24 / (4 + 0{,}67 + 1{,}50) \approx 23$\newline $O = 0{,}67 \times 23 \approx 15$\newline $P = 6 \times 24 - 15 - 4 \times 23 = 37$ & 40,333 \\
1.6.6.4 & PERT baja, E = 16 HH.\newline N = 14 informes.\newline H = 16 × 14 = 224 HH. & $M = 6 \times 16 / (4 + 0{,}80 + 1{,}25) \approx 16$\newline $O = 0{,}80 \times 16 \approx 13$\newline $P = 6 \times 16 - 13 - 4 \times 16 = 19$ & 14,000 \\
1.6.6.5 & PERT baja, E = 24 HH.\newline N = 6 olas (tres por etapa).\newline H = 24 × 6 = 144 HH. & $M = 6 \times 24 / (4 + 0{,}80 + 1{,}25) \approx 24$\newline $O = 0{,}80 \times 24 \approx 19$\newline $P = 6 \times 24 - 19 - 4 \times 24 = 29$ & 10,000 \\
1.7.1.1 & PERT media, E = 40 HH.\newline N = 10 sprints.\newline H = 40 × 10 = 400 HH. & $M = 6 \times 40 / (4 + 0{,}67 + 1{,}50) \approx 39$\newline $O = 0{,}67 \times 39 \approx 26$\newline $P = 6 \times 40 - 26 - 4 \times 39 = 58$ & 53,333 \\
1.7.1.2 & PERT baja, E = 16 HH.\newline N = 4 informes (línea base y tres evaluaciones).\newline H = 16 × 4 = 64 HH. & $M = 6 \times 16 / (4 + 0{,}80 + 1{,}25) \approx 16$\newline $O = 0{,}80 \times 16 \approx 13$\newline $P = 6 \times 16 - 13 - 4 \times 16 = 19$ & 4,000 \\
1.7.2.1 & PERT alta, E = 880 HH.\newline N = 1, paquete completo. & $M = 6 \times 880 / (4 + 0{,}50 + 2{,}00) \approx 812$\newline $O = 0{,}50 \times 812 \approx 406$\newline $P = 6 \times 880 - 406 - 4 \times 812 = 1.626$ & 203,333 \\
1.7.2.2 & PERT baja, E = 16 HH.\newline N = 3 informes.\newline H = 16 × 3 = 48 HH. & $M = 6 \times 16 / (4 + 0{,}80 + 1{,}25) \approx 16$\newline $O = 0{,}80 \times 16 \approx 13$\newline $P = 6 \times 16 - 13 - 4 \times 16 = 19$ & 3,000 \\
1.7.2.3 & PERT baja, E = 456 HH.\newline N = 1, paquete completo. & $M = 6 \times 456 / (4 + 0{,}80 + 1{,}25) \approx 452$\newline $O = 0{,}80 \times 452 \approx 362$\newline $P = 6 \times 456 - 362 - 4 \times 452 = 566$ & 34,000 \\
1.7.2.4 & PERT media, E = 80 HH.\newline N = 1, paquete completo. & $M = 6 \times 80 / (4 + 0{,}67 + 1{,}50) \approx 78$\newline $O = 0{,}67 \times 78 \approx 52$\newline $P = 6 \times 80 - 52 - 4 \times 78 = 116$ & 10,667 \\
1.7.3.1 & PERT media, E = 24 HH.\newline N = 3 instrumentos.\newline H = 24 × 3 = 72 HH. & $M = 6 \times 24 / (4 + 0{,}67 + 1{,}50) \approx 23$\newline $O = 0{,}67 \times 23 \approx 15$\newline $P = 6 \times 24 - 15 - 4 \times 23 = 37$ & 11,000 \\
1.7.3.2 & PERT media, E = 40 HH.\newline N = 12 sprints.\newline H = 40 × 12 = 480 HH. & $M = 6 \times 40 / (4 + 0{,}67 + 1{,}50) \approx 39$\newline $O = 0{,}67 \times 39 \approx 26$\newline $P = 6 \times 40 - 26 - 4 \times 39 = 58$ & 64,000 \\
1.7.3.3 & PERT baja, E = 16 HH.\newline N = 3 informes.\newline H = 16 × 3 = 48 HH. & $M = 6 \times 16 / (4 + 0{,}80 + 1{,}25) \approx 16$\newline $O = 0{,}80 \times 16 \approx 13$\newline $P = 6 \times 16 - 13 - 4 \times 16 = 19$ & 3,000 \\
1.7.4.1 & PERT media, E = 56 HH.\newline N = 1, paquete completo. & $M = 6 \times 56 / (4 + 0{,}67 + 1{,}50) \approx 54$\newline $O = 0{,}67 \times 54 \approx 36$\newline $P = 6 \times 56 - 36 - 4 \times 54 = 84$ & 8,000 \\
1.7.4.2 & PERT baja, E = 24 HH.\newline N = 4 actas.\newline H = 24 × 4 = 96 HH. & $M = 6 \times 24 / (4 + 0{,}80 + 1{,}25) \approx 24$\newline $O = 0{,}80 \times 24 \approx 19$\newline $P = 6 \times 24 - 19 - 4 \times 24 = 29$ & 6,667 \\
1.7.5.1 & PERT media, E = 40 HH.\newline N = 6 sprints.\newline H = 40 × 6 = 240 HH. & $M = 6 \times 40 / (4 + 0{,}67 + 1{,}50) \approx 39$\newline $O = 0{,}67 \times 39 \approx 26$\newline $P = 6 \times 40 - 26 - 4 \times 39 = 58$ & 32,000 \\
1.7.5.2 & PERT baja, E = 16 HH.\newline N = 5 informes.\newline H = 16 × 5 = 80 HH. & $M = 6 \times 16 / (4 + 0{,}80 + 1{,}25) \approx 16$\newline $O = 0{,}80 \times 16 \approx 13$\newline $P = 6 \times 16 - 13 - 4 \times 16 = 19$ & 5,000 \\
1.8.1.1 & \multicolumn{3}{I{0.896}}{Sin PERT. Mesa de ayuda 24x7 en M21–M56 (1.095 días), con refuerzo en las 6 detenciones de planta.\newline $\mathrm{Turnos} \mathrm{regulares} = 2 \times 1.095 = 2.190$\newline $\mathrm{Turnos} \mathrm{de} \mathrm{refuerzo} = 12 \times 6 = 72$\newline $H = 12 \times 2.190 + 12 \times 72 = 27.144 \mathrm{HH}$} \\
1.8.1.2 & \multicolumn{3}{I{0.896}}{Sin PERT. Técnico en faena 24x7 en M21–M56 (1.095 días), acreditación anual de 5 técnicos por 3 años y una atención por trimestre en el puerto.\newline $\mathrm{Turnos} = 2 \times 1.095 = 2.190$\newline $\mathrm{Renovaciones} = 5 \times 3 = 15$\newline $\mathrm{Atenciones} \mathrm{en} \mathrm{puerto} = 4 \times 3 = 12$\newline $H = 12 \times 2.190 + 8 \times 15 + 12 \times 12 = 26.544 \mathrm{HH}$} \\
1.8.2.1 & \multicolumn{3}{I{0.896}}{Sin PERT. NOC 24x7 en M21–M56 (1.095 días).\newline $\mathrm{Turnos} = 2 \times 1.095 = 2.190$\newline $H = 12 \times 2.190 = 26.280 \mathrm{HH}$} \\
1.8.2.2 & \multicolumn{3}{I{0.896}}{Sin PERT. SOC 24x7 en M21–M56 (1.095 días), más 8 HH al mes de revisión del ESI durante 36 meses.\newline $\mathrm{Turnos} = 2 \times 1.095 = 2.190$\newline $H = 12 \times 2.190 + 8 \times 36 = 26.568 \mathrm{HH}$} \\
1.8.2.3 & PERT baja, E = 312 HH.\newline N = 1, paquete completo. & $M = 6 \times 312 / (4 + 0{,}80 + 1{,}25) \approx 309$\newline $O = 0{,}80 \times 309 \approx 247$\newline $P = 6 \times 312 - 247 - 4 \times 309 = 389$ & 23,667 \\
1.8.3.1 & \multicolumn{3}{I{0.896}}{Sin PERT. Bolsa de mantención correctiva de 80 HH al mes durante los 36 meses de operación.\newline $H = 80 \times 36 = 2.880 \mathrm{HH}$} \\
1.8.3.2 & PERT baja, E = 16 HH.\newline N = 12 informes estimados (uno por trimestre).\newline H = 16 × 12 = 192 HH. & $M = 6 \times 16 / (4 + 0{,}80 + 1{,}25) \approx 16$\newline $O = 0{,}80 \times 16 \approx 13$\newline $P = 6 \times 16 - 13 - 4 \times 16 = 19$ & 12,000 \\
1.8.4.1 & PERT baja, E = 80 HH.\newline N = 3 planes.\newline H = 80 × 3 = 240 HH. & $M = 6 \times 80 / (4 + 0{,}80 + 1{,}25) \approx 79$\newline $O = 0{,}80 \times 79 \approx 63$\newline $P = 6 \times 80 - 63 - 4 \times 79 = 101$ & 19,000 \\
1.8.4.2 & PERT baja, E = 216 HH.\newline N = 1, paquete completo. & $M = 6 \times 216 / (4 + 0{,}80 + 1{,}25) \approx 214$\newline $O = 0{,}80 \times 214 \approx 171$\newline $P = 6 \times 216 - 171 - 4 \times 214 = 269$ & 16,333 \\
1.8.4.3 & PERT alta, E = 80 HH.\newline N = 3 adecuaciones estimadas (una por año).\newline H = 80 × 3 = 240 HH. & $M = 6 \times 80 / (4 + 0{,}50 + 2{,}00) \approx 74$\newline $O = 0{,}50 \times 74 \approx 37$\newline $P = 6 \times 80 - 37 - 4 \times 74 = 147$ & 55,000 \\
1.8.5.1 & \multicolumn{3}{I{0.896}}{Sin PERT. Mantención evolutiva de 80 HH al mes durante 36 meses, más 40 HH de la prueba de replicación.\newline $H = 80 \times 36 + 40 = 2.920 \mathrm{HH}$} \\
1.8.5.2 & PERT baja, E = 8 HH.\newline N = 36 meses.\newline H = 8 × 36 = 288 HH. & $M = 6 \times 8 / (4 + 0{,}80 + 1{,}25) \approx 8$\newline $O = 0{,}80 \times 8 \approx 6$\newline $P = 6 \times 8 - 6 - 4 \times 8 = 10$ & 24,000 \\
1.8.6.1 & PERT baja, E = 32 HH.\newline N = 3 actualizaciones.\newline H = 32 × 3 = 96 HH. & $M = 6 \times 32 / (4 + 0{,}80 + 1{,}25) \approx 32$\newline $O = 0{,}80 \times 32 \approx 26$\newline $P = 6 \times 32 - 26 - 4 \times 32 = 38$ & 6,000 \\
1.8.6.2 & PERT baja, E = 32 HH.\newline N = 3 actualizaciones.\newline H = 32 × 3 = 96 HH. & $M = 6 \times 32 / (4 + 0{,}80 + 1{,}25) \approx 32$\newline $O = 0{,}80 \times 32 \approx 26$\newline $P = 6 \times 32 - 26 - 4 \times 32 = 38$ & 6,000 \\
1.8.7.1 & PERT baja, E = 16 HH.\newline N = 36 informes.\newline H = 16 × 36 = 576 HH. & $M = 6 \times 16 / (4 + 0{,}80 + 1{,}25) \approx 16$\newline $O = 0{,}80 \times 16 \approx 13$\newline $P = 6 \times 16 - 13 - 4 \times 16 = 19$ & 36,000 \\
1.8.7.2 & PERT baja, E = 64 HH.\newline N = 1, paquete completo. & $M = E = 64$\newline $O = 0{,}80 \times 64 \approx 51$\newline $P = 6 \times 64 - 51 - 4 \times 64 = 77$ & 4,333 \\
1.8.7.3 & PERT baja, E = 544 HH.\newline N = 1, paquete completo. & $M = 6 \times 544 / (4 + 0{,}80 + 1{,}25) \approx 540$\newline $O = 0{,}80 \times 540 \approx 432$\newline $P = 6 \times 544 - 432 - 4 \times 540 = 672$ & 40,000 \\
1.8.7.4 & \multicolumn{3}{I{0.896}}{Sin PERT. Comités de operación, 24 HH al mes durante los 36 meses de operación.\newline $H = 24 \times 36 = 864 \mathrm{HH}$} \\
1.8.7.5 & \multicolumn{3}{I{0.896}}{Sin PERT. Saldo del GS. 64 HH en M21 y 152 HH al mes en los 35 meses siguientes, menos 1.456 HH de sus otros paquetes.\newline $H = 64 + 152 \times 35 - 1.456 = 3.928 \mathrm{HH}$} \\
1.8.8.1 & PERT baja, E = 216 HH.\newline N = 1, paquete completo. & $M = 6 \times 216 / (4 + 0{,}80 + 1{,}25) \approx 214$\newline $O = 0{,}80 \times 214 \approx 171$\newline $P = 6 \times 216 - 171 - 4 \times 214 = 269$ & 16,333 \\
1.8.8.2 & PERT media, E = 480 HH.\newline N = 1, paquete completo. & $M = 6 \times 480 / (4 + 0{,}67 + 1{,}50) \approx 467$\newline $O = 0{,}67 \times 467 \approx 313$\newline $P = 6 \times 480 - 313 - 4 \times 467 = 699$ & 64,333 \\
1.8.8.3 & PERT baja, E = 40 HH.\newline N = 3 pruebas.\newline H = 40 × 3 = 120 HH. & $M = 6 \times 40 / (4 + 0{,}80 + 1{,}25) \approx 40$\newline $O = 0{,}80 \times 40 \approx 32$\newline $P = 6 \times 40 - 32 - 4 \times 40 = 48$ & 8,000 \\
1.8.8.4 & PERT baja, E = 16 HH.\newline N = 36 versiones.\newline H = 16 × 36 = 576 HH. & $M = 6 \times 16 / (4 + 0{,}80 + 1{,}25) \approx 16$\newline $O = 0{,}80 \times 16 \approx 13$\newline $P = 6 \times 16 - 13 - 4 \times 16 = 19$ & 36,000 \\
1.8.8.5 & PERT baja, E = 8 HH.\newline N = 36 pruebas.\newline H = 8 × 36 = 288 HH. & $M = 6 \times 8 / (4 + 0{,}80 + 1{,}25) \approx 8$\newline $O = 0{,}80 \times 8 \approx 6$\newline $P = 6 \times 8 - 6 - 4 \times 8 = 10$ & 24,000 \\
1.8.8.6 & PERT baja, E = 16 HH.\newline N = 2 reevaluaciones.\newline H = 16 × 2 = 32 HH. & $M = 6 \times 16 / (4 + 0{,}80 + 1{,}25) \approx 16$\newline $O = 0{,}80 \times 16 \approx 13$\newline $P = 6 \times 16 - 13 - 4 \times 16 = 19$ & 2,000 \\
1.8.9.1 & PERT baja, E = 24 HH.\newline N = 6 jornadas.\newline H = 24 × 6 = 144 HH. & $M = 6 \times 24 / (4 + 0{,}80 + 1{,}25) \approx 24$\newline $O = 0{,}80 \times 24 \approx 19$\newline $P = 6 \times 24 - 19 - 4 \times 24 = 29$ & 10,000 \\
1.8.9.2 & PERT baja, E = 8 HH.\newline N = 18 inducciones (una por trimestre y una por cada detención de planta).\newline H = 8 × 18 = 144 HH. & $M = 6 \times 8 / (4 + 0{,}80 + 1{,}25) \approx 8$\newline $O = 0{,}80 \times 8 \approx 6$\newline $P = 6 \times 8 - 6 - 4 \times 8 = 10$ & 12,000 \\
1.8.10.1 & PERT media, E = 40 HH.\newline N = 4 entregas (código, datos, configuraciones y credenciales).\newline H = 40 × 4 = 160 HH. & $M = 6 \times 40 / (4 + 0{,}67 + 1{,}50) \approx 39$\newline $O = 0{,}67 \times 39 \approx 26$\newline $P = 6 \times 40 - 26 - 4 \times 39 = 58$ & 21,333 \\
1.8.10.2 & \multicolumn{3}{I{0.896}}{Sin PERT. Acompañamiento de 90 días, 40 HH a la semana durante 13 semanas.\newline $H = 40 \times 13 = 520 \mathrm{HH}$} \\
PERT & Subtotal de 154 paquetes con tres escenarios. & $H = \Sigma H_{p} = 41.160 \mathrm{HH}$\newline $\sigma = \surd (\Sigma \sigma _{p}^{2})$ & 771,745 \\
Sin PERT & Subtotal de 21 paquetes de cobertura, capacidad o saldo de dedicación. & $H = \Sigma H_{p} = 141.906 \mathrm{HH}$ & No aplica \\
Total & 175 paquetes. & $H = 41.160 + 141.906 = 183.066 \mathrm{HH}$ & 771,745 \\
\end{tabla}

\apartado{Estructura de descomposición del trabajo}
La EDT divide el trabajo de implementación y operación en paquetes que se pueden estimar y asignar a un responsable. La Figura \ref{fig:edt-sitco-aranda-vista-general-ocho} resume sus ocho ramas y las figuras \ref{fig:edt-sitco-aranda-ramas-1-1} a \ref{fig:edt-sitco-aranda-ramas-1-7} muestran los entregables de dos ramas en cada figura, desde dirección y requerimientos hasta innovaciones y operación.

\figura{16.09}{fig-B-1-edt-sitco-aranda-vista-general}{EDT de SITCO Aranda, vista general de las ocho ramas.}{fig:edt-sitco-aranda-vista-general-ocho}{Fuente: elaboración propia a partir de la Guía del PMBOK (PMI, 2017, pp. 156–162).}
La operación 24x7 concentra el 67,0 \% de las HH y la construcción e implantación suman el 23,5 \%. El 9,5 \% restante incluye dirección, requerimientos, diseño, certificación e innovaciones. La continuidad del servicio determina la mayor parte de la dotación.

\figura{16.09}{fig-B-2-edt-sitco-aranda-ramas-1}{EDT de SITCO Aranda, ramas 1.1 y 1.2.}{fig:edt-sitco-aranda-ramas-1-1}{Fuente: elaboración propia a partir de la Guía del PMBOK (PMI, 2017, pp. 156–162).}
La EDT separa la gestión del proyecto de los requisitos de cada etapa. Los acuerdos con terceros se controlan como entregables para identificar qué construcción depende de ellos.

\figura{16.09}{fig-B-3-edt-sitco-aranda-ramas-1}{EDT de SITCO Aranda, ramas 1.3 y 1.4.}{fig:edt-sitco-aranda-ramas-1-3}{Fuente: elaboración propia a partir de la Guía del PMBOK (PMI, 2017, pp. 156–162).}
El diseño y la construcción tienen controles separados. Los módulos se agrupan por etapa, con entregables propios para sus canales e integraciones.

\figura{16.09}{fig-B-4-edt-sitco-aranda-ramas-1}{EDT de SITCO Aranda, ramas 1.5 y 1.6.}{fig:edt-sitco-aranda-ramas-1-5}{Fuente: elaboración propia a partir de la Guía del PMBOK (PMI, 2017, pp. 156–162).}
Cada paso a producción requiere la certificación previa. La implantación habilita la infraestructura y prepara a los usuarios mediante capacitación y acompañamiento en faena.

\figura{16.09}{fig-B-5-edt-sitco-aranda-ramas-1}{EDT de SITCO Aranda, ramas 1.7 y 1.8.}{fig:edt-sitco-aranda-ramas-1-7}{Fuente: elaboración propia a partir de la Guía del PMBOK (PMI, 2017, pp. 156–162).}
Las innovaciones se verifican por separado. Los diez servicios de operación generan la evidencia del pago mensual.

La secuencia prioriza trazabilidad y conciliación, terreno y personas, y analítica de mantenimiento (Compañía Minera Altos de Aranda [CMAA], 2026b, sección 13.1). La telemetría requiere el acuerdo con tres fabricantes (1.2.4.2); la verificación facial, el acuerdo sindical (1.1.5.2), y Cumplimiento Ambiental espera el caudalímetro de marzo de 2028. Estas dependencias obligan a iniciar las marchas blancas de las seis áreas por separado.

\apartado{Carta Gantt}
La Figura \ref{fig:carta-gantt-resumida-56-meses-contrato} reúne los 56 meses, las fases, ocho ramas, doce hitos, pagos mensuales de operación y detenciones de marzo y septiembre. El mes 1 se fija en enero de 2027; un inicio real distinto desplaza las fechas, conservando los meses contractuales.

\begin{horizontal}
\figura{22.42}{fig-C-1-carta-gantt-resumida-56-meses}{Carta Gantt resumida de los 56 meses del contrato.}{fig:carta-gantt-resumida-56-meses-contrato}{Fuente: elaboración propia a partir de la Guía del PMBOK (PMI, 2017, pp. 217–218) y de las Bases administrativas (CMAA, 2026a, Art. 17.1, p. 12).}
El desarrollo de la Etapa 2 coincide con la marcha blanca de la Etapa 1 en los meses 13 a 15. Su propia marcha blanca, meses 19 a 20, transcurre con la Etapa 1 en producción. La Figura \ref{fig:carta-gantt-resumida-56-meses-contrato} sitúa los pasos de los meses 16 y 21 y el inicio de Operación.
\end{horizontal}

Los trabajos que detienen procesos se hacen en las detenciones de planta. El diodo se conecta en septiembre de 2027. En marzo de 2028 se instala el caudalímetro (1.6.2.2) y se prueban la conmutación y el corte de 48 horas de la Etapa 1; en septiembre de 2028, los de la Etapa 2.

Cada ola de paso a producción se ejecuta entre las 03:00 y las 05:00 (CMAA, 2026b, RT-10.05), con al menos 72 horas entre olas. Va después del quinto día hábil, para no coincidir con el cierre contable, y fuera de las detenciones, de los embarques y de las 48 horas previas a cada uno. El Comité de Proyecto aprueba el paso y el acta acredita H7 o H12 (Subdocumento 9; Formulario T-17).

H1 ajusta el calendario a las detenciones y embarques reales. Para planificar, las olas se fijan en 10, 17 y 21 de abril y 20, 24 y 28 de septiembre de 2028. Las detenciones posteriores se fijan en el plan anual (1.8.4.1).

Si un embarque impide una ola, se usa la primera ventana posterior que cumpla estas condiciones, sin mover H12 del mes 21. Un atraso mayor queda sujeto al Art. 80°. La marcha blanca que no cumple el Art. 17.3 se extiende a costo de Overfore.

La Operación comienza con la primera ola del 20 de septiembre de 2028 (Art. 17.2.4). Desde ese día, los puestos 24x7 de la rama 1.8 atienden la Etapa 1 con los niveles del Art. 78.2.

Las firmas, talleres y validaciones con personal clave se programan fuera de los seis días de detención (Bases Técnicas, cap. 10, restricción 11). Diodo, caudalímetro, conmutación y corte se prueban durante las detenciones, para no intervenir la producción (Subdocumento 9; T-13).

La ruta crítica es la secuencia de actividades que fija el plazo. La Tabla \ref{tab:red-precedencias-calculo-ruta-critica-meses} la calcula con el método CPM, en meses del contrato (PMI, 2017). ES y EF son el inicio y el término más tempranos; LS y LF, los más tardíos, y la holgura, LS − ES, es el atraso que admite cada actividad. FC, CC y FF indican las relaciones fin-comienzo, comienzo-comienzo y fin-fin.

\begin{tabla}{I{0.274}I{0.188}D{0.125}D{0.074}D{0.075}D{0.074}D{0.075}D{0.113}}{Red de precedencias y cálculo de la ruta crítica, en meses del contrato.}{tab:red-precedencias-calculo-ruta-critica-meses}
  {\cab{Actividad (paquetes)} & \cab{Predecesoras} & \cab{Duración} & \cab{ES} & \cab{EF} & \cab{LS} & \cab{LF} & \cab{Holgura}}
  {Fuente: elaboración propia a partir de la Guía del PMBOK (PMI, 2017, pp. 210–211) y de las Bases administrativas (CMAA, 2026a, Art. 17.1, p. 12).}
A01 Primera acreditación del equipo clave en faena (1.1.7.1) & Ninguna & 1 & 1 & 1 & 2 & 2 & 1 \\
A02 Análisis de interesados (1.1.5.3) & Ninguna & 1 & 1 & 1 & 1 & 1 & 0 \\
A03 Inventario de los diez sistemas e interfaces (1.2.1.1 y 1.2.1.2) & A01 (CC) & 3 & 1 & 3 & 2 & 4 & 1 \\
A04 Línea base de la Etapa 1 (1.2.2.1 a 1.2.2.3) & Ninguna & 2 & 1 & 2 & 1 & 2 & 0 \\
A05 Decisiones del 16.1 ratificadas (1.2.3.1 y 1.2.3.2) & Ninguna & 3 & 1 & 3 & 1 & 3 & 0 \\
A06 Arquitectura, datos y seguridad (1.3.1.1, 1.3.2.1 y 1.3.3.1) & A04 (CC+1), A03 (FF), A05 (FF+1) & 3 & 2 & 4 & 2 & 4 & 0 \\
A07 Especificación de acceso a la telemetría (1.2.4.1) & A03 (CC) & 3 & 1 & 3 & 5 & 7 & 4 \\
A08 Acuerdos con los tres fabricantes (1.2.4.2) & A07 (CC+2) & 10 & 3 & 12 & 7 & 16 & 4 \\
A09 Acuerdo con el sindicato (1.1.5.2) & A02 (FC) & 13 & 2 & 14 & 2 & 14 & 0 \\
A10 Ambientes en la nube (1.6.1.1) & A06 (CC) & 5 & 2 & 6 & 2 & 6 & 0 \\
A11 Contrato y recepción del diodo de datos (1.1.6.2 y 1.6.1.3) & A06 (CC) & 4 & 2 & 5 & 5 & 8 & 3 \\
A12 Conexión del diodo (detención de sep. 2027) (1.6.1.3) & A11 (FC); fecha fija & 1 & 9 & 9 & 9 & 9 & 0 \\
A13 Factores de ajuste explicados (1.2.6.1) & A04 (CC+1) & 1 & 2 & 2 & 5 & 5 & 3 \\
A14 Plataformas, módulos y canales de la Etapa 1 (1.4.2.1, 1.4.2.2, 1.4.3.1 a 1.4.3.3, 1.4.3.5, 1.4.3.6 y 1.4.5) & A06 (CC+1), A10 (CC+1) & 8 & 3 & 10 & 3 & 10 & 0 \\
A15 Conciliación Metalúrgica (1.4.3.4) & A13 (FC), A06 (FC) & 5 & 5 & 9 & 6 & 10 & 1 \\
A16 Corrección de los defectos de las pruebas (1.4.2.1, 1.4.2.2, 1.4.3.2 a 1.4.3.4, 1.4.3.6, 1.4.5.1 y 1.4.5.2, mes 11) & A14 (FC), A15 (FC), A12 (FC) & 1 & 11 & 11 & 11 & 11 & 0 \\
A17 Certificación de la Etapa 1 (1.5.3.1 y 1.5.3.3) & A16 (FC) & 1 & 12 & 12 & 12 & 12 & 0 \\
A18 Perfilado y saneamiento (1.4.6.1) & A06 (CC+1) & 6 & 3 & 8 & 3 & 8 & 0 \\
A19 Carga de los datos de la Etapa 1 (1.4.6.2) & A18 (FC) & 7 & 9 & 15 & 9 & 15 & 0 \\
A20 Validación de prototipos con guantes (1.3.4.3) & A06 (CC+2) & 4 & 4 & 7 & 4 & 7 & 0 \\
A21 Dispositivos de la Etapa 1 instalados (1.6.2.1) & A20 (FC) & 5 & 8 & 12 & 8 & 12 & 0 \\
A22 Marcha blanca de la Etapa 1 (1.6.3.1) & A17 (FC), A21 (FC) & 3 & 13 & 15 & 13 & 15 & 0 \\
A23 Corte de 48 h y conmutación (detención de mar. 2028) (1.5.5.1 y 1.5.5.2) & A17 (FC); fecha fija & 1 & 15 & 15 & 15 & 15 & 0 \\
A24 Paso a producción de la Etapa 1 (1.6.4.1) & A22 (FC), A23 (FC), A19 (FC) & 1 & 16 & 16 & 16 & 16 & 0 \\
B01 Línea base y diseño de la Etapa 2 (1.2.5.1, 1.3.5.1 y 1.3.5.2) & A17 (FC) & 2 & 13 & 14 & 13 & 14 & 0 \\
B02 Conectores y módulos de la Etapa 2 (1.4.2.3, 1.4.4.1 a 1.4.4.3, 1.4.4.5 y 1.4.4.6) & B01 (CC) & 5 & 13 & 17 & 13 & 17 & 0 \\
B03 Conexión productiva de la telemetría (1.4.2.3 y 1.4.4.1) & A08 (FC), B02 (FF) & 1 & 17 & 17 & 17 & 17 & 0 \\
B04 Verificación facial en portería (1.4.4.4) & A09 (FC) & 3 & 15 & 17 & 15 & 17 & 0 \\
B05 Caudalímetro (detención de mar. 2028) (1.6.2.2) & B01 (FC); fecha fija & 1 & 15 & 15 & 15 & 15 & 0 \\
B06 Digitalización de los formularios de SSO (1.4.6.3) & A19 (CC+1), A18 (FC) & 6 & 10 & 15 & 10 & 15 & 0 \\
B07 Carga de los datos de la Etapa 2 (1.4.6.4) & B06 (FC) & 5 & 16 & 20 & 16 & 20 & 0 \\
B08 Certificación de la Etapa 2 (1.5.6.1 y 1.5.6.3) & B02 (FC), B03 (FC), B04 (FC), B05 (FC) & 1 & 18 & 18 & 18 & 18 & 0 \\
B09 Marcha blanca de la Etapa 2 (1.6.3.2) & B08 (FC) & 2 & 19 & 20 & 19 & 20 & 0 \\
B10 Corte de 48 h y conmutación (detención de sep. 2028) (1.5.5.1 y 1.5.5.2) & B09 (FC); fecha fija & 1 & 21 & 21 & 21 & 21 & 0 \\
B11 Paso a producción de la Etapa 2 (1.6.4.2) & B10 (CC), B07 (FC) & 1 & 21 & 21 & 21 & 21 & 0 \\
\end{tabla}
Las actividades con holgura cero ocupan todo el plazo entre hitos (Art. 17°; E-25). Los acuerdos con fabricantes permiten construir contra la especificación desde el mes 13, pero la conexión productiva del mes 17 espera la firma. La verificación facial depende del acuerdo sindical.

La Tabla \ref{tab:estimacion-pert-duracion-actividades-inciertas-red} estima con PERT las duraciones que dependen de terceros y la probabilidad de cumplir el plazo que deja la red.

\begin{tabla}{I{0.251}I{0.146}I{0.179}I{0.186}I{0.238}}{Estimación PERT de la duración de las actividades inciertas de la red.}{tab:estimacion-pert-duracion-actividades-inciertas-red}
  {\cab{Actividad} & \cab{O, M y P} & \cab{Duración esperada y σ} & \cab{Plazo que deja la red} & \cab{Probabilidad de cumplirlo}}
  {Fuente: elaboración propia a partir de la Guía del PMBOK (PMI, 2017, pp. 201–202).}
A01 Primera acreditación del equipo clave & 7, 9 y 15 días hábiles & 9,7 y 1,3 días hábiles & Meses 1 a 2 de la red & más de 99,9 \% (Z = 4,00) \\
A03 Inventario de los diez sistemas e interfaces & 8, 12 y 18 semanas & 12,3 y 1,7 semanas & 17,3 semanas, meses 1 a 4 & 99,9 \% (Z = 3,00) \\
A08 Acuerdos con los tres fabricantes & 20, 32 y 52 semanas & 33,3 y 5,3 semanas & 60,7 semanas, meses 3 a 16 & más de 99,9 \% (Z = 5,12) \\
A09 Acuerdo con el sindicato & 12, 24 y 48 semanas & 26,0 y 6,0 semanas & 56,3 semanas, meses 2 a 14 & más de 99,9 \% (Z = 5,06) \\
B06 Digitalización con tres puestos & 22, 24,8 y 30 semanas & 25,2 y 1,3 semanas & 25,6 semanas, los 128 días hábiles de los meses 10 a 15 & 61,8 \% (Z = 0,30) \\
\end{tabla}
Con tres puestos, la digitalización tiene 61,8 \% de probabilidad de terminar en 128 días hábiles; con cuatro, la duración esperada baja a 18,9 semanas. En las negociaciones, el riesgo es el rechazo del fabricante o del sindicato. Se mantienen estas contingencias:

\begin{itemize}
\item Diodo de datos: si se pierde la detención de septiembre de 2027, la siguiente es marzo de 2028, y hasta entonces el conector del historiador se prueba con su simulador y con extractos (1.4.2.1).
\item Telemetría: sin acuerdo, el conector se construye contra la especificación y la conexión productiva espera la firma (1.2.4.2).
\item Sindicato: sin acuerdo al inicio del mes 15, la portería sigue con la credencial de proximidad (1.1.5.2).
\item Digitalización: si el lote piloto rinde menos de lo previsto, el cuarto digitalizador abre un cuarto puesto (1.4.6.3).
\end{itemize}
Las cartas Gantt representan la red de precedencias. En la carta resumida, las flechas encadenan las fases del Art. 17° y el inicio de la Operación. En las cartas de las ramas, una flecha une cada paquete con su sucesor dentro de la misma figura, según la red y las dependencias que declara el diccionario de la EDT, y otra une la barra con el hito que gatilla. Un recuadro junto a la barra da el predecesor que está en otra figura, con tipo y desfase si no es fin a comienzo (CC+1); en rojo, a la derecha, la sucesora crítica de otra figura. Las barras con borde rojo corresponden a actividades con holgura cero, que forman la ruta crítica, y los hitos con recuadro rojo son sus cierres: H5, H7, H10 y H12. La ruta crítica termina en H12, en el mes 21, porque la Operación dura 36 meses fijos.

\begin{horizontal}
La rama 1.1, dirección del proyecto, se presenta en las figuras \ref{fig:carta-gantt-rama-1-1-entregables}, \ref{fig:carta-gantt-rama-1-1-entregables-2} y \ref{fig:carta-gantt-rama-1-1-entregables-3}.

\figura{22.42}{fig-C-2-carta-gantt-rama-1-1}{Carta Gantt de la rama 1.1, entregables 1.1.1 a 1.1.3.}{fig:carta-gantt-rama-1-1-entregables}{Fuente: elaboración propia a partir de la Guía del PMBOK (PMI, 2017, pp. 217–219) y de las Bases administrativas (CMAA, 2026a, Art. 17.1, p. 12).}
La dirección continúa hasta el mes 21, mientras el seguimiento del avance y de los riesgos llega al mes 20. El registro de cambios acompaña ese seguimiento para evaluar cómo cada ajuste afecta el plazo y las horas previstas.

\figura{22.42}{fig-C-3-carta-gantt-rama-1-1}{Carta Gantt de la rama 1.1, entregables 1.1.4 a 1.1.6.}{fig:carta-gantt-rama-1-1-entregables-2}{Fuente: elaboración propia a partir de la Guía del PMBOK (PMI, 2017, pp. 217–219) y de las Bases administrativas (CMAA, 2026a, Art. 17.1, p. 12).}
Los criterios de calidad y aceptación se preparan al inicio y se revisan durante la implementación. El acuerdo sindical debe preceder a la construcción de la verificación facial, y las adquisiciones deben terminar antes de utilizar cada bien; un atraso en esos acuerdos afecta el trabajo que depende de ellos.

\figura{22.42}{fig-C-4-carta-gantt-rama-1-1}{Carta Gantt de la rama 1.1, entregables 1.1.7 a 1.1.10.}{fig:carta-gantt-rama-1-1-entregables-3}{Fuente: elaboración propia a partir de la Guía del PMBOK (PMI, 2017, pp. 217–219) y de las Bases administrativas (CMAA, 2026a, Art. 17.1, p. 12).}
Los cierres de etapa reúnen las lecciones aprendidas y los libros de operación que se entregan al cliente. La salida se prepara desde el plan de reversibilidad, entregado dentro de noventa días y actualizado anualmente durante el contrato.

La rama 1.2, requerimientos del producto, se presenta en las figuras \ref{fig:carta-gantt-rama-1-2-entregables} y \ref{fig:carta-gantt-rama-1-2-entregables-2}.

\figura{22.42}{fig-C-5-carta-gantt-rama-1-2}{Carta Gantt de la rama 1.2, entregables 1.2.1 a 1.2.3.}{fig:carta-gantt-rama-1-2-entregables}{Fuente: elaboración propia a partir de la Guía del PMBOK (PMI, 2017, pp. 217–219) y de las Bases administrativas (CMAA, 2026a, Art. 17.1, p. 12).}
El inventario de sistemas y la matriz de integraciones sustentan las decisiones que el cliente ratifica en H2. Los criterios de aceptación de la Etapa 1 se fijan al inicio, de modo que las revisiones de desarrollo contrasten cada entrega con un resultado ya acordado.

\figura{22.42}{fig-C-6-carta-gantt-rama-1-2}{Carta Gantt de la rama 1.2, entregables 1.2.4 a 1.2.6.}{fig:carta-gantt-rama-1-2-entregables-2}{Fuente: elaboración propia a partir de la Guía del PMBOK (PMI, 2017, pp. 217–219) y de las Bases administrativas (CMAA, 2026a, Art. 17.1, p. 12).}
Los factores de ajuste se documentan al inicio con la persona aún empleada y, si es posible, la jubilada. La negociación de acceso a la telemetría avanza en paralelo. La Etapa 2 mantiene sus revisiones hasta el mes 21, pero desde el mes 19 admite sólo defectos y ajustes de marcha blanca.

La rama 1.3, diseño del sistema, se presenta en las figuras \ref{fig:carta-gantt-rama-1-3-entregables} y \ref{fig:carta-gantt-rama-1-3-entregables-2}.

\figura{22.42}{fig-C-7-carta-gantt-rama-1-3}{Carta Gantt de la rama 1.3, entregables 1.3.1 a 1.3.3.}{fig:carta-gantt-rama-1-3-entregables}{Fuente: elaboración propia a partir de la Guía del PMBOK (PMI, 2017, pp. 217–219) y de las Bases administrativas (CMAA, 2026a, Art. 17.1, p. 12).}
Arquitectura, datos y seguridad se entregan en H2 para dar una base común a la construcción. El estudio de cobertura vuelve a revisarse antes de las detenciones de la implementación, porque las zonas de sombra deben contrastarse con los cambios previstos en el rajo.

\figura{22.42}{fig-C-8-carta-gantt-rama-1-3}{Carta Gantt de la rama 1.3, entregables 1.3.4 a 1.3.6.}{fig:carta-gantt-rama-1-3-entregables-2}{Fuente: elaboración propia a partir de la Guía del PMBOK (PMI, 2017, pp. 217–219) y de las Bases administrativas (CMAA, 2026a, Art. 17.1, p. 12).}
La observación de los turnos alimenta los prototipos que se validan con los usuarios de terreno. La Etapa 2 se diseña sobre la plataforma de la Etapa 1; sus nuevos módulos y conectores se definen durante la primera marcha blanca.

La rama 1.4, software del sistema, se presenta en las figuras \ref{fig:carta-gantt-rama-1-4-entregables} y \ref{fig:carta-gantt-rama-1-4-entregables-2}.

\figura{22.42}{fig-C-9-carta-gantt-rama-1-4}{Carta Gantt de la rama 1.4, entregables 1.4.1 a 1.4.3.}{fig:carta-gantt-rama-1-4-entregables}{Fuente: elaboración propia a partir de la Guía del PMBOK (PMI, 2017, pp. 217–219) y de las Bases administrativas (CMAA, 2026a, Art. 17.1, p. 12).}
El repositorio y los controles de entrega se habilitan al inicio para revisar y desplegar las versiones que construyen los equipos. H4 recibe la versión disponible para pruebas; después continúa la corrección de defectos de la Etapa 1 durante su certificación, marcha blanca y estabilización.

\figura{22.42}{fig-C-10-carta-gantt-rama-1-4}{Carta Gantt de la rama 1.4, entregables 1.4.4 a 1.4.6.}{fig:carta-gantt-rama-1-4-entregables-2}{Fuente: elaboración propia a partir de la Guía del PMBOK (PMI, 2017, pp. 217–219) y de las Bases administrativas (CMAA, 2026a, Art. 17.1, p. 12).}
Los canales se construyen junto a sus módulos y la Etapa 2 entrega su versión para pruebas en H9. La migración se realiza por etapa: cada carga definitiva tiene dos ensayos completos y la conciliación continúa con los registros que se generan durante la marcha blanca.

La rama 1.5, sistema certificado, se presenta en las figuras \ref{fig:carta-gantt-rama-1-5-entregables} y \ref{fig:carta-gantt-rama-1-5-entregables-2}.

\figura{22.42}{fig-C-11-carta-gantt-rama-1-5}{Carta Gantt de la rama 1.5, entregables 1.5.1 a 1.5.3.}{fig:carta-gantt-rama-1-5-entregables}{Fuente: elaboración propia a partir de la Guía del PMBOK (PMI, 2017, pp. 217–219) y de las Bases administrativas (CMAA, 2026a, Art. 17.1, p. 12).}
Las pruebas automatizadas acompañan los ciclos de desarrollo, mientras la usabilidad se comprueba con personas en terreno. La certificación de la Etapa 1 reúne las pruebas de carga, resiliencia y aceptación antes de comenzar la marcha blanca.

\figura{22.42}{fig-C-12-carta-gantt-rama-1-5}{Carta Gantt de la rama 1.5, entregables 1.5.4 a 1.5.6.}{fig:carta-gantt-rama-1-5-entregables-2}{Fuente: elaboración propia a partir de la Guía del PMBOK (PMI, 2017, pp. 217–219) y de las Bases administrativas (CMAA, 2026a, Art. 17.1, p. 12).}
Las pruebas de intrusión independientes preceden a cada certificación y paso a producción. La conmutación y el corte de 48 horas verifican la recuperación y la operación sin enlace; la regresión de ambas etapas comprueba que las ampliaciones no alteren lo ya entregado.

La rama 1.6, solución implantada, se presenta en las figuras \ref{fig:carta-gantt-rama-1-6-entregables} y \ref{fig:carta-gantt-rama-1-6-entregables-2}.

\figura{22.42}{fig-C-13-carta-gantt-rama-1-6}{Carta Gantt de la rama 1.6, entregables 1.6.1 a 1.6.3.}{fig:carta-gantt-rama-1-6-entregables}{Fuente: elaboración propia a partir de la Guía del PMBOK (PMI, 2017, pp. 217–219) y de las Bases administrativas (CMAA, 2026a, Art. 17.1, p. 12).}
Los dispositivos quedan instalados antes de cada marcha blanca: Overfore los especifica y el cliente los compra. La instalación da paso a la validación en faena, con una persona por turno, de día y de noche, en cada área para acompañar a los usuarios.

\figura{22.42}{fig-C-14-carta-gantt-rama-1-6}{Carta Gantt de la rama 1.6, entregables 1.6.4 a 1.6.6.}{fig:carta-gantt-rama-1-6-entregables-2}{Fuente: elaboración propia a partir de la Guía del PMBOK (PMI, 2017, pp. 217–219) y de las Bases administrativas (CMAA, 2026a, Art. 17.1, p. 12).}
El soporte de la Etapa 1 continúa entre su estabilización y el inicio de Operación. Así se mantiene la atención de la etapa que ya está en producción mientras se completa la segunda.

La estabilización termina el 31 de mayo de 2028 en la Etapa 1 y el 8 de noviembre de 2028 en la Etapa 2. El personal de cada ola se capacita en las cuatro semanas previas a su marcha blanca. La Figura \ref{fig:carta-gantt-rama-1-7-innovaciones} presenta la rama 1.7, innovaciones.

\figura{22.42}{fig-C-15-carta-gantt-rama-1-7}{Carta Gantt de la rama 1.7, innovaciones.}{fig:carta-gantt-rama-1-7-innovaciones}{Fuente: elaboración propia a partir de la Guía del PMBOK (PMI, 2017, pp. 217–219) y de las Bases administrativas (CMAA, 2026a, Art. 17.1, p. 12).}
\end{horizontal}
El clasificador de detenciones se activa después de reunir seis meses de datos en cada área; hasta entonces se usan reglas. Las innovaciones 1, 2, 3 y 5 se certifican con H5 o H10, y la 4 con actas durante la Operación.

La rama 1.8, servicio de operación y mantención, se presenta en las figuras \ref{fig:carta-gantt-rama-1-8-entregables}, \ref{fig:carta-gantt-rama-1-8-entregables-2} y \ref{fig:carta-gantt-rama-1-8-entregables-3}.

\figura{16.09}{fig-C-16-carta-gantt-rama-1-8}{Carta Gantt de la rama 1.8, entregables 1.8.1 a 1.8.4.}{fig:carta-gantt-rama-1-8-entregables}{Fuente: elaboración propia a partir de la Guía del PMBOK (PMI, 2017, pp. 217–219) y de las Bases administrativas (CMAA, 2026a, Art. 17.1, p. 12).}
La mesa, el soporte en faena y los centros de red y seguridad se mantienen desde el mes 21 hasta el 56. La atención de incidentes se complementa con análisis de causa raíz y mantención preventiva para reducir su repetición.

\figura{16.09}{fig-C-17-carta-gantt-rama-1-8}{Carta Gantt de la rama 1.8, entregables 1.8.5 a 1.8.7.}{fig:carta-gantt-rama-1-8-entregables-2}{Fuente: elaboración propia a partir de la Guía del PMBOK (PMI, 2017, pp. 217–219) y de las Bases administrativas (CMAA, 2026a, Art. 17.1, p. 12).}
Las mejoras se priorizan con el cliente mediante el plan de mantención evolutiva. Los planes de continuidad se renuevan cada doce meses, y los informes de servicio y comités permiten revisar si los cambios resuelven los incumplimientos detectados durante la operación.

\figura{16.09}{fig-C-18-carta-gantt-rama-1-8}{Carta Gantt de la rama 1.8, entregables 1.8.8 a 1.8.10.}{fig:carta-gantt-rama-1-8-entregables-3}{Fuente: elaboración propia a partir de la Guía del PMBOK (PMI, 2017, pp. 217–219) y de las Bases administrativas (CMAA, 2026a, Art. 17.1, p. 12).}
Las pruebas y la capacitación continúan durante la Operación para comprobar la recuperación del sistema y preparar al personal que se incorpora. En los tres últimos meses se entregan código, datos, configuraciones y credenciales, junto con noventa días de acompañamiento al nuevo operador.

\apartado{Responsables}
Cada paquete tiene un responsable del equipo clave que responde por su entrega y aceptación, aunque participen otras personas. La Tabla \ref{tab:roles-equipo-clave-dedicacion-paquetes-cargo} identifica los diez roles mínimos de implementación y el Gerente de Servicio, sus titulares, dedicación, meses, paquetes y HH.

\begin{tabla}{I{0.084}I{0.186}I{0.162}I{0.237}I{0.104}D{0.130}D{0.098}}{Roles del equipo clave, dedicación y paquetes a cargo.}{tab:roles-equipo-clave-dedicacion-paquetes-cargo}
  {\cab{Sigla} & \cab{Rol} & \cab{Persona} & \cab{Dedicación mínima exigida} & \cab{Meses} & \cab{Paquetes} & \cab{HH en la EDT}}
  {Fuente: elaboración propia a partir del Subdocumento 1 (Overfore S.A., 2026a, sección 1.5), del Anexo E del Subdocumento 1 (Overfore S.A., 2026b, Tabla \refext{A1}{tab:roles-dedicacion-periodo-responsabilidad-equipo-clave} de los Anexos del Subdocumento 1) y de las Tablas \ref{tab:diccionario-entregable-1-1-1-acta} a \ref{tab:diccionario-entregable-1-8-10-acta}}
JP & Jefe de Proyecto & Claudio Cabello Vásquez & 100 \% en implementación & 1 a 21 & 22 & 3.248 \\
AS & Arquitecto de Solución & Benjamín Rodríguez Pavez & Alta en diseño, permanente en el Comité de Arquitectura & 1 a 56 & 13 & 1.176 \\
ESI & Encargado de Seguridad de la Información & David Martínez Canto & Permanente & 1 a 56 & 13 & 1.208 \\
LDA & Líder de Datos & Benjamín Velásquez Farías & Permanente en implementación; en la operación, sólo los informes de las innovaciones 1 y 5 & 1 a 54 & 16 & 1.220 \\
LD & Líder de Desarrollo & Gabriel Sanzana Cuibin & 100 \% en implementación & 1 a 56 & 22 & 4.258 \\
LF & Líder Funcional & Patricio Figueroa Gallardo & 100 \% en implementación & 1 a 21 & 18 & 3.312 \\
LCP & Líder de Calidad y Pruebas & Kavon Kermani Órdenes & Permanente & 1 a 56 & 16 & 1.320 \\
LI & Líder de Integración & Marcelo Flores Coloane & Permanente en implementación & 1 a 21 & 8 & 1.462 \\
SRE & Líder de Operación / SRE & Lucas Zamora Gatica & Desde el mes 6, permanente en Operación & 6 a 56 & 28 & 3.436 \\
LIG & Líder de Implantación y Gestión del Cambio & Benjamín Velásquez Farías & Desde el mes 8 (se incorpora en el mes 2) & 2 a 56 & 11 & 1.096 \\
GS & Gerente de Servicio & Claudio Cabello Vásquez & Dedicado durante la Operación & 21 a 56 & 8 & 5.384 \\
Total & 11 roles & 9 personas & No aplica & 1 a 56 & 175 & 27.120 \\
\end{tabla}
Nueve personas ejercen once roles: Claudio Cabello y Benjamín Velásquez asumen dos roles cada uno. La dedicación permanente implica participar todos los meses, según las horas de sus paquetes.

Claudio Cabello es Jefe de Proyecto con dedicación exclusiva hasta H12. Desde la primera ola del 20 de septiembre de 2028 también es Gerente de Servicio (Art. 17.2.4), y hasta el acta de H12 el Líder de Operación / SRE asume la operación diaria. Sus 152 HH del mes 21 se cuentan una sola vez.

Benjamín Velásquez combina Datos e Implantación y Gestión del Cambio desde el mes 2, con la dedicación exigida desde el mes 8 (CMAA, 2026c, 19.2). Sus horas combinadas no superan su capacidad mensual. Cada rol tiene suplente preparado (1.1.7.1).

La capacidad es 160 HH mensuales hasta marzo de 2028 y 152 desde abril, porque la Ley N° 21.561 (2023) reduce la jornada de 42 a 40 horas semanales el 26 de abril de 2028.

Los roles con dedicación completa llenan su jornada con un paquete continuo, que recibe las horas que no usan sus demás paquetes. El Formulario T-15 muestra la distribución mensual.

La Tabla \ref{tab:perfiles-apoyo-externos-participan-edt} identifica al personal de apoyo propio y especialistas externos, sus ramas y HH. Los externos se contratan solo en los paquetes y meses necesarios, con autorización previa y escrita del cliente.

\begin{tabla}{I{0.126}I{0.386}I{0.186}D{0.173}D{0.130}}{Perfiles de apoyo y externos que participan en la EDT.}{tab:perfiles-apoyo-externos-participan-edt}
  {\cab{Sigla} & \cab{Perfil} & \cab{Origen} & \cab{Ramas} & \cab{HH}}
  {Fuente: elaboración propia a partir del Subdocumento 1 (Overfore S.A., 2026a, sección 1.5), del Anexo F del Subdocumento 1 (Overfore S.A., 2026b, Tablas \refext{A1}{tab:dotacion-apoyo-overfore-asignada-proyecto} y \refext{A1}{tab:especialistas-externos} de los Anexos del Subdocumento 1) y de las Tablas \ref{tab:diccionario-entregable-1-1-1-acta} a \ref{tab:diccionario-entregable-1-8-10-acta}}
DEV & Desarrollador & Apoyo Overfore & 1.2, 1.3, 1.4, 1.5, 1.6, 1.7, 1.8 & 17.994 \\
QA & Analista de pruebas & Apoyo Overfore & 1.1, 1.2, 1.5, 1.8 & 2.024 \\
ADT & Analista de datos & Apoyo Overfore & 1.2, 1.3, 1.4, 1.6, 1.7, 1.8 & 2.860 \\
OPS & Analista de operación de la mesa de ayuda y del centro de operaciones de red (NOC) & Apoyo Overfore & 1.1, 1.5, 1.6, 1.8 & 63.316 \\
ASG & Analista del centro de operaciones de seguridad (SOC) & Apoyo Overfore & 1.3, 1.5, 1.8 & 26.384 \\
DIG & Digitalizador & Apoyo Overfore & 1.4 & 2.976 \\
TEC & Técnico de terreno & Apoyo Overfore & 1.5, 1.6, 1.8 & 35.500 \\
ACP & Analista de control de proyectos & Externo & 1.1 & 608 \\
ABG & Abogado en datos personales y derecho laboral & Externo & 1.1, 1.2, 1.3, 1.8 & 344 \\
UXR & Diseñador UX/UI & Externo & 1.3, 1.4, 1.5 & 640 \\
RFE & Especialista en radiofrecuencia & Externo & 1.3, 1.8 & 192 \\
PEN & Auditor de intrusión independiente & Externo & 1.5, 1.8 & 400 \\
MON & Montajista eléctrico y de comunicaciones & Externo & 1.6 & 192 \\
REL & Relator de capacitación & Externo & 1.6, 1.8 & 1.176 \\
RTC & Redactor técnico & Externo & 1.6 & 216 \\
CDA & Científico de datos & Externo & 1.7 & 552 \\
BCP & Especialista en continuidad del negocio & Externo & 1.3, 1.8 & 128 \\
GCA & Consultor en gestión del cambio & Externo & 1.6 & 228 \\
PRF & Especialista en pruebas de rendimiento & Externo & 1.5, 1.8 & 216 \\
\end{tabla}
La operación 24x7 y el acompañamiento concentran la mayor parte de las HH de apoyo. La Tabla \ref{tab:esfuerzo-estimado-rama-origen-horas} distribuye el esfuerzo por rama y separa HH propias y externas.

\begin{tabla}{I{0.314}D{0.133}D{0.132}I{0.285}D{0.136}}{Esfuerzo estimado por rama y origen de las horas.}{tab:esfuerzo-estimado-rama-origen-horas}
  {\cab{Rama} & \cab{Equipo clave (HH)} & \cab{Apoyo Propio (HH)} & \cab{Externos: total HH y perfiles} & \cab{Total}}
  {Fuente: elaboración propia a partir de las Tablas \ref{tab:roles-equipo-clave-dedicacion-paquetes-cargo}, \ref{tab:perfiles-apoyo-externos-participan-edt} y \ref{tab:diccionario-entregable-1-1-1-acta} a \ref{tab:diccionario-entregable-1-8-10-acta}}
1.1 Dirección del proyecto & 4.168 & 104 & 816; ACP 608, ABG 208 & 5.088 \\
1.2 Requerimientos del producto & 3.308 & 352 & 24; ABG 24 & 3.684 \\
1.3 Diseño del sistema & 1.152 & 192 & 456; UXR 248, RFE 112, ABG 64, BCP 32 & 1.800 \\
1.4 Software del sistema & 4.492 & 16.290 & 280; UXR 280 & 21.062 \\
1.5 Sistema certificado & 1.296 & 2.032 & 512; PEN 256, PRF 144, UXR 112 & 3.840 \\
1.6 Solución implantada & 2.108 & 18.096 & 1716; REL 1.080, GCA 228, RTC 216, MON 192 & 21.920 \\
1.7 Innovaciones & 1.016 & 1.432 & 552; CDA 552 & 3.000 \\
1.8 Servicio de operación y mantención & 9.580 & 112.556 & 536; PEN 144, BCP 96, REL 96, RFE 80, PRF 72, ABG 48 & 122.672 \\
Total & 27.120 & 151.054 & 4892; No aplica & 183.066 \\
\end{tabla}
Las 4.892 HH externas (2,7 \%) se destinan a especialidades que se requieren en tareas específicas. Desarrollo, migración, acompañamiento y operación 24x7 usan personal propio por su continuidad.

La dotación de apoyo coincide con la del Subdocumento 1, contada en el mes de mayor carga de cada período:

\begin{itemize}
\item Desarrolladores: ocho en los meses 3 a 11, siete en los 13 a 17, dos en los 12 y 18 a 20, y uno en la operación.
\item Pruebas y datos: un analista de pruebas en los meses 1 a 20, un analista de datos en los meses 2 a 12 y 18 a 20, y dos en los 13 a 17, en la operación, parte de la jornada de cada uno.
\end{itemize}
•Digitalizadores: cuatro en los meses 10 a 15, rotando en tres puestos.

\begin{itemize}
\item Analistas de operación: ocho en las semanas 6 a 13 de la marcha blanca de la Etapa 1, ocho técnicos de terreno en los meses 13 a 15, diez analistas y diez técnicos en julio y agosto de 2028, y en la operación diez analistas para la mesa y el NOC y cinco técnicos.
\item Analistas del SOC: cinco desde el mes 21.
\end{itemize}
Cada puesto 24x7 en turnos 7x7 necesita cuatro personas con jornada de 42 horas y cinco con 40 (1.1.7.2).

\apartado{Hitos del contrato y paquetes que los gatillan}
La Tabla \ref{tab:hitos-formulario-25-paquetes-edt-gatillan} vincula los doce hitos E-25 con los paquetes que los acreditan. El T-17 establece su protocolo de aceptación.

\begin{tabla}{I{0.080}D{0.079}I{0.393}D{0.448}}{Hitos del Formulario E-25 y paquetes de la EDT que los gatillan.}{tab:hitos-formulario-25-paquetes-edt-gatillan}
  {\cab{Hito} & \cab{Mes} & \cab{Entregable que lo gatilla} & \cab{Paquetes de la EDT}}
  {Fuente: elaboración propia a partir de las Bases administrativas (CMAA, 2026a, Formulario E-25)}
H1 & 2 & Línea base de alcance y matriz de trazabilidad aprobadas & 1.1.1.2, 1.1.2.1, 1.1.4.1, 1.1.4.3, 1.1.4.4, 1.2.2.1, 1.2.2.2, 1.2.2.3 \\
H2 & 4 & Arquitectura, plan de seguridad y modelo de datos aprobados & 1.2.1.1, 1.2.1.2, 1.2.3.1, 1.2.3.2, 1.3.1.1, 1.3.1.2, 1.3.1.3, 1.3.1.5, 1.3.2.1, 1.3.2.2, 1.3.2.3, 1.3.3.1, 1.3.3.2, 1.3.3.3, 1.3.3.4, 1.3.6.1, 1.3.6.2 \\
H3 & 6 & Infraestructura híbrida y ambientes habilitados, con observabilidad & 1.1.4.2, 1.1.6.1, 1.4.1.1, 1.4.1.2, 1.4.1.3, 1.6.1.1, 1.6.1.2 \\
H4 & 10 & Software de la Etapa 1 entregado para pruebas, con QA superado y evidencia de cobertura & 1.2.2.4, 1.2.2.5, 1.2.2.6, 1.4.1.4, 1.4.2.1, 1.4.2.2, 1.4.3.1, 1.4.3.2, 1.4.3.3, 1.4.3.4, 1.4.3.5, 1.4.3.6, 1.4.5.1, 1.4.5.2, 1.5.1.1, 1.5.1.2, 1.5.1.3 \\
H5 & 12 & Solución de la Etapa 1 certificada & 1.1.4.4, 1.2.2.4, 1.4.3.7, 1.5.2.1, 1.5.3.1, 1.5.3.2, 1.5.3.3, 1.5.4.1, 1.7.1.1, 1.7.3.1, 1.7.3.2 \\
H6 & 13 & Inicio de la marcha blanca de la Etapa 1, con plan de reversión activo y usuarios capacitados (la primera ola, portería, con dos ciclos del 7x7; Formulario T-17) & 1.6.2.1, 1.6.2.3, 1.6.3.1, 1.6.3.3, 1.6.5.1 \\
H7 & 16 & Paso a producción de la Etapa 1 y cierre de su marcha blanca & 1.1.10.2, 1.3.1.4, 1.4.3.7, 1.4.6.2, 1.4.6.5, 1.5.2.1, 1.5.4.1, 1.5.4.2, 1.5.5.1, 1.5.5.2, 1.6.3.1, 1.6.3.3, 1.6.3.5, 1.6.4.1, 1.6.5.1 \\
H8 & 14 & Línea base y diseño detallado de la Etapa 2 aprobados & 1.2.5.1, 1.2.5.2, 1.2.5.3, 1.3.5.1, 1.3.5.2 \\
H9 & 17 & Software de la Etapa 2 entregado para pruebas, con QA superado & 1.1.5.2, 1.2.4.2, 1.2.5.4, 1.4.1.4, 1.4.2.3, 1.4.4.1, 1.4.4.2, 1.4.4.3, 1.4.4.4, 1.4.4.5, 1.4.4.6, 1.5.1.1, 1.5.1.2, 1.5.1.3, 1.5.6.3 \\
H10 & 18 & Solución de la Etapa 2 certificada y cierre del desarrollo & 1.2.5.4, 1.4.4.7, 1.4.6.3, 1.5.1.1, 1.5.1.2, 1.5.1.3, 1.5.4.1, 1.5.6.1, 1.5.6.2, 1.5.6.3, 1.7.2.1, 1.7.5.1 \\
H11 & 19 & Inicio de la marcha blanca de la Etapa 2 & 1.6.2.2, 1.6.3.2, 1.6.3.4, 1.6.5.1 \\
H12 & 21 & Paso a producción de la Etapa 2 y aceptación final & 1.1.4.4, 1.1.6.4, 1.1.10.2, 1.4.4.7, 1.4.6.4, 1.4.6.5, 1.5.2.2, 1.5.4.1, 1.5.5.1, 1.5.5.2, 1.6.3.2, 1.6.3.4, 1.6.4.2, 1.6.4.3 \\
\end{tabla}
Cada hito recibe la versión disponible en su mes. H4 y H9 aceptan la versión de los sprints en el ambiente de pruebas, y las correcciones de los meses 11 y 18 se certifican en H5 y H10. El diodo se acepta con acta propia en el mes 9. La rama 1.8 acredita los 36 pagos mensuales de los meses 21 a 56.

\apartado{Exclusiones del alcance}
La Tabla \ref{tab:elementos-excluidos-trabajo-si-queda-edt} distingue las actividades excluidas del alcance y los trabajos de especificación, coordinación o instalación que Overfore realiza, con sus paquetes asociados.

\begin{tabla}{I{0.398}I{0.426}D{0.176}}{Elementos excluidos y trabajo que sí queda en la EDT.}{tab:elementos-excluidos-trabajo-si-queda-edt}
  {\cab{Elemento excluido} & \cab{Trabajo incluido en la EDT} & \cab{Paquete}}
  {Fuente: elaboración propia a partir de las Bases técnicas del caso Minería (CMAA, 2026b, capítulo 11)}
Reemplazar SAP S/4HANA, el despacho de flota, el laboratorio o la planificación minera & Integración por las interfaces que ya ofrecen, sin modificarlos. Con el laboratorio, en ambos sentidos & 1.4.2.1, 1.4.2.3 \\
Intervenir el control de proceso de la planta o modificar su lógica & Lectura del historiador a través del diodo de datos & 1.4.2.1, 1.6.1.3 \\
Construir la red de comunicaciones del rajo & Estudio de cobertura y especificación de lo que falta & 1.3.1.5 \\
Comprar el hardware de terreno & Especificación, instalación, configuración y gestión & 1.6.2.1, 1.6.2.2 \\
Automatizar equipos mineros (operación autónoma de camiones o perforadoras) & Lectura de la telemetría de los equipos que entregan los fabricantes & 1.2.4.1, 1.2.4.2, 1.4.4.1 \\
Diseñar el plan minero o estimar recursos y reservas & Integración con el software de planificación, que se mantiene & 1.4.2.1 \\
Gestionar remuneraciones o administrar personal, que residen en el ERP corporativo, SAP S/4HANA & Acreditación y control de acceso de personas, sin funciones de remuneraciones & 1.4.3.5 \\
Operar el terminal portuario, que administra un tercero & Intercambio de los datos de despacho y embarque con el operador del puerto & 1.4.3.3 \\
\end{tabla}
El cliente ejecuta la obra civil de sala y gabinetes (1.6.1.2) y contrata la telemetría con los fabricantes, como dueño de los equipos. Overfore define los datos y acompaña los acuerdos (1.2.4.1 y 1.2.4.2). Torniquetes y cámaras de fatiga existentes se integran sin reemplazo.

\apartado{Rama 1.1 Dirección del proyecto}
La rama 1.1 tiene 10 entregables y 29 paquetes de trabajo, que suman 5.088 HH. Las tablas \ref{tab:diccionario-entregable-1-1-1-acta} a \ref{tab:diccionario-entregable-1-1-10-documentacion} presentan la ficha de cada entregable y, bajo ella, la de cada uno de sus paquetes.

\begin{tabla}{I{0.269}I{0.731}}{Diccionario del entregable 1.1.1, Acta de constitución y plan para la dirección.}{tab:diccionario-entregable-1-1-1-acta}
  {\cab{1.1.1} & \cab{Acta de constitución y plan para la dirección}}
  {Fuente: elaboración propia a partir de la Guía del PMBOK (PMI, 2017, p. 162).}
Descripción del trabajo & Acta de constitución y plan para la dirección, con el espacio de consulta del cliente. Supone que la Contraparte Técnica y el Comité Ejecutivo quedan designados en la primera semana. \\
Responsable & Claudio Cabello (JP). \\
Hitos asociados & M1–M21. Aporta a H1 (mes 2). \\
Esfuerzo estimado & 1.464 HH, 40 externas (JP 1.416, ESI 8 y ACP externo 40). \\
1.1.1.1 & Acta de constitución del proyecto \\
Descripción del trabajo & Redactar el acta, revisarla con la Contraparte Técnica y firmarla en la primera sesión del Comité Ejecutivo. \\
Entregable & Acta que formaliza el proyecto SITCO Aranda, con sus doce hitos. \\
Criterio de aceptación & Acta firmada por el patrocinador del cliente y por la gerencia de Overfore, con los doce hitos del Formulario E-25, sus meses del Art. 17° y la autoridad del Jefe de Proyecto. \\
Responsable & Claudio Cabello (JP). \\
Hitos asociados & M1. \\
Esfuerzo estimado & 24 HH del JP. PERT baja. \\
Referencias & Art. 70, RT-19.02 y Subdocumento 6, 6.1. \\
1.1.1.2 & Plan para la dirección del proyecto \\
Descripción del trabajo & Consolidar los planes, fijar las líneas base de alcance, cronograma y esfuerzo, y someter el plan al Comité Ejecutivo. \\
Entregable & Plan para la dirección del proyecto con nueve planes de apoyo y tres líneas base. \\
Criterio de aceptación & Plan aprobado por el Comité Ejecutivo en el mes 2, con los nueve planes y las tres líneas base con versión y fecha. \\
Responsable & Claudio Cabello (JP). \\
Hitos asociados & M1–M2. Aporta a H1 (mes 2). \\
Esfuerzo estimado & 80 HH (JP 40 y ACP ext. 40). PERT media. \\
Referencias & Subdocumento 6, 6.1.1 y RT-19.01. \\
1.1.1.3 & Espacio colaborativo del proyecto habilitado \\
Descripción del trabajo & Crear el espacio en Microsoft Teams y SharePoint del cliente, con carpetas, permisos por rol y primeros documentos. \\
Entregable & Espacio colaborativo donde el cliente consulta documentos y su estado. \\
Criterio de aceptación & La Contraparte Técnica accede con su cuenta y ve al día acta, plan, entregables, registros de riesgos y cambios. Ningún dato personal queda fuera de carpetas restringidas. \\
Responsable & Claudio Cabello (JP). \\
Hitos asociados & M1. \\
Esfuerzo estimado & 16 HH (JP 8 y ESI 8). PERT baja. \\
Referencias & RT-19.05. \\
1.1.1.4 & Dirección y coordinación integrada del proyecto \\
Descripción del trabajo & Dirigir los frentes, resolver lo que supera a cada líder y representar a Overfore ante el cliente. \\
Entregable & Plan para la dirección ejecutado y actualizado cada semana. \\
Criterio de aceptación & Jefe de Proyecto con dedicación exclusiva del mes 1 al acta de H12 (mes 21) y facultades para comprometer a Overfore. Ningún compromiso vencido sin plan de recuperación registrado. \\
Responsable & Claudio Cabello (JP). \\
Hitos asociados & M1–M21, hasta el acta de H12. \\
Esfuerzo estimado & 1.344 HH del JP, saldo de su dedicación en los meses 1 a 21. Sin PERT. \\
Referencias & Art. 70.3 y RT-19.02. T-12: RT-19.01 y 19.02. RT-21.02. \\
\end{tabla}
\begin{tabla}{I{0.269}I{0.731}}{Diccionario del entregable 1.1.2, Cronograma e informes de valor ganado.}{tab:diccionario-entregable-1-1-2-cronograma}
  {\cab{1.1.2} & \cab{Cronograma e informes de valor ganado}}
  {Fuente: elaboración propia a partir de la Guía del PMBOK (PMI, 2017, p. 162).}
Descripción del trabajo & Cronograma base con los hitos del contrato y control mensual del avance medido en valor ganado. Supone que el mes 1 es enero de 2027. \\
Responsable & Claudio Cabello (JP). \\
Hitos asociados & M1–M20. Aporta a H1 (mes 2). \\
Esfuerzo estimado & 584 HH, 120 externas (JP 344, AS 60, LF 60 y ACP ext. 120). \\
1.1.2.1 & Documento de hitos H1 a H12 \\
Descripción del trabajo & Vincular cada hito con los paquetes que lo activan y aprobar el documento con la Contraparte Técnica. \\
Entregable & Documento con los doce hitos, sus meses y su evidencia de aceptación. \\
Criterio de aceptación & Los doce hitos con su mes y al menos un paquete que los gatilla. Sin diferencias de mes con el Artículo 17, el E-25, el T-17 y la carta Gantt. \\
Responsable & Claudio Cabello (JP). \\
Hitos asociados & M2. Aporta a H1 (mes 2). \\
Esfuerzo estimado & 24 HH del JP. PERT baja. \\
Referencias & Art. 17, Art. 18.1, RT-20.01, Formulario E-25 y Formulario T-18. \\
1.1.2.2 & Informes mensuales de avance con valor ganado \\
Descripción del trabajo & Medir el avance por trabajo terminado, no por horas, analizar desviaciones y revisarlas con la Contraparte Técnica. \\
Entregable & Informe mensual de avance con valor ganado y desviaciones. \\
Criterio de aceptación & Informe entregado en los cinco primeros días hábiles, con índices por paquete. Toda desviación mayor al 10 \% tiene plan de acción. Un índice bajo 0,90 exige plan de recuperación. \\
Responsable & Claudio Cabello (JP). \\
Hitos asociados & M1–M20. \\
Esfuerzo estimado & 320 HH (JP 200 y ACP ext. 120). PERT baja. \\
Referencias & RT-19.06, 19.07, 19.10; Art. 71; Subdocumento 6, 6.1. \\
1.1.2.3 & Registro de cambios con análisis de impacto \\
Descripción del trabajo & Registrar cada solicitud de cambio, analizar su impacto en alcance, plazo y costo, y ejecutar sólo lo aprobado. \\
Entregable & Registro de cambios con impacto y decisión del Comité Ejecutivo. \\
Criterio de aceptación & Ningún cambio se ejecuta sin aprobación previa del Comité Ejecutivo. Cada análisis se entrega en cinco días hábiles y el registro se publica en el espacio compartido. \\
Responsable & Claudio Cabello (JP). \\
Hitos asociados & M2–M20. \\
Esfuerzo estimado & 240 HH (JP 120, AS 60 y LF 60). PERT baja. \\
Referencias & Art. 72 y RT-19.03. \\
\end{tabla}
\begin{tabla}{I{0.269}I{0.731}}{Diccionario del entregable 1.1.3, Documentación de riesgos.}{tab:diccionario-entregable-1-1-3-documentacion}
  {\cab{1.1.3} & \cab{Documentación de riesgos}}
  {Fuente: elaboración propia a partir de la Guía del PMBOK (PMI, 2017, p. 162).}
Descripción del trabajo & Riesgos del proyecto con su análisis y su respuesta, actualizados durante toda la implementación. Supone que el cliente comparte la información de fabricantes, sindicato y personal clave. \\
Responsable & Claudio Cabello (JP). \\
Hitos asociados & M1–M20. Sin hito propio. \\
Esfuerzo estimado & 312 HH, todas del JP. \\
1.1.3.1 & Registro de riesgos con su análisis \\
Descripción del trabajo & Identificar los riesgos en talleres, analizarlos y actualizar el registro en cada Comité de Proyecto. \\
Entregable & Registro de riesgos con causa, probabilidad, impacto, dueño y estado. \\
Criterio de aceptación & Registro con al menos estos riesgos: telemetría, cobertura de red, sindicato, mantenimiento, personal clave e inventario técnico, cada uno con dueño. Cada revisión quincenal queda en acta. \\
Responsable & Claudio Cabello (JP). \\
Hitos asociados & M1–M20. \\
Esfuerzo estimado & 160 HH del JP. PERT baja. \\
Referencias & RT-19.04. \\
1.1.3.2 & Plan de respuesta a riesgos y reservas \\
Descripción del trabajo & Definir la respuesta a cada riesgo alto, calcular las reservas de tiempo y horas, y aprobar el plan. \\
Entregable & Plan de respuesta a riesgos con reservas de tiempo, horas y gestión. \\
Criterio de aceptación & El Comité Ejecutivo aprueba la respuesta de cada riesgo alto o extremo, con responsable, plazo, disparador y plan de contingencia. Cada uso de reserva queda registrado. \\
Responsable & Claudio Cabello (JP). \\
Hitos asociados & M2–M20. \\
Esfuerzo estimado & 152 HH del JP. PERT baja. \\
Referencias & Subdocumento 8, Art. 50.2 y Art. 71. T-12: RT-19.04. \\
\end{tabla}
\begin{tabla}{I{0.269}I{0.731}}{Diccionario del entregable 1.1.4, Plan de calidad y puertas de calidad.}{tab:diccionario-entregable-1-1-4-plan}
  {\cab{1.1.4} & \cab{Plan de calidad y puertas de calidad}}
  {Fuente: elaboración propia a partir de la Guía del PMBOK (PMI, 2017, p. 162).}
Descripción del trabajo & Plan de calidad, cinco puertas que detienen los cambios que no cumplen y evaluaciones de madurez. Supone que la Contraparte Técnica aprueba la definición de terminado en H1. \\
Responsable & Kavon Kermani (LCP). \\
Hitos asociados & M1–M21. Aporta a H1 (mes 2), H3 (mes 6), H5 (mes 12) y H12 (mes 21). \\
Esfuerzo estimado & 224 HH, todas propias (JP 8, LD 16, LCP 152 y QA 48). \\
1.1.4.1 & Plan de calidad del proyecto \\
Descripción del trabajo & Fijar las características de calidad y sus umbrales, definir las métricas y aprobar el plan. \\
Entregable & Plan de calidad del proyecto con sus métricas y su control. \\
Criterio de aceptación & Plan aprobado con un umbral medible por cada característica de calidad que aplica. Umbrales iguales a los del Subdocumento 9 y del Formulario T-13. \\
Responsable & Kavon Kermani (LCP). \\
Hitos asociados & M1–M2. Aporta a H1 (mes 2). \\
Esfuerzo estimado & 56 HH (LCP 48 y QA 8). PERT media. \\
Referencias & Subdocumento 9, 9.1. \\
1.1.4.2 & Puertas de calidad y definición de terminado \\
Descripción del trabajo & Configurar qué bloquea cada puerta, documentar la definición de terminado y comprobar que se detiene un cambio que no cumple. \\
Entregable & Cinco puertas de calidad con umbrales y definición de terminado. \\
Criterio de aceptación & La Contraparte Técnica acepta la definición de terminado. Cobertura de pruebas de 70 \% en lógica de negocio y 80 \% en código nuevo. Complejidad máxima de 10 por método, con excepción hasta 15 con aprobación del par; cero hallazgos críticos o altos pendientes. \\
Responsable & Kavon Kermani (LCP). \\
Hitos asociados & M2–M4. Aporta a H3 (mes 6). \\
Esfuerzo estimado & 64 HH (LD 16, LCP 32 y QA 16). PERT media. \\
Referencias & Subdocumento 9, 9.2.1 y RT-20.07. T-12: RT-04.05 y 04.11; RT-21.20, RNF-TRV-04.11. \\
1.1.4.3 & Protocolo de aceptación de los hitos \\
Descripción del trabajo & Ajustar el protocolo con la Contraparte Técnica, acordar los plazos de revisión y firmar el modelo de acta. \\
Entregable & Protocolo de aceptación de los doce hitos, con modelo de acta. \\
Criterio de aceptación & Protocolo aceptado para los doce hitos, con plazo de revisión de diez días hábiles. Cada criterio se verifica sin interpretación y tiene su evidencia. \\
Responsable & Kavon Kermani (LCP). \\
Hitos asociados & M2. Aporta a H1 (mes 2). \\
Esfuerzo estimado & 32 HH (JP 8 y LCP 24). PERT baja. \\
Referencias & RT-20.08 y T-17. \\
1.1.4.4 & Evaluaciones de madurez OWASP SAMM \\
Descripción del trabajo & Evaluar las quince prácticas del modelo de madurez de seguridad, con el mismo método, y registrar las brechas. \\
Entregable & Informes de madurez inicial, nivel 2 y nivel 3. \\
Criterio de aceptación & Nivel 2 en las quince prácticas en el mes 12, con brechas cerradas. Nivel 3 en las seis de implementación y verificación en el mes 21. \\
Responsable & Kavon Kermani (LCP). \\
Hitos asociados & M2, M12 y M21. Aporta a H1 (mes 2), H5 (mes 12) y H12 (mes 21). \\
Esfuerzo estimado & 72 HH (LCP 48 y QA 24). PERT baja. \\
Referencias & Subdocumento 9, 9.1.3, y Formulario T-17. T-12: RT-11.28. \\
\end{tabla}
\begin{tabla}{I{0.269}I{0.731}}{Diccionario del entregable 1.1.5, Registro de interesados y plan de comunicaciones.}{tab:diccionario-entregable-1-1-5-registro}
  {\cab{1.1.5} & \cab{Registro de interesados y plan de comunicaciones}}
  {Fuente: elaboración propia a partir de la Guía del PMBOK (PMI, 2017, p. 162).}
Descripción del trabajo & Interesados del proyecto, instancias de gobierno, comunicaciones y el acuerdo con el sindicato que condiciona la verificación facial. Supone que el sindicato acepta negociar desde el mes 2. \\
Responsable & Claudio Cabello (JP). \\
Hitos asociados & M1–M20. Aporta a H9 (mes 17). \\
Esfuerzo estimado & 1.312 HH, 264 externas (JP 640, AS 80, ESI 80, LDA 80, LF 56, LI 80, SRE 32, ACP ext. 160 y ABG ext. 104). \\
1.1.5.1 & Actas de comité, coordinación semanal y tablero de estado \\
Descripción del trabajo & Convocar los cuatro comités del Art. 71 (Ejecutivo, de Proyecto, de Arquitectura y de Operación), conducir la coordinación semanal y mantener actas y tablero de estado. \\
Entregable & Actas de cuatro comités, coordinación semanal y tablero de estado. \\
Criterio de aceptación & Cada acta, en dos días hábiles desde la sesión, registra acuerdos, responsables y plazos. El tablero muestra el estado de cada paquete. \\
Responsable & Claudio Cabello (JP). \\
Hitos asociados & M1–M20. \\
Esfuerzo estimado & 992 HH (JP 480, AS 80, ESI 80, LDA 80, LI 80, SRE 32 y ACP externo 160). Sin PERT. \\
Supuestos & Desde el mes 13, el Líder de Operación conduce el Comité de Operación. \\
Referencias & Art. 71 y RT-19.08, 19.09. T-12: RNF-DIS-004. \\
1.1.5.2 & Acuerdo con el sindicato \\
Descripción del trabajo & Presentar la propuesta al sindicato, negociar sus condiciones, incorporarlas al diseño y firmar el acuerdo. \\
Entregable & Acuerdo firmado con el sindicato sobre verificación facial en portería. \\
Criterio de aceptación & Acuerdo firmado por el sindicato y el cliente antes del mes 15, con condiciones de consentimiento y de conservación de datos conformes con la Ley N° 21.719. \\
Responsable & Claudio Cabello (JP). \\
Hitos asociados & M2–M14. Aporta a H9 (mes 17). \\
Esfuerzo estimado & 208 HH (JP 104 y ABG ext. 104). PERT alta. \\
Supuestos & Sin acuerdo al inicio del mes 15, la portería de la Etapa 2 usa sólo credencial de proximidad. \\
Referencias & Subdocumento 6, 6.1.4 y Bases Técnicas, cap. 8 (p. 13) y cap. 19 (p. 33). T-12: RF-ACR-009. \\
1.1.5.3 & Análisis de interesados y calendario de comunicaciones \\
Descripción del trabajo & Identificar y analizar los grupos de interés de Aranda y definir sus mensajes, canales y calendario. \\
Entregable & Registro de interesados y plan de comunicaciones con calendario. \\
Criterio de aceptación & Registro y plan aprobados por el Comité Ejecutivo en el mes 1, con todos los grupos de interés (incluido el turno de noche) y una actualización trimestral en acta. \\
Responsable & Claudio Cabello (JP). \\
Hitos asociados & M1–M20. \\
Esfuerzo estimado & 112 HH (JP 56 y LF 56). PERT baja. \\
Referencias & Art. 89 y Subdocumento 2. \\
\end{tabla}
\begin{tabla}{I{0.269}I{0.731}}{Diccionario del entregable 1.1.6, Plan de adquisiciones y contratos.}{tab:diccionario-entregable-1-1-6-plan}
  {\cab{1.1.6} & \cab{Plan de adquisiciones y contratos}}
  {Fuente: elaboración propia a partir de la Guía del PMBOK (PMI, 2017, p. 162).}
Descripción del trabajo & Compra y contratación de bienes y servicios del proyecto: nube, sala técnica, equipamiento de faena, enlaces, pruebas de intrusión y licencias. Supone proveedores con soporte en Antofagasta. \\
Responsable & Claudio Cabello (JP). \\
Hitos asociados & M1–M14 y M21. Aporta a H3 (mes 6) y H12 (mes 21). \\
Esfuerzo estimado & 320 HH, 120 externas (JP 200, ACP ext. 16 y ABG ext. 104). \\
1.1.6.1 & Plan de adquisiciones \\
Descripción del trabajo & Listar los bienes y servicios, definir el tipo de contrato y fijar las fechas de orden de compra. \\
Entregable & Plan de compras, con tipo de contrato y mes de necesidad. \\
Criterio de aceptación & Plan aprobado por el Comité Ejecutivo, con cotizaciones de la sala técnica y del nodo de faena emitidas en el mes 1. \\
Responsable & Claudio Cabello (JP). \\
Hitos asociados & M1–M2. Aporta a H3 (mes 6). \\
Esfuerzo estimado & 40 HH (JP 24 y ACP ext. 16). PERT media. \\
Referencias & Subdocumento 6, 6.1.6. \\
1.1.6.2 & Contratos con proveedores y subcontratistas \\
Descripción del trabajo & Especificar, adjudicar y administrar los contratos. Excluye la obra civil y los dispositivos de terreno, a cargo del cliente. \\
Entregable & Siete contratos firmados con proveedores y subcontratistas. \\
Criterio de aceptación & Cada contrato se firma antes de la fecha en que se requiere el bien o servicio. Cada recepción se inspecciona contra su especificación y queda en acta. \\
Responsable & Claudio Cabello (JP). \\
Hitos asociados & M2–M14. \\
Esfuerzo estimado & 224 HH (JP 128 y ABG externo 96). PERT alta. \\
Referencias & Subdocumento 6, 6.1.6 Art. 73; Subdocumento 4, 4.2.3.2, Subdocumento 1, Anexo F del Subdocumento 1, RT-03.17 y RT-06.33. T-12: RT-06.06 y 06.33. \\
1.1.6.3 & Registro de licencias y garantías \\
Descripción del trabajo & Registrar cada licencia y garantía, asignar el titular y programar las renovaciones. \\
Entregable & Registro de licencias y garantías, con vencimientos y plan de reposición. \\
Criterio de aceptación & Registro con titular, vencimiento y condiciones de cada licencia y garantía, y el soporte de cinco años del equipamiento. Ninguna licencia vence sin renovación programada. \\
Responsable & Claudio Cabello (JP). \\
Hitos asociados & M2–M14. \\
Esfuerzo estimado & 24 HH del JP. PERT baja. \\
Referencias & Subdocumento 4, 4.2.1.4, Subdocumento 6, 6.1.6, y RT-08.06, 08.13. \\
1.1.6.4 & Garantías y seguros del contrato \\
Descripción del trabajo & Constituir y mantener las garantías de fiel cumplimiento y de correcto funcionamiento, y las pólizas de seguro del contrato. \\
Entregable & Garantías de fiel cumplimiento y correcto funcionamiento, y pólizas de seguro. \\
Criterio de aceptación & Cada instrumento con su monto y vigencia, recibido conforme por el cliente. La garantía de fiel cumplimiento se entrega en diez días hábiles desde la notificación de la adjudicación (Art. 36.1) y la de correcto funcionamiento es condición de H12 (Art. 37.1). \\
Responsable & Claudio Cabello (JP). \\
Hitos asociados & M1, M12 y M21. Aporta a H12 (mes 21). \\
Esfuerzo estimado & 32 HH (JP 24 y ABG externo 8). PERT baja. \\
Referencias & Art. 36, 37, 38, 67 y 68 de las Bases administrativas. \\
\end{tabla}
\begin{tabla}{I{0.269}I{0.731}}{Diccionario del entregable 1.1.7, Equipo acreditado y logística en faena.}{tab:diccionario-entregable-1-1-7-equipo}
  {\cab{1.1.7} & \cab{Equipo acreditado y logística en faena}}
  {Fuente: elaboración propia a partir de la Guía del PMBOK (PMI, 2017, p. 162).}
Descripción del trabajo & Habilitar al equipo del proyecto para trabajar en la faena, a 3.180 metros sobre el mar y 145 km de Antofagasta. Supone que el cliente resuelve las acreditaciones en el plazo de su procedimiento. \\
Responsable & Claudio Cabello (JP). \\
Hitos asociados & M1–M20. Sin hito propio. \\
Esfuerzo estimado & 496 HH, 272 externas (JP 224 y ACP ext. 272). \\
1.1.7.1 & Equipo del proyecto acreditado en faena \\
Descripción del trabajo & Programar exámenes médicos, la inducción de seguridad minera, controlar credenciales y designar suplentes. \\
Entregable & Equipo con exámenes médicos, inducción de seguridad y credencial de acceso. \\
Criterio de aceptación & Cada persona tiene credencial vigente antes de su primer turno en terreno. Todo reemplazo de personal clave tiene autorización escrita del cliente y quince días hábiles de traslape. \\
Responsable & Claudio Cabello (JP). \\
Hitos asociados & M1–M18. \\
Esfuerzo estimado & 336 HH (JP 64 y ACP externo 272). PERT baja. \\
Supuestos & Un contratista nuevo tarda unos nueve días en acreditarse, los trámites del equipo clave empiezan en la primera semana. \\
Referencias & Bases Técnicas, 2.4 (p. 5); RT-15.07 a 15.09; Art. 76.1 y 76.2, Tabla \ref{tab:roles-equipo-clave-dedicacion-paquetes-cargo}. T-12: RP-001 y 008. \\
1.1.7.2 & Plan de logística y turnos en faena \\
Descripción del trabajo & Definir los turnos del equipo, reservar traslados y alojamiento y actualizar el plan cada mes. \\
Entregable & Plan de traslados, alojamiento y turnos en faena. \\
Criterio de aceptación & Presencia en faena que exige cada paquete, sin turnos que excedan la jornada excepcional autorizada, y al menos una persona del proyecto por turno, de día y de noche, en cada área en marcha blanca. \\
Responsable & Claudio Cabello (JP). \\
Hitos asociados & M1–M20. \\
Esfuerzo estimado & 160 HH del JP. PERT baja. \\
Supuestos & Los turnos de doce horas en 7x7 promedian 42 horas semanales hasta abril de 2028 y 40 desde entonces (Ley N° 21.561), con un turno libre adicional cada seis semanas. \\
Referencias & Bases Técnicas, RT-21.16, p. 26; cap. 19, p. 33. \\
\end{tabla}
\begin{tabla}{I{0.269}I{0.731}}{Diccionario del entregable 1.1.8, Actas de cierre de etapa y lecciones aprendidas.}{tab:diccionario-entregable-1-1-8-actas}
  {\cab{1.1.8} & \cab{Actas de cierre de etapa y lecciones aprendidas}}
  {Fuente: elaboración propia a partir de la Guía del PMBOK (PMI, 2017, p. 162).}
Descripción del trabajo & Cierre administrativo de cada etapa y registro de lecciones aprendidas. Supone que los hitos H7 y H12 quedan aceptados antes del cierre de su etapa. \\
Responsable & Claudio Cabello (JP). \\
Hitos asociados & M16–M21. Aporta a H7 (mes 16) y H12 (mes 21). \\
Esfuerzo estimado & 80 HH, todas propias (JP 56 y GS 24). \\
1.1.8.1 & Acta de cierre de la Etapa 1 \\
Descripción del trabajo & Verificar la aceptación de H1 a H7, cerrar los contratos de la etapa y firmar el acta. \\
Entregable & Acta que cierra la Etapa 1 con los pendientes traspasados. \\
Criterio de aceptación & Acta firmada por la Contraparte Técnica con los siete hitos de la Etapa 1 aceptados. Sin observaciones abiertas de nivel crítico o alto. \\
Responsable & Claudio Cabello (JP). \\
Hitos asociados & M16, una vez aceptado H7. \\
Esfuerzo estimado & 24 HH del JP. PERT baja. \\
Referencias & Art. 18. \\
1.1.8.2 & Acta de cierre de la Etapa 2 \\
Descripción del trabajo & Verificar la aceptación de H8 a H12, cerrar los contratos de la etapa y firmar el acta. \\
Entregable & Acta que cierra la Etapa 2 y la implementación. \\
Criterio de aceptación & Acta firmada por la Contraparte Técnica con los cinco hitos de la Etapa 2 aceptados. Sin observaciones abiertas de nivel crítico o alto. \\
Responsable & Claudio Cabello (GS). \\
Hitos asociados & M21, una vez aceptado H12. \\
Esfuerzo estimado & 24 HH del GS. PERT baja. \\
Referencias & Art. 18. \\
1.1.8.3 & Registro de lecciones aprendidas \\
Descripción del trabajo & Realizar la sesión de lecciones con cada frente, consolidar el registro y publicarlo. \\
Entregable & Registro de aciertos y fallas de cada etapa, con acciones. \\
Criterio de aceptación & Registro con al menos una lección por frente y su acción para la etapa siguiente, revisado con la Contraparte Técnica al cierre de cada etapa. \\
Responsable & Claudio Cabello (JP). \\
Hitos asociados & M16 y M21. \\
Esfuerzo estimado & 32 HH del JP. PERT baja. \\
Referencias & PMI (2017), sección 4.7. \\
\end{tabla}
\begin{tabla}{I{0.269}I{0.731}}{Diccionario del entregable 1.1.9, Plan de reversibilidad y salida.}{tab:diccionario-entregable-1-1-9-plan}
  {\cab{1.1.9} & \cab{Plan de reversibilidad y salida}}
  {Fuente: elaboración propia a partir de la Guía del PMBOK (PMI, 2017, p. 162).}
Descripción del trabajo & Plan para que el cliente cambie de operador o internalice el servicio sin perder datos ni continuidad. Supone que el plan inicial se entrega dentro de los primeros 90 días y se actualiza cada año. \\
Responsable & Claudio Cabello (JP hasta el acta de H12, GS desde la primera ola de la Etapa 2). \\
Hitos asociados & M3, M15, M27, M39, M51. Sin hito propio. \\
Esfuerzo estimado & 112 HH, todas propias (JP 32, AS 8, LDA 8 y SRE 64). \\
1.1.9.1 & Plan de reversibilidad inicial \\
Descripción del trabajo & Inventariar código, datos, configuraciones y credenciales, declarar qué se puede trasladar y aprobar el plan. \\
Entregable & Plan de salida con inventario, formatos, plazos y acompañamiento. \\
Criterio de aceptación & Plan aceptado por la Contraparte Técnica dentro de los primeros 90 días. Incluye inventario de activos, traspaso de credenciales y exportación íntegra de datos en formatos abiertos. Acompañamiento de al menos 90 días (Art. 77.2). \\
Responsable & Claudio Cabello (JP). \\
Hitos asociados & M3. \\
Esfuerzo estimado & 48 HH (JP 32, AS 8 y LDA 8). PERT media. \\
Referencias & Art. 77.2. T-12: RT-03.07; RT-05.06. \\
1.1.9.2 & Actualizaciones anuales del plan de reversibilidad \\
Descripción del trabajo & Revisar el inventario, actualizar el plan y aprobarlo con la Contraparte Técnica. \\
Entregable & Revisión anual del plan de reversibilidad con los cambios del año. \\
Criterio de aceptación & Una actualización aprobada por año, en los meses 15, 27, 39 y 51. El inventario incluye todo componente puesto en producción durante el año. \\
Responsable & Lucas Zamora (SRE). \\
Hitos asociados & M15, M27, M39, M51. \\
Esfuerzo estimado & 64 HH del SRE. PERT baja. \\
Referencias & Art. 77.2. \\
\end{tabla}
\begin{tabla}{I{0.269}I{0.731}}{Diccionario del entregable 1.1.10, Documentación contractual del proyecto.}{tab:diccionario-entregable-1-1-10-documentacion}
  {\cab{1.1.10} & \cab{Documentación contractual del proyecto}}
  {Fuente: elaboración propia a partir de la Guía del PMBOK (PMI, 2017, p. 162).}
Descripción del trabajo & Índice de los documentos contractuales del Artículo 91 y de los documentos de operación. Supone que la lista del Artículo 91 no cambia durante el contrato. \\
Responsable & Claudio Cabello (JP). \\
Hitos asociados & M3–M20. Aporta a H7 (mes 16) y H12 (mes 21). \\
Esfuerzo estimado & 184 HH, todas propias (JP 16, SRE 112 y OPS 56). \\
1.1.10.1 & Índice de documentación contractual \\
Descripción del trabajo & Listar los documentos, asignar el paquete de origen y mantener el índice en cada hito. \\
Entregable & Índice de los documentos del Artículo 91, con su paquete y hito. \\
Criterio de aceptación & Cada documento del Artículo 91 aparece con su paquete y su hito. Índice actualizado en cada acta de conformidad. \\
Responsable & Claudio Cabello (JP). \\
Hitos asociados & M3. \\
Esfuerzo estimado & 16 HH del JP. PERT baja. \\
Referencias & Art. 91. \\
1.1.10.2 & Libros de operación y matriz de escalamiento \\
Descripción del trabajo & Redactar un libro de operación por escenario de falla, definir la escalada de incidentes y probar los libros. \\
Entregable & Manual, cinco libros con guías de resolución y matriz de escalamiento. \\
Criterio de aceptación & El Comité de Operación aprueba el manual, la matriz y un libro para cada una de las cinco fallas de resiliencia, probados antes de aceptarse. \\
Responsable & Lucas Zamora (SRE). \\
Hitos asociados & M14–M20. Aporta a H7 (mes 16) y H12 (mes 21). \\
Esfuerzo estimado & 168 HH (SRE 112 y OPS 56). PERT media. \\
Referencias & Art. 91 y RT-21.04. T-12: RT-14.05; RP-009. \\
\end{tabla}

\apartado{Rama 1.2 Requerimientos del producto}
La rama 1.2 tiene 6 entregables y 19 paquetes de trabajo, que suman 3.684 HH. Las tablas \ref{tab:diccionario-entregable-1-2-1-inventario} a \ref{tab:diccionario-entregable-1-2-6-registro} presentan la ficha de cada entregable y, bajo ella, la de cada uno de sus paquetes.

\begin{tabla}{I{0.269}I{0.731}}{Diccionario del entregable 1.2.1, Inventario técnico de los sistemas existentes.}{tab:diccionario-entregable-1-2-1-inventario}
  {\cab{1.2.1} & \cab{Inventario técnico de los sistemas existentes}}
  {Fuente: elaboración propia a partir de la Guía del PMBOK (PMI, 2017, p. 162).}
Descripción del trabajo & Inventario técnico de los diez sistemas existentes con los que SITCO Aranda convive o se integra, sin modificar ninguno. Supone que cada dueño entrega su información en los tres primeros meses. \\
Responsable & Marcelo Flores (LI). \\
Hitos asociados & M1–M3. Aporta a H2 (mes 4). \\
Esfuerzo estimado & 208 HH, todas propias (AS 8, LDA 24, LI 144 y ADT 32). \\
1.2.1.1 & Fichas técnicas de los sistemas existentes \\
Descripción del trabajo & Entrevistar a cada dueño, revisar la documentación, completar y validar las fichas. \\
Entregable & Ficha de cada sistema existente con versión, interfaces, volumen y dueño. \\
Criterio de aceptación & Diez fichas validadas por el dueño de cada sistema. Cada ficha indica versión, responsable, dependencias y si se lee por interfaz o por base de lectura. \\
Responsable & Marcelo Flores (LI). \\
Hitos asociados & M1–M3. Aporta a H2 (mes 4). \\
Esfuerzo estimado & 160 HH (LDA 24, LI 104 y ADT 32). PERT alta. \\
Referencias & Bases Técnicas, cap. 5 (p. 10), Subdocumento 3, 3.3.2.2, y Subdocumento 2, Anexo A, RP-005. \\
1.2.1.2 & Matriz de integraciones \\
Descripción del trabajo & Listar las conexiones entre sistemas, definir su forma de intercambio y horario, y validarlas con los dueños. \\
Entregable & Matriz de integraciones con modo, dirección, volumen, frecuencia y ventana. \\
Criterio de aceptación & Matriz con toda integración del Subdocumento 4. Cada una declara qué hace la solución si la contraparte no responde. Ninguna escribe en la red de control de proceso. \\
Responsable & Marcelo Flores (LI). \\
Hitos asociados & M2–M3. Aporta a H2 (mes 4). \\
Esfuerzo estimado & 48 HH (AS 8 y LI 40). PERT alta. \\
Referencias & RT-05.21. \\
\end{tabla}
\begin{tabla}{I{0.269}I{0.731}}{Diccionario del entregable 1.2.2, Línea base de alcance de la Etapa 1.}{tab:diccionario-entregable-1-2-2-linea}
  {\cab{1.2.2} & \cab{Línea base de alcance de la Etapa 1}}
  {Fuente: elaboración propia a partir de la Guía del PMBOK (PMI, 2017, p. 162).}
Descripción del trabajo & Alcance de referencia de la Etapa 1: catálogo de requerimientos, matriz de trazabilidad y criterios de aceptación. Supone acceso a los usuarios principales de mina, planta y puerto en los dos primeros meses. \\
Responsable & Patricio Figueroa (LF). \\
Hitos asociados & M1–M12. Aporta a H1 (mes 2), H4 (mes 10) y H5 (mes 12). \\
Esfuerzo estimado & 1.840 HH, todas propias (LD 16, LF 1.616, QA 192 y ADT 16). \\
1.2.2.1 & Catálogo de requerimientos de la Etapa 1 \\
Descripción del trabajo & Realizar talleres por proceso, redactar los requerimientos y priorizarlos con el cliente. \\
Entregable & Catálogo de requerimientos de la Etapa 1 con identificador y origen. \\
Criterio de aceptación & El 100 \% de los requerimientos tiene identificador, origen, resultado esperado y forma de verificarlo. Sin ambigüedades ni duplicados. \\
Responsable & Patricio Figueroa (LF). \\
Hitos asociados & M1–M2. Aporta a H1 (mes 2). \\
Esfuerzo estimado & 240 HH (LF 160 y QA 80). PERT media. \\
Referencias & Subdocumento 3, 3.2.4. \\
1.2.2.2 & Matriz de trazabilidad de la Etapa 1 \\
Descripción del trabajo & Asignar cada requerimiento a un componente y a una prueba, y registrarlo en el Formulario T-12. \\
Entregable & Matriz que relaciona cada requerimiento con su componente y su prueba. \\
Criterio de aceptación & Ningún requerimiento queda sin componente ni prueba, y cada uno llega hasta el despliegue realizado. Se verifica en cada puerta 4. \\
Responsable & Patricio Figueroa (LF). \\
Hitos asociados & M2. Aporta a H1 (mes 2). \\
Esfuerzo estimado & 56 HH (LF 16 y QA 40). PERT media. \\
Referencias & RT-04.04. \\
1.2.2.3 & Criterios de aceptación de la Etapa 1 \\
Descripción del trabajo & Redactar un criterio por requerimiento funcional y acordarlos con la Contraparte Técnica. \\
Entregable & Criterios verificables con que se acepta el alcance de la Etapa 1. \\
Criterio de aceptación & La Contraparte Técnica acuerda los casos de aceptación de los 92 requerimientos funcionales de la Etapa 1, trazados a componente y prueba. \\
Responsable & Patricio Figueroa (LF). \\
Hitos asociados & M2. Aporta a H1 (mes 2). \\
Esfuerzo estimado & 96 HH (LF 24 y QA 72). PERT media. \\
Referencias & Subdocumento 3, 3.2.5. \\
1.2.2.4 & Backlog del producto y revisiones de sprint de la Etapa 1 \\
Descripción del trabajo & Ordenar la lista única de trabajo por etapa, dependencia y valor, y revisarla cada semana con los referentes del cliente. \\
Entregable & Lista de trabajo priorizada y actas de revisión firmadas. \\
Criterio de aceptación & Cada sprint empieza con historias terminadas según la definición de listo y cierra con acta firmada por el referente del área. Toda historia aceptada queda trazada a su requerimiento. \\
Responsable & Patricio Figueroa (LF). \\
Hitos asociados & M1–M12. Aporta a H4 (mes 10) y H5 (mes 12). \\
Esfuerzo estimado & 1.344 HH del LF, saldo de su dedicación en los meses 1 a 12. Sin PERT. \\
Supuestos & Sprints de dos semanas. En los meses 1 y 2 la lista inicial sale del catálogo de requerimientos. \\
Referencias & Subdocumento 6, 6.2 y RT-04.04. \\
1.2.2.5 & Taxonomía única de demoras y detenciones \\
Descripción del trabajo & Hacer talleres con supervisores de mina y de planta y acordar una sola lista de causas de demora. \\
Entregable & Lista única de causas de demoras con código y definición. \\
Criterio de aceptación & Las superintendencias de mina y de planta aprueban la lista antes del mes 6. Es un catálogo cerrado, sin texto libre, con versiones y vigencia. \\
Responsable & Patricio Figueroa (LF). \\
Hitos asociados & M3–M5. Aporta a H4 (mes 10). \\
Esfuerzo estimado & 64 HH (LF 32, LD 16 y ADT 16). PERT media. \\
Supuestos & Los talleres del mes 3 se hacen fuera de los seis días de la detención de marzo de 2027. \\
Referencias & Subdocumento 3, Tabla 3.3, y Bases Técnicas, 7.2 y cap. 18, criterio 5. T-12: RF-PLA-005; RNF-USA-004, RP-011. \\
1.2.2.6 & Requisitos de cadena de custodia y de aseguramiento del cobre \\
Descripción del trabajo & Levantar con compradores la evidencia de cadena de custodia del cobre y traducirla a campos del módulo. \\
Entregable & Matriz de exigencias por esquema y comprador con su evidencia. \\
Criterio de aceptación & Cada exigencia tiene un campo o registro en 1.4.3.3, validado por el área comercial en acta antes del mes 5. \\
Responsable & Patricio Figueroa (LF). \\
Hitos asociados & M3–M4. Aporta a H4 (mes 10). \\
Esfuerzo estimado & 40 HH del LF. PERT media. \\
Referencias & Bases Técnicas, 16.2 (p. 27) y cap. 18, criterio 1. \\
\end{tabla}
\begin{tabla}{I{0.269}I{0.731}}{Diccionario del entregable 1.2.3, Acta de ratificación de las decisiones de diseño.}{tab:diccionario-entregable-1-2-3-acta}
  {\cab{1.2.3} & \cab{Acta de ratificación de las decisiones de diseño}}
  {Fuente: elaboración propia a partir de la Guía del PMBOK (PMI, 2017, p. 162).}
Descripción del trabajo & Acta en que el cliente ratifica por escrito las catorce decisiones de diseño que dejó abiertas. Supone que el dato oficial de tonelaje queda ratificado en el mes 3. \\
Responsable & Patricio Figueroa (LF). \\
Hitos asociados & M1–M3. Aporta a H2 (mes 4). \\
Esfuerzo estimado & 184 HH, todas propias (AS 84, LDA 28, LF 44 y LI 28). \\
1.2.3.1 & Análisis de alternativas de las decisiones pendientes \\
Descripción del trabajo & Describir las alternativas de cada decisión, evaluar su efecto en el proceso, los datos e integraciones y recomendar una. \\
Entregable & Análisis de cada decisión pendiente con su recomendación. \\
Criterio de aceptación & Catorce análisis, uno por decisión de las Bases Técnicas, 16.1. Indican la decisión adoptada, su paquete y su efecto en el operador, los datos y las HH. \\
Responsable & Benjamín Rodríguez (AS). \\
Hitos asociados & M1–M3. Aporta a H2 (mes 4). \\
Esfuerzo estimado & 168 HH (AS 84, LDA 28, LF 28 y LI 28). PERT media. \\
Referencias & Bases Técnicas, 16.1 (p. 26). \\
1.2.3.2 & Acta de ratificación firmada por el cliente \\
Descripción del trabajo & Presentar los análisis al Comité Ejecutivo, registrar la decisión y firmar el acta. \\
Entregable & Acta en que el cliente ratifica las catorce decisiones de diseño. \\
Criterio de aceptación & Acta firmada al término del mes 3, con la fuente del tonelaje ratificada. Las decisiones de terceros tienen fecha de cierre: sindicato antes del mes 15 y fabricantes antes del mes 13. \\
Responsable & Patricio Figueroa (LF). \\
Hitos asociados & M3. Aporta a H2 (mes 4). \\
Esfuerzo estimado & 16 HH del LF. PERT media. \\
Referencias & Bases Técnicas, 16.1 (p. 26) y Subdocumento 3. T-12: RF-CON-002. \\
\end{tabla}
\begin{tabla}{I{0.269}I{0.731}}{Diccionario del entregable 1.2.4, Acuerdos de acceso a la telemetría.}{tab:diccionario-entregable-1-2-4-acuerdos}
  {\cab{1.2.4} & \cab{Acuerdos de acceso a la telemetría}}
  {Fuente: elaboración propia a partir de la Guía del PMBOK (PMI, 2017, p. 162).}
Descripción del trabajo & Acuerdos con los tres fabricantes de equipos para recibir su telemetría. El cliente firma cada acuerdo, los cobros de los fabricantes son de su cargo. \\
Responsable & Marcelo Flores (LI). \\
Hitos asociados & M1–M12. Aporta a H9 (mes 17). \\
Esfuerzo estimado & 192 HH, de las que 24 HH son externas (LI 152, DEV 16 y ABG externo 24). \\
1.2.4.1 & Especificación de acceso a la telemetría \\
Descripción del trabajo & Definir qué datos necesita el sistema de cada equipo y el protocolo de envío. No incluye pagos a fabricantes. \\
Entregable & Especificación de datos, protocolo y frecuencia de cada fabricante. \\
Criterio de aceptación & El cliente aprueba por escrito la especificación. Cada fabricante queda declarado con su forma de envío de datos. \\
Responsable & Marcelo Flores (LI). \\
Hitos asociados & M1–M3. \\
Esfuerzo estimado & 72 HH (LI 56 y DEV 16). PERT alta. \\
Referencias & Subdocumento 4, Tabla \refext{S4}{tab:conectores-integrador-sistemas-etapa-lugar-ejecucion}, y Bases Técnicas, cap. 5 (p. 10). T-12: RT-05.23; RF-EQU-002, RNF-INT-005. \\
1.2.4.2 & Acuerdos firmados con los fabricantes \\
Descripción del trabajo & Acompañar la negociación con los fabricantes, verificar el acceso de prueba y registrar los acuerdos firmados. \\
Entregable & Tres acuerdos firmados por el cliente, con acceso de prueba verificado. \\
Criterio de aceptación & Tres acuerdos firmados antes del inicio del mes 13, con acceso de prueba verificado. Cada acuerdo indica datos, frecuencia, formato y propiedad. \\
Responsable & Marcelo Flores (LI). \\
Hitos asociados & M3–M12. Aporta a H9 (mes 17). \\
Esfuerzo estimado & 120 HH (LI 96 y ABG externo 24). PERT alta. \\
Supuestos & Si un fabricante no firma antes del mes 13, su conector se construye contra la especificación y la conexión productiva espera la firma del cliente. \\
Referencias & Bases Técnicas, cap. 8 (p. 13). T-12: RF-EQU-003. \\
\end{tabla}
\begin{tabla}{I{0.269}I{0.731}}{Diccionario del entregable 1.2.5, Línea base de alcance de la Etapa 2.}{tab:diccionario-entregable-1-2-5-linea}
  {\cab{1.2.5} & \cab{Línea base de alcance de la Etapa 2}}
  {Fuente: elaboración propia a partir de la Guía del PMBOK (PMI, 2017, p. 162).}
Descripción del trabajo & Base de referencia del alcance de la Etapa 2: lista de requerimientos, cómo se verifica cada uno y cómo se acepta. Supone la Etapa 1 en marcha blanca. \\
Responsable & Patricio Figueroa (LF). \\
Hitos asociados & M13–M21. Aporta a H8 (mes 14), H9 (mes 17) y H10 (mes 18). \\
Esfuerzo estimado & 1.196 HH, todas propias (LDA 16, LF 1.100, QA 56 y ADT 24). \\
1.2.5.1 & Catálogo de requerimientos de la Etapa 2 \\
Descripción del trabajo & Realizar talleres con mantenimiento, seguridad y medio ambiente, y redactar y priorizar los requerimientos. \\
Entregable & Catálogo con los requerimientos de la Etapa 2. \\
Criterio de aceptación & Los 34 requerimientos funcionales de la Etapa 2 tienen identificador, origen y método de verificación, según ISO/IEC/IEEE 29148. \\
Responsable & Patricio Figueroa (LF). \\
Hitos asociados & M13. Aporta a H8 (mes 14). \\
Esfuerzo estimado & 96 HH (LF 72 y QA 24). PERT media. \\
Referencias & Subdocumento 3, 3.2.2.3. \\
1.2.5.2 & Matriz de trazabilidad de la Etapa 2 \\
Descripción del trabajo & Asignar a cada requerimiento un componente y una prueba, y actualizar el Formulario T-12. \\
Entregable & Tabla que vincula cada requerimiento con su componente y su prueba. \\
Criterio de aceptación & La tabla incluye los 34 requerimientos funcionales y los no funcionales que les aplican, y las pruebas de la Etapa 1 que se repiten en cada entrega. \\
Responsable & Patricio Figueroa (LF). \\
Hitos asociados & M13–M14. Aporta a H8 (mes 14). \\
Esfuerzo estimado & 32 HH (LF 16 y QA 16). PERT media. \\
Referencias & RT-04.04. \\
1.2.5.3 & Criterios de aceptación de la Etapa 2 \\
Descripción del trabajo & Redactar un caso de aceptación por requerimiento y acordarlo con la Contraparte Técnica. \\
Entregable & Casos de aceptación verificables para cada requerimiento de la Etapa 2. \\
Criterio de aceptación & Los 34 requerimientos funcionales de la Etapa 2 tienen caso de aceptación acordado con la Contraparte Técnica. \\
Responsable & Patricio Figueroa (LF). \\
Hitos asociados & M14. Aporta a H8 (mes 14). \\
Esfuerzo estimado & 40 HH (LF 24 y QA 16). PERT media. \\
Referencias & Subdocumento 3, 3.2.5. \\
1.2.5.4 & Backlog del producto y revisiones de sprint de la Etapa 2 \\
Descripción del trabajo & Ordenar el backlog de la Etapa 2 y conducir sus revisiones con mantenimiento, seguridad y medio ambiente. \\
Entregable & Lista priorizada de pendientes y actas firmadas de cada revisión. \\
Criterio de aceptación & Cada período de dos semanas cierra con acta firmada por el referente del área. Ninguna historia se acepta si rompe las pruebas de la Etapa 1. \\
Responsable & Patricio Figueroa (LF). \\
Hitos asociados & M13–M21. Aporta a H9 (mes 17) y H10 (mes 18). \\
Esfuerzo estimado & 972 HH del LF, saldo de su dedicación en los meses 13 a 21. Sin PERT. \\
Supuestos & Desde el mes 19 sólo entran ajustes de la marcha blanca. \\
Referencias & Subdocumento 6, 6.2 y RT-04.04. \\
1.2.5.5 & Revisión de un lote piloto con el comprador \\
Descripción del trabajo & Reconstruir en SITCO el origen de un lote embarcado y revisarlo con el comprador. Registrar sus observaciones. \\
Entregable & Informe del lote piloto con origen, evidencia y observaciones del comprador. \\
Criterio de aceptación & El origen del lote se reconstruye en menos de 2 horas, con su evidencia completa. Revisión con el comprador registrada en acta. Cada observación pasa a la lista de pendientes. \\
Responsable & Patricio Figueroa (LF). \\
Hitos asociados & M17. \\
Esfuerzo estimado & 56 HH (LF 16, LDA 16 y ADT 24). PERT alta. \\
Supuestos & El cliente elige el lote y convoca al comprador. \\
Referencias & Subdocumento 3, Tabla \refext{S3}{tab:correspondencia-entre-problema-solucion}, Subdocumento 6, 6.1.4 y Bases Técnicas, 13.2 (p. 22). \\
\end{tabla}
\begin{tabla}{I{0.269}I{0.731}}{Diccionario del entregable 1.2.6, Registro documentado de los factores de ajuste.}{tab:diccionario-entregable-1-2-6-registro}
  {\cab{1.2.6} & \cab{Registro documentado de los factores de ajuste}}
  {Fuente: elaboración propia a partir de la Guía del PMBOK (PMI, 2017, p. 162).}
Descripción del trabajo & Explicación escrita de los factores de ajuste del balance metalúrgico, registrada con la persona que hoy los conoce en la compañía. Si es posible, también con la persona jubilada. \\
Responsable & Benjamín Velásquez (LDA). \\
Hitos asociados & M2. Sin hito propio. \\
Esfuerzo estimado & 64 HH, todas propias (LDA 48 y ADT 16). \\
1.2.6.1 & Explicación escrita de los factores de ajuste \\
Descripción del trabajo & Entrevistar a los metalurgistas, reconstruir cada factor y validar el documento. \\
Entregable & Documento con origen, fórmula y vigencia de cada factor desde 2013. \\
Criterio de aceptación & Metalurgia valida cada factor vigente. Doce cierres mensuales anteriores, recalculados con esos factores, se contrastan con el cierre original (TRA-002). \\
Responsable & Benjamín Velásquez (LDA). \\
Hitos asociados & M2. \\
Esfuerzo estimado & 40 HH (LDA 32 y ADT 8). PERT alta. \\
Referencias & Bases Técnicas, cap. 8 (p. 14) y Subdocumento 5, 5.3.2. T-12: RF-CON-007; TRA-002. \\
1.2.6.2 & Registro en video validado por metalurgia \\
Descripción del trabajo & Grabar las sesiones, indexar el video por factor y obtener la validación. \\
Entregable & Video de la explicación de los factores, validado por metalurgia. \\
Criterio de aceptación & Video ordenado por factor y validado por metalurgia al cierre del mes 2. Queda en el espacio del cliente con acceso restringido. \\
Responsable & Benjamín Velásquez (LDA). \\
Hitos asociados & M2. \\
Esfuerzo estimado & 24 HH (LDA 16 y ADT 8). PERT alta. \\
Referencias & Subdocumento 5, 5.3.2. \\
\end{tabla}

\apartado{Rama 1.3 Diseño del sistema}
La rama 1.3 tiene 6 entregables y 19 paquetes de trabajo, que suman 1.800 HH. Las tablas \ref{tab:diccionario-entregable-1-3-1-arquitectura} a \ref{tab:diccionario-entregable-1-3-6-matriz} presentan la ficha de cada entregable y, bajo ella, la de cada uno de sus paquetes.

\begin{tabla}{I{0.269}I{0.731}}{Diccionario del entregable 1.3.1, Arquitectura de la solución.}{tab:diccionario-entregable-1-3-1-arquitectura}
  {\cab{1.3.1} & \cab{Arquitectura de la solución}}
  {Fuente: elaboración propia a partir de la Guía del PMBOK (PMI, 2017, p. 162).}
Descripción del trabajo & Documento de arquitectura de la solución híbrida, con vistas, decisiones, capacidad, red y continuidad. Supone las decisiones de diseño ratificadas en el mes 3. \\
Responsable & Benjamín Rodríguez (AS). \\
Hitos asociados & M2–M20, arquitectura en M4. Aporta a H2 (mes 4) y H7 (mes 16). \\
Esfuerzo estimado & 464 HH, 144 externas (AS 272, ESI 16, LDA 8, LD 24, RFE ext. 112 y BCP ext. 32). \\
1.3.1.1 & Arquitectura lógica y registro de decisiones \\
Descripción del trabajo & Modelar las capas y módulos del sistema, registrar cada decisión con sus alternativas y aprobarla en el Comité de Arquitectura. \\
Entregable & Vista de capas y componentes, y registro de decisiones con alternativas. \\
Criterio de aceptación & Registro de decisiones fechadas, aprobado por el Comité de Arquitectura. La vista muestra las ocho capas del numeral 2.1. \\
Responsable & Benjamín Rodríguez (AS). \\
Hitos asociados & M2–M4. Aporta a H2 (mes 4). \\
Esfuerzo estimado & 80 HH (AS 56 y LD 24). PERT media. \\
Referencias & RT-02.01 a 02.04. T-12: RT-02.01 a 02.04 y 02.14; RT-05.02, RT-26.01. \\
1.3.1.2 & Arquitectura física y plan de capacidad \\
Descripción del trabajo & Dimensionar la capacidad desde la volumetría, diseñar los respaldos y clasificar los servicios. \\
Entregable & Plano de equipos en nube y faena, con respaldos y prioridad. \\
Criterio de aceptación & La Contraparte Técnica acepta el diseño con prueba de 300 transacciones por segundo sin cambiar la topología. El relevo de turno suma 135: 101 de personas y 34 de periféricos. \\
Responsable & Benjamín Rodríguez (AS). \\
Hitos asociados & M2–M4. Aporta a H2 (mes 4). \\
Esfuerzo estimado & 80 HH (AS 72 y LDA 8). PERT media. \\
Referencias & Subdocumento 4, 4.2.3.5, 4.2.4, RT-07.09 a 07.14, cap. 9 de las Bases técnicas transversales; Bases Técnicas, 14.1; RT-09.02, RT-10.01, 10.02 y Subdocumento 5, 5.4.1. T-12: RT-02.03, 02.10 y 02.11; RT-03.05, 03.08 y 03.09, RT-06.34, RT-08.01 y 08.05; RT-09.01, 09.03 y 09.05; RT-15.01, 15.04 y 15.05; RNF-CAP-001 a 005 y 008; RNF-OPE-003, RP-003. \\
1.3.1.3 & Diseño de red, enlaces y zonas de seguridad \\
Descripción del trabajo & Diseñar la topología, definir las zonas y sus flujos permitidos y especificar los enlaces. \\
Entregable & Diseño de red, enlaces y zonas de seguridad de la faena. \\
Criterio de aceptación & Diseño de zonas aprobado según la norma IEC 62443-3-2. Ningún flujo de datos va desde la red corporativa hacia la red de control de proceso. \\
Responsable & Benjamín Rodríguez (AS). \\
Hitos asociados & M2–M4. Aporta a H2 (mes 4). \\
Esfuerzo estimado & 72 HH (AS 56 y ESI 16). PERT media. \\
Referencias & Subdocumento 3, 3.3.4, RT-03.17 a 03.24, IEC 62443-3-2 (International Electrotechnical Commission [IEC], 2020) e IEC 62264-1, modelo ISA-95 (IEC, 2013). T-12: RT-03.04 y 03.17 a 03.24; RT-05.23, RT-08.03, RNF-INT-005. \\
1.3.1.4 & Plan de continuidad y recuperación ante desastres \\
Descripción del trabajo & Analizar el impacto de una falla en el negocio, fijar los tiempos de recuperación de cada servicio y escribir los procedimientos de conmutación y retorno. \\
Entregable & Plan de continuidad y recuperación ante desastres, con sus procedimientos. \\
Criterio de aceptación & La Contraparte Técnica aprueba el plan antes de H7. Servicios críticos de nube: recuperación en 4 horas o menos y pérdida máxima de 15 minutos. Nodo de Faena opera 48 horas sin enlace. \\
Responsable & Benjamín Rodríguez (AS). \\
Hitos asociados & M4–M15, antes de H7. \\
Esfuerzo estimado & 80 HH (AS 48 y BCP ext. 32). PERT media. \\
Referencias & RT-07.01 a 07.03, 07.05, 07.06 y 07.08; RT-03.10 y 03.13 (Bases Técnicas, pp. 24 y 25); RT-10.03, 10.04 y 10.08; Subdocumento 4, 4.2.3.4; ISO 22301:2019 (ISO, 2019); ISO/IEC 27031:2011 (ISO/IEC, 2011). T-12: RT-03.01; RNF-OPE-005, RP-009. \\
1.3.1.5 & Estudio de cobertura de la red del rajo \\
Descripción del trabajo & Medir cobertura en rajo, botaderos y rampas, y especificar lo que falta. No incluye construir la red. \\
Entregable & Estudio de cobertura del rajo, con zonas sin señal y proyección. \\
Criterio de aceptación & La Contraparte Técnica aprueba el estudio en H2. Se actualiza antes de las detenciones de septiembre de 2027, marzo de 2028 y septiembre de 2028. \\
Responsable & Benjamín Rodríguez (AS). \\
Hitos asociados & M2–M4, M8, M14 y M20. Aporta a H2 (mes 4). \\
Esfuerzo estimado & 152 HH (AS 40 y RFE ext. 112). PERT alta. \\
Referencias & Bases Técnicas, cap. 15, RT-03.24 (p. 25), y Subdocumento 3, 3.2.3.1. T-12: RT-03.23 y 03.24; RP-004. \\
\end{tabla}
\begin{tabla}{I{0.269}I{0.731}}{Diccionario del entregable 1.3.2, Modelo de datos.}{tab:diccionario-entregable-1-3-2-modelo}
  {\cab{1.3.2} & \cab{Modelo de datos}}
  {Fuente: elaboración propia a partir de la Guía del PMBOK (PMI, 2017, p. 162).}
Descripción del trabajo & Modelo de datos de la solución, con sus dominios, su diccionario y sus reglas de retención. Supone que la volumetría del Subdocumento 4 se confirma en el levantamiento. \\
Responsable & Benjamín Velásquez (LDA). \\
Hitos asociados & M2–M4. Aporta a H2 (mes 4). \\
Esfuerzo estimado & 248 HH, 8 externas (LDA 168, ADT 72 y ABG ext. 8). \\
1.3.2.1 & Modelo de datos por dominio \\
Descripción del trabajo & Modelar cada dominio, definir los eventos entre módulos y revisar el modelo con el Arquitecto. \\
Entregable & Modelo de datos de las nueve áreas, con relaciones y eventos. \\
Criterio de aceptación & Modelo de las nueve áreas aprobado en H2, con un responsable por área. Ningún módulo lee datos de otro sin pasar por su contrato. \\
Responsable & Benjamín Velásquez (LDA). \\
Hitos asociados & M2–M4. Aporta a H2 (mes 4). \\
Esfuerzo estimado & 144 HH (LDA 108 y ADT 36). PERT media. \\
Referencias & Subdocumento 5, 5.1, RT-02.13 y RT-05.01. T-12: RT-02.03 y 02.13; RT-05.01, 05.02, 05.05 y 05.09. \\
1.3.2.2 & Diccionario de datos y clasificación de la información \\
Descripción del trabajo & Documentar los atributos, clasificar la información y validar con los propietarios. \\
Entregable & Diccionario de atributos, dominios de valores y clasificación de sensibilidad. \\
Criterio de aceptación & Diccionario completo: nombre, tipo, valores permitidos, obligatoriedad, dueño y sensibilidad de cada atributo. Los datos personales sensibles quedan marcados para ocultarlos fuera de producción. \\
Responsable & Benjamín Velásquez (LDA). \\
Hitos asociados & M2–M4. Aporta a H2 (mes 4). \\
Esfuerzo estimado & 72 HH (LDA 36 y ADT 36). PERT media. \\
Referencias & Subdocumento 5, 5.1.3 y Anexo A, y RT-05.01. T-12: RT-05.01, 05.08 y 05.10; RT-11.03. \\
1.3.2.3 & Política de retención y eliminación \\
Descripción del trabajo & Fijar los plazos legales y operativos, definir el borrado y aprobar la política. \\
Entregable & Plazos de guardado por tipo de información y procedimiento de borrado seguro. \\
Criterio de aceptación & Política aprobada, con plazo por categoría: 10 años para seguridad y salud, ambiente y trazabilidad; 5 para operacional, 7 para tributaria. Trazas técnicas: 90 días. Cada borrado con certificado (RT-05.07). \\
Responsable & Benjamín Velásquez (LDA). \\
Hitos asociados & M3–M4. Aporta a H2 (mes 4). \\
Esfuerzo estimado & 32 HH (LDA 24 y ABG ext. 8). PERT media. \\
Referencias & Subdocumento 5, 5.2.4, 5.2.5 y RT-05.07. Ley N° 21.719 (2024). RT-05.10, Bases Técnicas, p. 25. T-12: RT-16.10 y 16.16; RF-DAT-002, RF-CUM-006, RF-ADM-001, RNF-CUM-001, RNF-SEG-002 y 005. \\
\end{tabla}
\begin{tabla}{I{0.269}I{0.731}}{Diccionario del entregable 1.3.3, Plan de seguridad, amenazas y respuesta a incidentes.}{tab:diccionario-entregable-1-3-3-plan}
  {\cab{1.3.3} & \cab{Plan de seguridad, amenazas y respuesta a incidentes}}
  {Fuente: elaboración propia a partir de la Guía del PMBOK (PMI, 2017, p. 162).}
Descripción del trabajo & Plan de seguridad de la solución: confianza cero, amenazas, controles, identidad y respuesta a incidentes. Supone el directorio corporativo del cliente disponible para la federación de identidades. \\
Responsable & David Martínez (ESI). \\
Hitos asociados & M2–M4. Aporta a H2 (mes 4). \\
Esfuerzo estimado & 240 HH, todas propias (AS 16, ESI 200 y ASG 24). \\
1.3.3.1 & Arquitectura de seguridad de confianza cero \\
Descripción del trabajo & Definir dónde se verifica quién entra, diseñar el cifrado y documentar la arquitectura de seguridad. \\
Entregable & Arquitectura de seguridad de confianza cero, según NIST SP 800-207. \\
Criterio de aceptación & Arquitectura aprobada en H2. Quedan expuestos el Portal de Contratistas, el Portal para Compradores y el Portal Web SITCO. Tránsito con TLS 1.3; no se admiten TLS 1.0 ni 1.1. \\
Responsable & David Martínez (ESI). \\
Hitos asociados & M2–M4. Aporta a H2 (mes 4). \\
Esfuerzo estimado & 72 HH (AS 8 y ESI 64). PERT media. \\
Referencias & RT-11.01. T-12: RT-02.03; RT-03.22, RT-05.08, RT-11.01, 11.08, 11.09 y 11.13, RNF-SEG-001 y 009. \\
1.3.3.2 & Modelo de amenazas y matriz de controles \\
Descripción del trabajo & Modelar las amenazas, asignar controles y registrar el riesgo residual. \\
Entregable & Modelo de amenazas por componente y matriz de controles. \\
Criterio de aceptación & Modelo con las amenazas y controles de cada componente, y su riesgo residual aceptado por el cliente. La matriz indica implementación y evidencia de cada control. \\
Responsable & David Martínez (ESI). \\
Hitos asociados & M2–M4. Aporta a H2 (mes 4). \\
Esfuerzo estimado & 72 HH (ESI 56 y ASG 16). PERT media. \\
Referencias & RT-11.02, 11.05. ISO/IEC 27001 (ISO/IEC, 2022a); ISO/IEC 27002 (ISO/IEC, 2022b), ISO/IEC 27017 (ISO/IEC, 2015); ISO/IEC 27018 (ISO/IEC, 2019). T-12: RT-11.02, 11.05 y 11.06; RT-26.07. \\
1.3.3.3 & Modelo de identidad y acceso \\
Descripción del trabajo & Definir los roles, diseñar la federación y definir las políticas de sesión. \\
Entregable & Modelo de identidad con roles, doble factor y sesiones. \\
Criterio de aceptación & Modelo aprobado en H2. El doble factor rige para todo el personal. La baja de cada identidad se hace en 24 horas o menos desde la desvinculación. \\
Responsable & David Martínez (ESI). \\
Hitos asociados & M2–M4. Aporta a H2 (mes 4). \\
Esfuerzo estimado & 56 HH (AS 8 y ESI 48). PERT media. \\
Referencias & Cap. 12 de las Bases técnicas transversales (CMAA, 2026c, Art. 22 y Bases Técnicas, cap. 10, restricción 9. T-12: RT-03.22); RT-11.01 y 11.27; RT-12.01 a 12.11 y 12.13; RF-ACR-012 y 016; RF-ADM-007, RNF-SEG-003, 006 y 012 a 014. \\
1.3.3.4 & Plan de respuesta a incidentes \\
Descripción del trabajo & Definir los tipos de incidente, asignar roles, fijar la comunicación al cliente y a la autoridad. \\
Entregable & Plan de respuesta a incidentes con roles, escalamiento y comunicación. \\
Criterio de aceptación & Plan aprobado, con severidad, escalamiento, plazos y responsables. El cliente recibe aviso dentro de 2 horas de un incidente crítico. Toda brecha de datos personales se notifica en 24 horas con informe preliminar. \\
Responsable & David Martínez (ESI). \\
Hitos asociados & M4. Aporta a H2 (mes 4). \\
Esfuerzo estimado & 40 HH (ESI 32 y ASG 8). PERT media. \\
Referencias & RT-11.18. T-12: RT-11.18 y 11.19. \\
\end{tabla}
\begin{tabla}{I{0.269}I{0.731}}{Diccionario del entregable 1.3.4, Investigación de usuarios, sistema de diseño y prototipos.}{tab:diccionario-entregable-1-3-4-investigacion}
  {\cab{1.3.4} & \cab{Investigación de usuarios, sistema de diseño y prototipos}}
  {Fuente: elaboración propia a partir de la Guía del PMBOK (PMI, 2017, p. 162).}
Descripción del trabajo & Investigación con usuarios en terreno, sistema de diseño y prototipos navegables de los canales. Supone acceso a operadores, supervisores y portería en ambos grupos del 7x7. \\
Responsable & Patricio Figueroa (LF). \\
Hitos asociados & M3–M7. Sin hito propio. \\
Esfuerzo estimado & 440 HH, 248 externas (LF 96, DEV 96 y UXR ext. 248). \\
1.3.4.1 & Informe de investigación con usuarios \\
Descripción del trabajo & Observar los turnos, entrevistar por perfil y sintetizar los hallazgos. \\
Entregable & Hallazgos de la observación y entrevistas en terreno, por perfil de usuario. \\
Criterio de aceptación & Informe con los hallazgos de cada perfil de la Etapa 1, observados en ambos grupos del 7x7. Cada hallazgo se traza a un requerimiento o a una decisión de diseño. \\
Responsable & Patricio Figueroa (LF). \\
Hitos asociados & M3–M4. \\
Esfuerzo estimado & 80 HH (LF 32 y UXR ext. 48). PERT media. \\
Supuestos & La observación se hace fuera de los seis días de detención de marzo de 2027. \\
Referencias & RT-13.03. \\
1.3.4.2 & Sistema de diseño documentado \\
Descripción del trabajo & Definir los componentes, documentar su uso y publicar la biblioteca. \\
Entregable & Reglas de pantalla con colores, tipografía, íconos y componentes reutilizables. \\
Criterio de aceptación & Biblioteca publicada y usada por los prototipos, con contraste y tamaños conformes con WCAG 2.2 nivel AA. Ninguna función principal a más de tres toques. Usable con guantes y entre −6 °C y 34 °C. \\
Responsable & Patricio Figueroa (LF). \\
Hitos asociados & M3–M5. \\
Esfuerzo estimado & 72 HH (LF 16 y UXR ext. 56). PERT media. \\
Referencias & RT-13.09 y Bases Técnicas, RT-13.08 (p. 25). WCAG 2.2 (World Wide Web Consortium, 2023). T-12: RT-13.02, 13.05 a 13.07 y 13.09 a 13.11; RT-25.05, RNF-USA-006, 010 y 012 a 015. \\
1.3.4.3 & Prototipos navegables de los canales \\
Descripción del trabajo & Construir prototipos de los seis canales y probarlos con al menos cinco usuarios por perfil. \\
Entregable & Prototipos navegables de los seis canales, incluidos los tres portales. \\
Criterio de aceptación & Prototipo de cada canal de la Etapa 1, validado con al menos cinco usuarios del perfil. Tres interacciones como máximo en las tareas críticas de terreno. El prototipo del Portal para Compradores se valida con cinco usuarios del perfil antes del mes 7. El lote piloto del mes 17 (1.2.5.5) lo verifica con datos reales. \\
Responsable & Patricio Figueroa (LF). \\
Hitos asociados & M4–M7. \\
Esfuerzo estimado & 288 HH (LF 48, DEV 96 y UXR ext. 144). PERT media. \\
Referencias & Subdocumento 4, 4.1.5. T-12: RT-13.03, 13.05 y 13.12; RT-25.01 a 25.04, 25.06 y 25.08 a 25.10. \\
\end{tabla}
\begin{tabla}{I{0.269}I{0.731}}{Diccionario del entregable 1.3.5, Diseño detallado de la Etapa 2.}{tab:diccionario-entregable-1-3-5-diseno}
  {\cab{1.3.5} & \cab{Diseño detallado de la Etapa 2}}
  {Fuente: elaboración propia a partir de la Guía del PMBOK (PMI, 2017, p. 162).}
Descripción del trabajo & Diseño de los módulos y conectores de la Etapa 2, sin cambiar la base de la Etapa 1. Supone que la Etapa 2 no cambia la arquitectura. \\
Responsable & Benjamín Rodríguez (AS). \\
Hitos asociados & M13–M14. Aporta a H8 (mes 14). \\
Esfuerzo estimado & 288 HH, todas propias (AS 144, LD 48 y LI 96). \\
1.3.5.1 & Diseño de los módulos de la Etapa 2 \\
Descripción del trabajo & Diseñar los módulos, definir sus contratos y revisar en el Comité de Arquitectura. \\
Entregable & Diseño funcional y técnico de los módulos de la Etapa 2. \\
Criterio de aceptación & Diseño aprobado en H8 sin cambios a la arquitectura de la Etapa 1. Mismas reglas de acoplamiento de la Etapa 1. \\
Responsable & Benjamín Rodríguez (AS). \\
Hitos asociados & M13–M14. Aporta a H8 (mes 14). \\
Esfuerzo estimado & 192 HH (AS 144 y LD 48). PERT media. \\
Referencias & Subdocumento 4, 4.1.6.2, 4.1.6.3. \\
1.3.5.2 & Diseño de los conectores de la Etapa 2 \\
Descripción del trabajo & Definir los contratos de cada conector y validar con los dueños de sistema. \\
Entregable & Diseño de interfaces de la Etapa 2, con un adaptador por fabricante. \\
Criterio de aceptación & Cada contrato aprobado por el dueño del sistema. Todas las interfaces se publican en OpenAPI 3.1 o AsyncAPI 3.0, con versionado semántico y compatibilidad hacia atrás. Aviso de seis meses antes de retirar una versión. \\
Responsable & Marcelo Flores (LI). \\
Hitos asociados & M13–M14. Aporta a H8 (mes 14). \\
Esfuerzo estimado & 96 HH del LI. PERT alta. \\
Referencias & Subdocumento 4, Tabla \refext{S4}{tab:conectores-integrador-sistemas-etapa-lugar-ejecucion} y \refext{S4}{tab:selecciones-principales-capas-1-3}.7.4 (OpenAPI Initiative, 2021; AsyncAPI Initiative, 2023). T-12: RF-EQU-003; RNF-INT-006. \\
\end{tabla}
\begin{tabla}{I{0.269}I{0.731}}{Diccionario del entregable 1.3.6, Matriz de cumplimiento normativo y de datos personales.}{tab:diccionario-entregable-1-3-6-matriz}
  {\cab{1.3.6} & \cab{Matriz de cumplimiento normativo y de datos personales}}
  {Fuente: elaboración propia a partir de la Guía del PMBOK (PMI, 2017, p. 162).}
Descripción del trabajo & Obligaciones legales que aplican a la solución, con el tratamiento de los datos personales y la biometría. Supone la Ley N° 21.719 vigente durante el contrato. \\
Responsable & David Martínez (ESI). \\
Hitos asociados & M2–M4. Aporta a H2 (mes 4). \\
Esfuerzo estimado & 120 HH, suma de sus paquetes, de las que 56 HH son externas. Se reparten en ESI 64 HH y ABG externo 56 HH. \\
1.3.6.1 & Matriz de cumplimiento normativo \\
Descripción del trabajo & Identificar cada obligación legal y normativa que aplica a SITCO Aranda (Art. 27), asignarle un control y una evidencia. \\
Entregable & Matriz con cada obligación legal, su control y su evidencia. \\
Criterio de aceptación & Matriz aprobada con un control y una evidencia por obligación. Revisada por el Encargado de Seguridad de la Información y el área legal del cliente. \\
Responsable & David Martínez (ESI). \\
Hitos asociados & M2–M4. Aporta a H2 (mes 4). \\
Esfuerzo estimado & 56 HH (ESI 32 y ABG externo 24). PERT media. \\
Referencias & Art. 27. \\
1.3.6.2 & Registro de tratamiento y evaluación de impacto \\
Descripción del trabajo & Registrar cada tratamiento, evaluar el impacto y definir las medidas. \\
Entregable & Registro de tratamientos y evaluación de impacto de la biometría. \\
Criterio de aceptación & Registro y evaluación aprobados antes de diseñar los controles biométricos. Cada tratamiento declara su base legal. La copia en Brazil South tiene aprobación del cliente y cláusulas de transferencia. \\
Responsable & David Martínez (ESI). \\
Hitos asociados & M2–M4. Aporta a H2 (mes 4). \\
Esfuerzo estimado & 64 HH (ESI 32 y ABG ext. 32). PERT media. \\
Referencias & Art. 27, 85, RT-05.08 y RT-06.20. Ley N° 21.719 (2024). T-12: RNF-CUM-001. \\
\end{tabla}

\apartado{Rama 1.4 Software del sistema}
La rama 1.4 tiene 6 entregables y 28 paquetes de trabajo, que suman 21.062 HH. Los módulos y canales se construyen con los tres equipos Scrum en sprints de dos semanas, los ambientes de Desarrollo y QA y el flujo de entrega continua, el mes 11 corrige los defectos de las pruebas de la Etapa 1 sin agregar requerimientos. Las tablas \ref{tab:diccionario-entregable-1-4-1-flujo} a \ref{tab:diccionario-entregable-1-4-6-datos} presentan la ficha de cada entregable y, bajo ella, la de cada uno de sus paquetes.

\begin{tabla}{I{0.269}I{0.731}}{Diccionario del entregable 1.4.1, Flujo de integración y entrega continua.}{tab:diccionario-entregable-1-4-1-flujo}
  {\cab{1.4.1} & \cab{Flujo de integración y entrega continua}}
  {Fuente: elaboración propia a partir de la Guía del PMBOK (PMI, 2017, p. 162).}
Descripción del trabajo & Flujo automático que construye, prueba y despliega el software. El repositorio queda en la organización del cliente desde el primer día. \\
Responsable & Gabriel Sanzana (LD). \\
Hitos asociados & M1–M21. Aporta a H3 (mes 6). \\
Esfuerzo estimado & 1.710 HH, todas propias (ESI 24, LD 1.614 y DEV 72). \\
1.4.1.1 & Repositorio y flujo de integración continua \\
Descripción del trabajo & Crear la organización del cliente, proteger la rama principal y configurar los análisis. \\
Entregable & Repositorio de código con revisión por pares y análisis automáticos. \\
Criterio de aceptación & La Contraparte Técnica recibe el repositorio cuando muestra la rama principal protegida y análisis que bloquean hallazgos críticos. Umbrales de la puerta 1 del Subdocumento 9 activos. \\
Responsable & Gabriel Sanzana (LD). \\
Hitos asociados & M1–M2. Aporta a H3 (mes 6). \\
Esfuerzo estimado & 48 HH (LD 24 y DEV 24). PERT media. \\
Referencias & RT-04.03, 04.05 y RT-11.22. T-12: RT-11.22, 11.26 y 11.27. \\
1.4.1.2 & Flujo de entrega continua con reversión \\
Descripción del trabajo & Configurar la entrega automática en la nube, en el Nodo de Faena y en terreno, y probar la reversión. \\
Entregable & Entrega automática por etapas, con reversión y software firmado. \\
Criterio de aceptación & La Contraparte Técnica recibe el flujo cuando una prueba en Preproducción lo promueve y lo revierte solo.La reversión se completa en no más de 30 minutos por función y área, sin registros perdidos. \\
Responsable & Gabriel Sanzana (LD). \\
Hitos asociados & M2–M3. Aporta a H3 (mes 6). \\
Esfuerzo estimado & 64 HH (LD 32 y DEV 32). PERT media. \\
Referencias & RT-04.06, 04.07. T-12: RT-04.06 a 04.08, 04.10, 04.12 y 04.14; RT-10.06, RNF-DIS-008, TRA-006. \\
1.4.1.3 & Gestor de secretos y firma de artefactos \\
Descripción del trabajo & Configurar el gestor, generar la lista de componentes y firmar los artefactos. \\
Entregable & Almacén de claves, lista de componentes y firma de programas. \\
Criterio de aceptación & Almacén de claves con cambio automático y registro de cada acceso. Todo programa desplegado está firmado, con lista de componentes CycloneDX o SPDX y procedencia SLSA de nivel 3. \\
Responsable & David Martínez (ESI). \\
Hitos asociados & M2–M3. Aporta a H3 (mes 6). \\
Esfuerzo estimado & 40 HH (ESI 24 y DEV 16). PERT media. \\
Referencias & RT-04.09 y RT-11.23, 11.24. T-12: RNF-POR-003; RNF-SEG-009. \\
1.4.1.4 & Coordinación técnica de los equipos de desarrollo \\
Descripción del trabajo & Definir los estándares de codificación, conducir la coordinación de los equipos y resolver sus impedimentos técnicos. \\
Entregable & Estándares de codificación, revisiones técnicas y actas de coordinación. \\
Criterio de aceptación & La Contraparte Técnica recibe los estándares publicados en el repositorio y aplicados en la puerta 1. Cada dependencia entre equipos queda registrada con responsable y fecha. \\
Responsable & Gabriel Sanzana (LD). \\
Hitos asociados & M1–M21. Aporta a H4 (mes 10) y H9 (mes 17). \\
Esfuerzo estimado & 1.558 HH del LD, saldo de su dedicación en los meses 1 a 21. Sin PERT. \\
Supuestos & Tres equipos trabajan con un solo listado de pendientes y un solo flujo de entrega. \\
Referencias & Art. 91 y Subdocumento 6, 6.2. T-12: RT-05.16 y 05.17; RT-11.23, RP-005. \\
\end{tabla}
\begin{tabla}{I{0.269}I{0.731}}{Diccionario del entregable 1.4.2, Plataformas de integración y de faena.}{tab:diccionario-entregable-1-4-2-plataformas}
  {\cab{1.4.2} & \cab{Plataformas de integración y de faena}}
  {Fuente: elaboración propia a partir de la Guía del PMBOK (PMI, 2017, p. 162).}
Descripción del trabajo & Integrador de Sistemas con los conectores de ambas etapas y Nodo de Faena que opera sin enlace. Supone que los sistemas existentes permiten leer por interfaz; si no, exportación programada. \\
Responsable & Marcelo Flores (LI). \\
Hitos asociados & M3–M17. Aporta a H4 (mes 10) y H9 (mes 17). \\
Esfuerzo estimado & 3.536 HH, todas propias (LI 810 y DEV 2.726). \\
1.4.2.1 & Integrador de Sistemas con conectores de la Etapa 1 \\
Descripción del trabajo & Construir los conectores de la Etapa 1 sin modificar esos sistemas, y probar sus contratos. Incluye los periféricos. \\
Entregable & Conexiones con siete sistemas de la Etapa 1 y radio de faena. \\
Criterio de aceptación & Conectores con pruebas de contrato aprobadas contra el simulador y el sistema real. La prueba con el laboratorio envía 200 muestras sin pérdidas ni duplicados. Ninguno de 100 intentos de escribir en el historiador pasa el diodo. \\
Responsable & Marcelo Flores (LI). \\
Hitos asociados & M3–M11. Aporta a H4 (mes 10). \\
Esfuerzo estimado & 1.520 HH (LI 342 y DEV 1.178). PERT alta. \\
Referencias & Subdocumento 4, Tabla \refext{S4}{tab:conectores-integrador-sistemas-etapa-lugar-ejecucion}; Bases Técnicas, RT-17.06. T-12: RT-02.06, 02.07 y 02.14; RT-05.18 a 05.20 y 05.23; RT-10.08, RF-TRZ-001 y 018; RF-CON-001, 009, 012, 015 y 017; RF-TER-003 y 004; RF-PLA-001 y 009; RF-NOT-001, RNF-DIS-007, RNF-INT-001 a 004 y 006; RNF-SEG-004 y 007; RNF-TRV-02.06 y 02.07. \\
1.4.2.2 & Nodo de Faena con operación sin enlace \\
Descripción del trabajo & Construir los servicios del Nodo de Faena y de sus cuatro gabinetes, y bloquear accesos a distancia en 15 minutos. \\
Entregable & Nodo de Faena que opera sin conexión y se sincroniza al volver. \\
Criterio de aceptación & Corte de 48 horas sin conexión aprobado, con sincronización en 4 horas o menos y sin pérdida de registros. Las pantallas de terreno responden en 2 segundos o menos y la portería en 1,5 segundos. \\
Responsable & Marcelo Flores (LI). \\
Hitos asociados & M3–M11. Aporta a H4 (mes 10). \\
Esfuerzo estimado & 1.216 HH (LI 228 y DEV 988). PERT media. \\
Referencias & Subdocumento 3, 3.3.6, Subdocumento 4, Tabla \refext{S4}{tab:modulo-plataforma-sostiene-cada-funcion-durante}, y Bases Técnicas, RT-03.10 a 03.13 (pp. 24 y 25). T-12: RT-03.10 a 03.13, 03.16 y 03.19; RT-11.11, RT-12.10, RT-17.01, RF-TER-005, RF-EQU-006, RF-CNX-001, 002, 004 y 005; RNF-CAP-006, RNF-DES-001, 003 y 012; RNF-DIS-001, 003 y 006; RNF-SEG-006, RNF-OPE-002, RNF-TRV-11.11. \\
1.4.2.3 & Integrador de Sistemas con conectores de la Etapa 2 \\
Descripción del trabajo & Construir las conexiones de la Etapa 2 para telemetría y control de fatiga, sin modificar SAP S/4HANA. \\
Entregable & Conectores de la Etapa 2 para telemetría, fatiga, mantenimiento y portería. \\
Criterio de aceptación & Pruebas de contrato y regresión de la Etapa 1 sin fallas. La ampliación de portería entra sólo con acuerdo firmado con el sindicato. \\
Responsable & Marcelo Flores (LI). \\
Hitos asociados & M13–M17. Aporta a H9 (mes 17). \\
Esfuerzo estimado & 800 HH (LI 240 y DEV 560). PERT alta. \\
Referencias & Subdocumento 4, Tabla \refext{S4}{tab:conectores-integrador-sistemas-etapa-lugar-ejecucion}; RT-05.18, RT-05.20 y RT-02.06. T-12: RF-CON-013, RF-CUM-005, RF-EQU-002 y 004; RNF-TRV-02.08. \\
\end{tabla}
\begin{tabla}{I{0.269}I{0.731}}{Diccionario del entregable 1.4.3, Módulos de la Etapa 1.}{tab:diccionario-entregable-1-4-3-modulos}
  {\cab{1.4.3} & \cab{Módulos de la Etapa 1}}
  {Fuente: elaboración propia a partir de la Guía del PMBOK (PMI, 2017, p. 162).}
Descripción del trabajo & Módulos de software de la Etapa 1, que entran en producción en el mes 16. Supone tres equipos que ejecutan 19 sprints de dos semanas entre los meses 3 y 11. \\
Responsable & Gabriel Sanzana (LD). \\
Hitos asociados & M3–M17. Aporta a H4 (mes 10), H5 (mes 12) y H7 (mes 16). \\
Esfuerzo estimado & 6.120 HH, todas propias (LDA 224, LD 752, DEV 4.888 y ADT 256). \\
1.4.3.1 & Módulo de Administración y Servicios Comunes \\
Descripción del trabajo & Construir datos maestros, configuración, roles, notificaciones y firma electrónica, y probarlos en QA. \\
Entregable & Módulo de administración con datos, roles, notificaciones y firma electrónica. \\
Criterio de aceptación & Módulo aprobado en QA para datos maestros, configuración, roles, notificaciones y firma. Una unidad operacional nueva se configura sin cambiar código en 20 días hábiles o menos. \\
Responsable & Gabriel Sanzana (LD). \\
Hitos asociados & M3–M8. Aporta a H4 (mes 10). \\
Esfuerzo estimado & 768 HH (LD 120 y DEV 648). PERT media. \\
Referencias & Subdocumento 3, 3.4.3 y RT-16.01 a 16.16. Bases técnicas transversales, RT-16.14 (p. 26). T-12: RT-02.12, RT-05.03, 05.06 y 05.22; RT-16.01 a 16.23 y 16.25 a 16.30; RF-TRZ-017, RF-DAT-001 a 003; RF-ACR-004 y 014; RF-ADM-002 a 013; RF-NOT-001 a 003; RNF-CUM-002, RNF-DES-018, RNF-POR-001 y 002. \\
1.4.3.2 & Módulo de Operación Mina \\
Descripción del trabajo & Construir por sprint el ciclo de carguío y transporte y la declaración de material, e integrar con despacho. \\
Entregable & Módulo de operación mina con carguío, transporte y declaración de material. \\
Criterio de aceptación & Módulo con pruebas aprobadas en QA, también sin conexión. Cálculo crítico con 90 \% de cobertura de pruebas. Desde el mes 16, el 100 \% de las detenciones tiene causa dentro del turno siguiente. \\
Responsable & Gabriel Sanzana (LD). \\
Hitos asociados & M4–M11. Aporta a H4 (mes 10). \\
Esfuerzo estimado & 1.280 HH (LD 160 y DEV 1.120). PERT media. \\
Referencias & Subdocumento 3, 3.4.2.1. T-12: RT-05.29; RF-TRZ-001, 002, 010 y 011; RF-TER-001, 003, 004 y 006 a 009; RF-PLA-001 a 006, 008 y 010; RNF-DES-008 y 009; RNF-USA-004 y 009. \\
1.4.3.3 & Módulo de Trazabilidad y Embarque \\
Descripción del trabajo & Construir stock, lotes, despacho y embarque, conectarlos con báscula y laboratorio, y probarlos en QA. \\
Entregable & Módulo de trazabilidad con stock, lotes, despacho y embarque. \\
Criterio de aceptación & Módulo recibido cuando el origen de un lote se reconstruye en menos de 2 horas. Registro en báscula confirma en 3 segundos o menos, medido en 200 pesajes seguidos. \\
Responsable & Gabriel Sanzana (LD). \\
Hitos asociados & M5–M11. Aporta a H4 (mes 10). \\
Esfuerzo estimado & 1.064 HH (LD 140 y DEV 924). PERT media. \\
Referencias & Subdocumento 3, 3.4.2.2; Bases Técnicas, 16.2; SUP-03, SUP-11 y SUP-21. T-12: RF-TRZ-002 a 009, 012 a 015 y 017 a 019; RNF-DES-002 y 005. \\
1.4.3.4 & Módulo de Conciliación Metalúrgica \\
Descripción del trabajo & Construir por sprint el balance de turno y la conciliación mensual, y cargar los factores versionados. \\
Entregable & Módulo de conciliación metalúrgica con balance de turno y conciliación mensual. \\
Criterio de aceptación & Módulo recibido cuando la conciliación de un período de prueba coincide con el cálculo de referencia del cliente. Toda diferencia queda explicada. El cierre de prueba se completa en 3 días hábiles o menos. \\
Responsable & Gabriel Sanzana (LD). \\
Hitos asociados & M6–M11. Aporta a H4 (mes 10). \\
Esfuerzo estimado & 960 HH (LDA 96, LD 96 y DEV 768). PERT media. \\
Supuestos & Hasta que el conector del historiador pasa sus pruebas, el módulo usa el simulador y extractos enmascarados. \\
Referencias & Subdocumento 3, 3.4.2.3. T-12: RF-CON-001 a 008, 010, 011, 012 y 016; RNF-DES-006 y 014; TRA-002. \\
1.4.3.5 & Módulo de Acreditación y Acceso \\
Descripción del trabajo & Construir por ciclo la habilitación de personas, la credencial y el control en portería, integrándolo con los torniquetes. \\
Entregable & Módulo de acreditación y acceso con credencial, control en portería y bitácora. \\
Criterio de aceptación & Pruebas aprobadas en QA. Un contratista nuevo queda acreditado en menos de 48 horas desde su solicitud completa (CMAA, 2026b, criterio 8). Nadie es rechazado en portería por documento vencido que el sistema pudo advertir (CMAA, 2026b, criterio 7). \\
Responsable & Gabriel Sanzana (LD). \\
Hitos asociados & M4–M10. Aporta a H4 (mes 10). \\
Esfuerzo estimado & 896 HH (LD 140 y DEV 756). PERT media. \\
Referencias & Subdocumento 3, 3.4.2.5. T-12: RT-06.20, 06.21 y 06.23 a 06.25; RT-11.10, RF-ACR-001 a 004, 006 a 008, 012, 013, 017 y 018; RNF-DES-015. \\
1.4.3.6 & Indicadores y Centro Integrado: tableros de turno y vista del CIO \\
Descripción del trabajo & Construir los modelos analíticos, diseñar los tableros y validarlos con los usuarios. \\
Entregable & Tableros de turno y vista del Centro Integrado con los indicadores. \\
Criterio de aceptación & La Contraparte Técnica recibe los tableros cuando los indicadores de turno tienen latencia de 15 minutos o menos en QA. Se publica la pérdida del cambio de turno frente a la meta de 25 minutos. \\
Responsable & Benjamín Velásquez (LDA). \\
Hitos asociados & M8–M11. Aporta a H4 (mes 10). \\
Esfuerzo estimado & 384 HH (LDA 128 y ADT 256). PERT media. \\
Referencias & RT-05.25 a 05.29 y Bases Técnicas, 9.7 (p. 19). T-12: RT-05.04 y 05.25 a 05.29; RT-14.02, RF-PLA-008 y 010; RF-ANA-001, RNF-DES-010 y 017; RNF-OPE-002. \\
1.4.3.7 & Defectos de la certificación, la marcha blanca y la estabilización de la Etapa 1 corregidos \\
Descripción del trabajo & Clasificar los defectos por gravedad, corregirlos con el mismo flujo de entrega y verificar que no reaparecen. No agrega requerimientos. \\
Entregable & Correcciones desplegadas de los defectos de certificación, marcha blanca y estabilización. \\
Criterio de aceptación & Correcciones sin defectos críticos ni altos abiertos al cierre de la marcha blanca. Pasan las puertas de calidad y la regresión. \\
Responsable & Gabriel Sanzana (LD). \\
Hitos asociados & M12–M17. Aporta a H5 (mes 12) y H7 (mes 16). \\
Esfuerzo estimado & 768 HH (LD 96 y DEV 672). PERT media. \\
Supuestos & Los especialistas de nivel 2 quedan disponibles en ambos turnos durante la estabilización. \\
Referencias & RT-20.06 y Subdocumento 9, 9.2.6. \\
\end{tabla}
\begin{tabla}{I{0.269}I{0.731}}{Diccionario del entregable 1.4.4, Módulos de la Etapa 2.}{tab:diccionario-entregable-1-4-4-modulos}
  {\cab{1.4.4} & \cab{Módulos de la Etapa 2}}
  {Fuente: elaboración propia a partir de la Guía del PMBOK (PMI, 2017, p. 162).}
Descripción del trabajo & Módulos de software de la Etapa 2 en producción en el mes 21, con el desarrollo cerrado en el mes 18. \\
Responsable & Gabriel Sanzana (LD). \\
Hitos asociados & M13–M21. Aporta a H9 (mes 17), H10 (mes 18) y H12 (mes 21). \\
Esfuerzo estimado & 3.552 HH, todas propias (ESI 48, LDA 48, LD 516, LI 64, DEV 2.540 y ADT 336). \\
1.4.4.1 & Módulo de Equipos y Mantenimiento \\
Descripción del trabajo & Construir avisos, pautas, alertas y registro de paradas de planta, con SAP y telemetría. No incluye modelos predictivos. \\
Entregable & Módulo de Equipos y Mantenimiento con avisos, pautas, alertas y paradas. \\
Criterio de aceptación & Caso de aceptación aprobado en QA y regresión de la Etapa 1 sin fallas. Reglas de clasificación aprobadas en QA. Desde el mes 21, el 100 \% de las detenciones queda clasificado con causa dentro del turno siguiente. \\
Responsable & Gabriel Sanzana (LD). \\
Hitos asociados & M13–M17. Aporta a H9 (mes 17). \\
Esfuerzo estimado & 800 HH (LD 100 y DEV 700). PERT alta. \\
Referencias & Subdocumento 3, 3.4.2.4. Bases Técnicas, cap. 18, criterios 5 y 6. T-12: RT-05.29; RF-PLA-006, 007 y 009; RF-EQU-001, 003 a 005, 007 y 008; RNF-DES-009 a 011; RNF-DIS-005. \\
1.4.4.2 & Módulo de Seguridad y Salud Ocupacional \\
Descripción del trabajo & Construir incidentes, inspecciones y controles de salud, con versión local en el Nodo de Faena. \\
Entregable & Módulo de Seguridad y Salud Ocupacional con incidentes, inspecciones y controles. \\
Criterio de aceptación & Caso aprobado en QA. Cada consulta o cambio de datos de salud, fatiga, alcohol y drogas queda registrado con quién, cuándo y sobre quién; se conserva 10 años. Alerta de fatiga sin acción escala a los 30 minutos. \\
Responsable & Gabriel Sanzana (LD). \\
Hitos asociados & M13–M16. Aporta a H9 (mes 17). \\
Esfuerzo estimado & 512 HH (LD 80 y DEV 432). PERT media. \\
Referencias & Subdocumento 3, 3.4.2.6. Bases Técnicas, RT-16.14 y cap. 18, criterios 9 y 10. T-12: RT-11.10; RT-16.09 y 16.30; RF-ACR-007 y 010; RF-CUM-001 a 007; RNF-SEG-001 y 002; RNF-USA-007. \\
1.4.4.3 & Módulo de Cumplimiento Ambiental \\
Descripción del trabajo & Construir compromisos ambientales, mediciones, informes y huella por tonelada de cobre fino. Sin lotes de los meses 18 a 21. \\
Entregable & Módulo de Cumplimiento Ambiental con compromisos, mediciones, informes y huella. \\
Criterio de aceptación & Caso de aceptación aprobado en QA. Cada compromiso tiene frecuencia, punto de medición y destinatario, y el módulo alerta su vencimiento. La huella mensual está disponible dentro de 5 días hábiles desde el cierre. \\
Responsable & Gabriel Sanzana (LD). \\
Hitos asociados & M14–M17. Aporta a H9 (mes 17). \\
Esfuerzo estimado & 576 HH (LD 88 y DEV 488). PERT media. \\
Referencias & Subdocumento 3, 3.4.2.7. Bases Técnicas, RT-17.06 y cap. 18, criterios 11 y 12. T-12: RT-05.23; RF-CUM-008 a 013 y 015 a 017; RNF-DES-013. \\
1.4.4.4 & Acreditación y Acceso: verificación facial en portería \\
Descripción del trabajo & Construir la verificación facial, integrarla con la consola de portería y probarla con usuarios. \\
Entregable & Verificación facial en portería, con consentimiento y alternativa sin biometría. \\
Criterio de aceptación & Caso de aceptación aprobado en QA. No guarda imágenes, solo una plantilla cifrada. Sin plantilla o con lector fuera de servicio, entra con credencial. Se activa con acuerdo del sindicato. \\
Responsable & Gabriel Sanzana (LD). \\
Hitos asociados & M15–M17. Aporta a H9 (mes 17). \\
Esfuerzo estimado & 384 HH (ESI 48, LD 60 y DEV 276). PERT alta. \\
Referencias & RT-18.01 a 18.09 de las Bases técnicas transversales (CMAA, 2026c, Art. 86 y 1.7.2.1. T-12: RF-ACR-009 y 010. \\
1.4.4.5 & Portal para Compradores \\
Descripción del trabajo & Construir el portal para compradores, probar su carga y desplegarlo en QA. \\
Entregable & Portal para compradores con consulta de trazabilidad de sus lotes. \\
Criterio de aceptación & Requerimientos aprobados en QA, prueba de carga con picos superada sin degradar servicios internos, y acceso de compradores con doble factor); entrega el certificado de análisis de cada lote. \\
Responsable & Gabriel Sanzana (LD). \\
Hitos asociados & M15–M17. Aporta a H9 (mes 17). \\
Esfuerzo estimado & 384 HH (LD 60 y DEV 324). PERT media. \\
Referencias & Bases Técnicas, RT-16.30 (p. 26); Bases técnicas transversales, RT-25.04. T-12: RT-12.12; RT-16.31 y 16.34; RF-TRZ-016, RF-ACR-015. \\
1.4.4.6 & Indicadores y Centro Integrado: indicadores de gestión, costo por tonelada y autoservicio de informes \\
Descripción del trabajo & Construir modelos analíticos y un modelo semántico documentado, diseñar tableros y habilitar informes de autoservicio. \\
Entregable & Indicadores de gestión, costo por tonelada e informes de autoservicio. \\
Criterio de aceptación & El costo por tonelada del día anterior queda disponible en 4 horas o menos desde la entrega de SAP. El tablero de calidad muestra completitud, exactitud y consistencia, con metas de 95 \%, 98 \% y 98 \%. \\
Responsable & Benjamín Velásquez (LDA). \\
Hitos asociados & M15–M17. Aporta a H9 (mes 17). \\
Esfuerzo estimado & 384 HH (LDA 48 y ADT 336). PERT media. \\
Referencias & RT-05.25 a 05.29. T-12: RT-05.04 y 05.25 a 05.29; RF-DAT-001, RF-CON-013 y 014; RF-CUM-014, RF-ANA-001 a 003; RNF-DES-007 y 010; RNF-OPE-002. \\
1.4.4.7 & Defectos de la certificación, la marcha blanca y la estabilización de la Etapa 2 corregidos \\
Descripción del trabajo & Corregir los defectos de la Etapa 2 sin degradar la Etapa 1. No agrega requerimientos nuevos. \\
Entregable & Correcciones desplegadas de defectos de certificación, marcha blanca y estabilización. \\
Criterio de aceptación & Ningún defecto crítico ni alto abierto al cerrar la marcha blanca de la Etapa 2, y la Etapa 1 sigue sin fallas tras cada corrección. \\
Responsable & Gabriel Sanzana (LD). \\
Hitos asociados & M18–M21. Aporta a H10 (mes 18) y H12 (mes 21). \\
Esfuerzo estimado & 512 HH (LD 128, LI 64 y DEV 320). PERT media. \\
Supuestos & Ciclos de dos semanas de corrección de defectos, de los meses 18 a 21, sin requerimientos nuevos. \\
Referencias & RT-20.06 y Art. 17.2. \\
\end{tabla}
\begin{tabla}{I{0.269}I{0.731}}{Diccionario del entregable 1.4.5, Canales de uso.}{tab:diccionario-entregable-1-4-5-canales}
  {\cab{1.4.5} & \cab{Canales de uso}}
  {Fuente: elaboración propia a partir de la Guía del PMBOK (PMI, 2017, p. 162).}
Descripción del trabajo & Aplicaciones con que los usuarios operan SITCO Aranda en terreno, en oficina y en portería. \\
Responsable & Gabriel Sanzana (LD). \\
Hitos asociados & M3–M11. Aporta a H4 (mes 10). \\
Esfuerzo estimado & 2.240 HH, 280 externas (LD 176, DEV 1.784 y UXR ext. 280). \\
1.4.5.1 & Botonera de Pala y App SITCO Terreno \\
Descripción del trabajo & Construir la App SITCO Terreno, para faena, y la Botonera de Pala, para declarar cambios de material. \\
Entregable & Botonera de Pala y App SITCO Terreno, con captura sin conexión. \\
Criterio de aceptación & La Botonera no agrega segundos de pala detenida, medido en 500 ciclos o más. Dispositivo perdido bloqueado o borrado en 15 minutos o menos. Sesión con credencial en 2 segundos o menos, en el 95 \% de intentos. \\
Responsable & Gabriel Sanzana (LD). \\
Hitos asociados & M3–M11. Aporta a H4 (mes 10). \\
Esfuerzo estimado & 1.216 HH (LD 96, DEV 968 y UXR ext. 152). PERT media. \\
Referencias & Subdocumento 3, 3.3.3; Subdocumento 4, ADR-06; RT-17.01 a 17.08. T-12: RT-03.12, 03.18 y 03.19; RT-08.14, RT-12.11, RT-13.05 y 13.08; RF-TER-001, 002, 005 y 006; RF-ACR-011 y 012; RF-CUM-001, RF-EQU-005 y 006; RF-CNX-001, 003 y 004; RNF-DES-004 y 016; RNF-POR-003 y 004; RNF-SEG-016, RNF-USA-001, 002, 008 y 012. \\
1.4.5.2 & Portales y Consola de Portería \\
Descripción del trabajo & Construir los portales, probar su accesibilidad y desplegarlos en QA. \\
Entregable & Portal Web SITCO, Portal de Contratistas y Consola de Portería. \\
Criterio de aceptación & Requerimientos aprobados en QA y WCAG 2.2 nivel AA sin incumplimientos. El Portal de Contratistas exige doble factor. Cada empresa administra la identidad de su personal y consulta su acreditación sin llamar por teléfono. Resiste 1,5 veces su pico. \\
Responsable & Gabriel Sanzana (LD). \\
Hitos asociados & M4–M11. Aporta a H4 (mes 10). \\
Esfuerzo estimado & 1.024 HH (LD 80, DEV 816 y UXR ext. 128). PERT media. \\
Referencias & Subdocumento 4, 4.1.4.1, y Bases Técnicas, RT-16.30 (p. 26). T-12: RT-12.03 y 12.12; RT-13.02, RT-16.31, 16.32 y 16.34; RT-25.04, RF-ACR-004, 005, 013 y 015; RNF-SEG-003, RNF-USA-010 y 011. \\
\end{tabla}
\begin{tabla}{I{0.269}I{0.731}}{Diccionario del entregable 1.4.6, Datos migrados.}{tab:diccionario-entregable-1-4-6-datos}
  {\cab{1.4.6} & \cab{Datos migrados}}
  {Fuente: elaboración propia a partir de la Guía del PMBOK (PMI, 2017, p. 162).}
Descripción del trabajo & Datos históricos migrados, saneados y conciliados en las dos etapas, y un repositorio de consulta para los datos que no se migran. \\
Responsable & Benjamín Velásquez (LDA). \\
Hitos asociados & M3–M20. Aporta a H7 (mes 16), H10 (mes 18) y H12 (mes 21). \\
Esfuerzo estimado & 3.904 HH, todas propias (LDA 216, DEV 128, ADT 584 y DIG 2.976). \\
1.4.6.1 & Informe de perfilado y datos saneados \\
Descripción del trabajo & Revisar cada fuente de datos, decidir cómo tratar cada defecto y corregirlo. \\
Entregable & Informe de fuentes con decisión sobre cada defecto y datos limpios. \\
Criterio de aceptación & Informe cerrado antes del primer ensayo, con decisión registrada para cada defecto, recibido por la Contraparte Técnica. Muestra de al menos 200 registros por entidad crítica. Plan de migración aprobado antes del ensayo. \\
Responsable & Benjamín Velásquez (LDA). \\
Hitos asociados & M3–M8. \\
Esfuerzo estimado & 168 HH (LDA 56 y ADT 112). PERT alta. \\
Referencias & Subdocumento 5, 5.3.3 y RT-05.12. Bases Técnicas, RT-05.15 (p. 25); TRA-MIG-005. T-12: RT-05.11 y 05.12; TRA-MIG-001, 002 y 005. \\
1.4.6.2 & Datos de la Etapa 1 cargados y conciliados \\
Descripción del trabajo & Hacer dos ensayos de migración en Preproducción y la carga definitiva, y conciliar los datos. \\
Entregable & Datos de la Etapa 1 cargados, con acta de conciliación. \\
Criterio de aceptación & Acta firmada por la Contraparte Técnica. Completitud del 100 \% de los registros, sin vínculos rotos. Cero diferencias en sumas de control. Sin errores en la muestra de 200 registros por entidad crítica, con tasa de error bajo 1,5 \%. \\
Responsable & Benjamín Velásquez (LDA). \\
Hitos asociados & M9–M15. Aporta a H7 (mes 16). \\
Esfuerzo estimado & 240 HH (LDA 48, DEV 48 y ADT 144). PERT alta. \\
Referencias & Subdocumento 5, 5.3.2, 5.3.3 y 5.3.4; RT-05.13, 05.14, TRA-MIG-003. T-12: RT-05.11, 05.13 y 05.14; RT-11.25, TRA-MIG-001, 003 y 004. \\
1.4.6.3 & Formularios de seguridad digitalizados \\
Descripción del trabajo & Escanear y indexar formularios de seguridad en lotes semanales. Los ilegibles se registran como no migrables. \\
Entregable & Formularios de seguridad en papel escaneados, indexados y cargados. \\
Criterio de aceptación & Los 193.000 formularios quedan indexados en PDF/A antes del primer ensayo de la Etapa 2. Error de 2 \% o menos en muestra de 200 formularios por lote. Cada formulario indexado por tipo, fecha, área, empresa, persona y folio. \\
Responsable & Benjamín Velásquez (LDA). \\
Hitos asociados & M10–M15. Aporta a H10 (mes 18). \\
Esfuerzo estimado & 3.184 HH (LDA 56, ADT 152 y DIG 2.976). PERT baja. \\
Supuestos & Tres puestos de escaneo durante 124 días hábiles, con cuatro digitalizadores rotando. El primer lote es piloto, si rinde menos, se abre un cuarto puesto. \\
Referencias & Subdocumento 5, 5.3.1 (Tabla \refext{S5}{tab:volumen-migracion-origen}), 5.3.2 y 5.3.3; TRA-MIG-005. \\
1.4.6.4 & Datos de la Etapa 2 cargados y conciliados \\
Descripción del trabajo & Hacer dos ensayos en Preproducción y la carga final al cierre del mes 18, y conciliarlos. \\
Entregable & Datos de la Etapa 2 cargados, con acta de conciliación. \\
Criterio de aceptación & Ensayos conciliados y acta firmada por la Contraparte Técnica. Cumple los criterios de la Etapa 1. Campos personales sensibles enmascarados de forma irreversible. \\
Responsable & Benjamín Velásquez (LDA). \\
Hitos asociados & M16–M20. Aporta a H10 (mes 18) y H12 (mes 21). \\
Esfuerzo estimado & 192 HH (LDA 48, DEV 24 y ADT 120). PERT alta. \\
Referencias & Subdocumento 5, 5.2.6, 5.3.2 y 5.3.3; TRA-MIG-003. T-12: RT-05.11, 05.13 y 05.14; TRA-MIG-003 y 004. \\
1.4.6.5 & Repositorio de consulta histórica habilitado \\
Descripción del trabajo & Crear una copia de solo lectura por módulo con los registros anteriores a la migración. No recalcula períodos cerrados. \\
Entregable & Repositorio de consulta con los registros no migrados de ambas etapas. \\
Criterio de aceptación & Cada registro se consulta solo en lectura, con su origen y su limitación identificados. Mismo control de acceso y auditoría que el registro vivo. Probada con usuarios de cada área. \\
Responsable & Benjamín Velásquez (LDA). \\
Hitos asociados & M10–M20. Aporta a H7 (mes 16) y H12 (mes 21). \\
Esfuerzo estimado & 120 HH (LDA 8, DEV 56 y ADT 56). PERT media. \\
Referencias & Bases técnicas transversales, RT-05.15, Formulario T-12, RF-ADM-001, y Subdocumento 5, 5.3.5. T-12: RNF-POR-002; RNF-SEG-005, TRA-MIG-005. \\
\end{tabla}

\apartado{Rama 1.5 Sistema certificado}
La rama 1.5 tiene 6 entregables y 16 paquetes de trabajo, que suman 3.840 HH. Las tablas \ref{tab:diccionario-entregable-1-5-1-suite} a \ref{tab:diccionario-entregable-1-5-6-acta} presentan la ficha de cada entregable y, bajo ella, la de cada uno de sus paquetes.

\begin{tabla}{I{0.269}I{0.731}}{Diccionario del entregable 1.5.1, Suite de pruebas automatizadas e informes de integración.}{tab:diccionario-entregable-1-5-1-suite}
  {\cab{1.5.1} & \cab{Suite de pruebas automatizadas e informes de integración}}
  {Fuente: elaboración propia a partir de la Guía del PMBOK (PMI, 2017, p. 162).}
Descripción del trabajo & Pruebas automáticas que corren con cada cambio, e informes de integración y de análisis. Los ambientes de desarrollo y QA usan datos sintéticos. \\
Responsable & Kavon Kermani (LCP). \\
Hitos asociados & M3–M18. Aporta a H4 (mes 10), H9 (mes 17) y H10 (mes 18). \\
Esfuerzo estimado & 2.048 HH, todas propias (LCP 512, DEV 384 y QA 1.152). \\
1.5.1.1 & Pruebas unitarias y de contrato automatizadas \\
Descripción del trabajo & Escribir pruebas con cada historia, mantener los contratos con sus consumidores y medir la cobertura. \\
Entregable & Pruebas unitarias y de contrato que corren en cada cambio. \\
Criterio de aceptación & Reporte con los contratos verificados. Cobertura de 70 \% o más en la lógica de negocio y 80 \% o más en código nuevo. El cálculo crítico: 90 \% de líneas y 85 \% de ramas. \\
Responsable & Kavon Kermani (LCP). \\
Hitos asociados & M3–M18. Aporta a H4 (mes 10), H9 (mes 17) y H10 (mes 18). \\
Esfuerzo estimado & 768 HH (LCP 128, DEV 384 y QA 256). PERT media. \\
Referencias & RT-04.11. T-12: RT-04.05 y 04.11; RT-21.20. \\
1.5.1.2 & Pruebas de integración y regresión automatizadas \\
Descripción del trabajo & Automatizar los flujos completos, repetir las pruebas de regresión y restaurar datos antes de cada ciclo. \\
Entregable & Batería de integración y de regresión que se ejecuta en cada promoción. \\
Criterio de aceptación & Regresión sin fallas en cada paso a Preproducción. La regresión de la Etapa 1 dura 60 minutos o menos. Una prueba confirma que todo intento de escribir en el historiador se rechaza y se registra. \\
Responsable & Kavon Kermani (LCP). \\
Hitos asociados & M3–M18. Aporta a H4 (mes 10), H9 (mes 17) y H10 (mes 18). \\
Esfuerzo estimado & 1.024 HH (LCP 256 y QA 768). PERT media. \\
Referencias & Subdocumento 9, 9.2.4. T-12: RT-09.10; RT-21.20. \\
1.5.1.3 & Informes de análisis estático y dinámico \\
Descripción del trabajo & Configurar los análisis automáticos, revisar sus hallazgos y publicar los informes. \\
Entregable & Informes de análisis de seguridad del código, claves, dependencias, imágenes y ejecución. \\
Criterio de aceptación & Cada entrega a pruebas sin hallazgos nuevos críticos ni altos. Pantallas conformes con WCAG 2.2 nivel AA. Hallazgos medios corregidos en 30 días corridos. \\
Responsable & Kavon Kermani (LCP). \\
Hitos asociados & M3–M18. Aporta a H4 (mes 10), H9 (mes 17) y H10 (mes 18). \\
Esfuerzo estimado & 256 HH (LCP 128 y QA 128). PERT baja. \\
Referencias & RT-11.22 y Subdocumento 9, 9.2.3. T-12: RT-04.05; RT-13.01 y 13.11; RT-21.20, RNF-USA-010 y 015; RNF-TRV-11.04. \\
\end{tabla}
\begin{tabla}{I{0.269}I{0.731}}{Diccionario del entregable 1.5.2, Informes de usabilidad en terreno.}{tab:diccionario-entregable-1-5-2-informes}
  {\cab{1.5.2} & \cab{Informes de usabilidad en terreno}}
  {Fuente: elaboración propia a partir de la Guía del PMBOK (PMI, 2017, p. 162).}
Descripción del trabajo & Pruebas de usabilidad con usuarios reales en las condiciones de la faena, de M5 a M20. Supone acceso a usuarios en terreno. \\
Responsable & Kavon Kermani (LCP). \\
Hitos asociados & M5–M20. Aporta a H5 (mes 12), H7 (mes 16), H10 (mes 18) y H12 (mes 21). \\
Esfuerzo estimado & 336 HH, 112 externas (LCP 112, QA 112 y UXR ext. 112). \\
1.5.2.1 & Informe de usabilidad de la Etapa 1 \\
Descripción del trabajo & Preparar escenarios, probar con usuarios, registrar y corregir los hallazgos. \\
Entregable & Informe de pruebas en terreno de los canales de la Etapa 1. \\
Criterio de aceptación & Informe parcial en el mes 11 y final en el mes 15. Al menos 90 \% de los participantes completa sin ayuda las tareas críticas. Se prueba con guantes, con sol directo y entre −6 °C y 34 °C. \\
Responsable & Kavon Kermani (LCP). \\
Hitos asociados & M5–M15. Aporta a H5 (mes 12) y H7 (mes 16). \\
Esfuerzo estimado & 192 HH (LCP 64, QA 64 y UXR ext. 64). PERT media. \\
Referencias & RT-13.04 y Bases Técnicas, RT-13.08 (p. 25). T-12: RT-13.01, 13.04 y 13.08; RNF-USA-001 a 003. \\
1.5.2.2 & Informe de usabilidad de la Etapa 2 \\
Descripción del trabajo & Preparar escenarios, probar con usuarios, registrar y corregir los hallazgos. \\
Entregable & Informe de pruebas con usuarios de los canales de la Etapa 2. \\
Criterio de aceptación & Informe parcial en el mes 17 y final en el mes 20. Al menos 90 \% de los participantes completa sin ayuda las tareas críticas. Un registro de inspección sin conexión se completa en 3 minutos o menos. \\
Responsable & Kavon Kermani (LCP). \\
Hitos asociados & M15–M20. Aporta a H10 (mes 18) y H12 (mes 21). \\
Esfuerzo estimado & 144 HH (LCP 48, QA 48 y UXR ext. 48). PERT media. \\
Referencias & RT-13.04, 13.08. T-12: RNF-USA-007. \\
\end{tabla}
\begin{tabla}{I{0.269}I{0.731}}{Diccionario del entregable 1.5.3, Acta de certificación de la Etapa 1.}{tab:diccionario-entregable-1-5-3-acta}
  {\cab{1.5.3} & \cab{Acta de certificación de la Etapa 1}}
  {Fuente: elaboración propia a partir de la Guía del PMBOK (PMI, 2017, p. 162).}
Descripción del trabajo & Acta que certifica la solución de la Etapa 1 con las pruebas de carga, de resiliencia y de aceptación. Supone que Preproducción es equivalente a producción. \\
Responsable & Kavon Kermani (LCP). \\
Hitos asociados & M11–M12. Aporta a H5 (mes 12). \\
Esfuerzo estimado & 440 HH, 72 externas (LF 32, LCP 112, SRE 80, QA 144 y PRF ext. 72). \\
1.5.3.1 & Informe de pruebas de carga de la Etapa 1 \\
Descripción del trabajo & Preparar la prueba, ejecutarla en Preproducción y analizar resultados, con cambio de turno y detención mayor. \\
Entregable & Informe de carga, estrés y resistencia de la Etapa 1. \\
Criterio de aceptación & Prueba aprobada si se sostienen 300 transacciones por segundo durante 60 minutos, con punto de quiebre identificado. Sin degradación sobre 10 \% en 72 horas. En el 95 \% de las respuestas, 2 segundos o menos en terreno. \\
Responsable & Kavon Kermani (LCP). \\
Hitos asociados & M12. Aporta a H5 (mes 12). \\
Esfuerzo estimado & 120 HH (LCP 48 y PRF externo 72). PERT media. \\
Referencias & RT-09.06, 09.07, Subdocumento 9, 9.2.4.5, Subdocumento 5, 5.4 y las Bases técnicas transversales (CMAA, 2026c, 9.1. T-12: RT-05.05); RT-09.02, 09.03, 09.06 a 09.08 y 09.10; RNF-CAP-004, 007 y 008; RNF-DES-017, RNF-DIS-009. \\
1.5.3.2 & Informe de pruebas de resiliencia de la Etapa 1 \\
Descripción del trabajo & Provocar cinco fallas en Preproducción con Azure Chaos Studio y registrar cómo se recupera. \\
Entregable & Informe de pruebas de fallas: instancia, zona, dependencia, latencia y disco. \\
Criterio de aceptación & La Contraparte Técnica recibe el informe si las cinco fallas se recuperan sin intervención ni pérdida de transacciones o sesión. Cada prueba registra hipótesis y resultado. \\
Responsable & Lucas Zamora (SRE). \\
Hitos asociados & M12. Aporta a H5 (mes 12). \\
Esfuerzo estimado & 80 HH del SRE. PERT media. \\
Referencias & RT-10.07. T-12: RT-02.05 y 02.09; RT-03.02, RT-10.07 y 10.08. \\
1.5.3.3 & Informe de pruebas de aceptación de la Etapa 1 \\
Descripción del trabajo & Ejecutar los casos con los usuarios, incluido el Portal de Contratistas, registrar defectos y obtener la firma. \\
Entregable & Informe de pruebas de aceptación de la Etapa 1, según T-17. \\
Criterio de aceptación & Casos de aceptación aprobados y firmados por la Contraparte Técnica. Sin defectos críticos ni altos abiertos. \\
Responsable & Kavon Kermani (LCP). \\
Hitos asociados & M11–M12. Aporta a H5 (mes 12). \\
Esfuerzo estimado & 240 HH (LF 32, LCP 64 y QA 144). PERT media. \\
Referencias & Art. 18.1 y 18.2, y Formulario T-17. \\
\end{tabla}
\begin{tabla}{I{0.269}I{0.731}}{Diccionario del entregable 1.5.4, Informes de pruebas de seguridad ofensiva.}{tab:diccionario-entregable-1-5-4-informes}
  {\cab{1.5.4} & \cab{Informes de pruebas de seguridad ofensiva}}
  {Fuente: elaboración propia a partir de la Guía del PMBOK (PMI, 2017, p. 162).}
Descripción del trabajo & Pruebas de intrusión hechas por un tercero independiente, simulacros de incidentes y control de vulnerabilidades. Supone que el auditor independiente se contrata por campaña. \\
Responsable & David Martínez (ESI). \\
Hitos asociados & M3–M20, campañas en M11, M15, M17 y M20. Aporta a H5 (mes 12), H7 (mes 16), H10 (mes 18) y H12 (mes 21). \\
Esfuerzo estimado & 512 HH, 256 externas (ESI 176, ASG 80 y PEN ext. 256). \\
1.5.4.1 & Informes de pruebas de intrusión independientes \\
Descripción del trabajo & Definir el alcance, ejecutar la campaña, corregir y verificar las correcciones. \\
Entregable & Informes de cuatro campañas de intrusión independientes. \\
Criterio de aceptación & Cada campaña se recibe sin hallazgos críticos ni altos abiertos, con informe íntegro y plan de corrección con plazos. Aplicaciones probadas según OWASP ASVS 4.0.3, nivel 2. \\
Responsable & David Martínez (ESI). \\
Hitos asociados & M11, M15, M17 y M20. Aporta a H5 (mes 12), H7 (mes 16), H10 (mes 18) y H12 (mes 21). \\
Esfuerzo estimado & 336 HH (ESI 80 y PEN ext. 256). PERT media. \\
Supuestos & El auditor independiente se contrata por campaña. \\
Referencias & RT-11.20; OWASP ASVS 4.0.3 (OWASP Foundation, 2021); IEC 62443. \\
1.5.4.2 & Informe de simulación de incidente \\
Descripción del trabajo & Preparar el escenario, ejecutar el ejercicio y registrar las mejoras. \\
Entregable & Ejercicio de respuesta a un incidente de seguridad realizado con el cliente. \\
Criterio de aceptación & Ejercicio con el cliente y mejoras al plan 1.3.3.4. El cliente recibe aviso en dos horas desde la detección de un incidente crítico. \\
Responsable & David Martínez (ESI). \\
Hitos asociados & M15, fuera de la detención de marzo de 2028. Aporta a H7 (mes 16). \\
Esfuerzo estimado & 32 HH (ESI 24 y ASG 8). PERT media. \\
Referencias & RT-11.18 y RT-11.21. \\
1.5.4.3 & Informe de gestión de vulnerabilidades \\
Descripción del trabajo & Escanear dependencias con Dependabot e imágenes con Trivy en cada versión, clasificar hallazgos y corregirlos dentro de su plazo. \\
Entregable & Escaneo continuo de vulnerabilidades y seguimiento de su corrección. \\
Criterio de aceptación & Registro con fecha de detección y de cierre de cada vulnerabilidad. Críticas corregidas en 7 días, altas en 15 y medias en 30 días corridos. \\
Responsable & David Martínez (ESI). \\
Hitos asociados & M3–M20. \\
Esfuerzo estimado & 144 HH (ESI 72 y ASG 72). PERT baja. \\
Referencias & RT-11.04. \\
\end{tabla}
\begin{tabla}{I{0.269}I{0.731}}{Diccionario del entregable 1.5.5, Informe de recuperación y corte de 48 horas.}{tab:diccionario-entregable-1-5-5-informe}
  {\cab{1.5.5} & \cab{Informe de recuperación y corte de 48 horas}}
  {Fuente: elaboración propia a partir de la Guía del PMBOK (PMI, 2017, p. 162).}
Descripción del trabajo & Pruebas reales de recuperación ante desastres y de un corte de enlace de 48 horas en el Nodo de Faena. Supone que la región secundaria Brazil South está habilitada. \\
Responsable & Lucas Zamora (SRE). \\
Hitos asociados & M15 y M21. Aporta a H7 (mes 16) y H12 (mes 21). \\
Esfuerzo estimado & 240 HH, todas propias (SRE 112, OPS 96 y TEC 32). \\
1.5.5.1 & Informe de conmutación y retorno \\
Descripción del trabajo & Cambiar a la región secundaria y volver, según el procedimiento 1.3.1.4, midiendo tiempo y pérdida de datos. \\
Entregable & Informe del cambio real a la región secundaria y del retorno. \\
Criterio de aceptación & Informe con recuperación en 4 horas o menos y pérdida de datos de 15 minutos o menos. Retorno al sitio principal y datos reconciliados. Se hace dentro de la detención de planta y antes del paso a producción. \\
Responsable & Lucas Zamora (SRE). \\
Hitos asociados & M15 y M21, dentro de las detenciones de marzo y septiembre de 2028. Aporta a H7 (mes 16) y H12 (mes 21). \\
Esfuerzo estimado & 80 HH (SRE 48 y OPS 32). PERT media. \\
Referencias & Subdocumento 4, 4.3.2.2, RT-07.05 a 07.07 y RT-20.02. T-12: RT-07.04 a 07.07; RNF-OPE-005. \\
1.5.5.2 & Informe del corte de 48 horas \\
Descripción del trabajo & Cortar el enlace durante un cambio de turno, operar con la punta local y sincronizar fuera de la tronadura. \\
Entregable & Informe del Nodo de Faena en un corte de 48 horas. \\
Criterio de aceptación & Cero registros perdidos y sincronización en 4 horas o menos, sin intervención manual. Tras reiniciar un servidor, el registro en terreno vuelve a operar en 1 hora o menos, sin perder eventos. \\
Responsable & Lucas Zamora (SRE). \\
Hitos asociados & M15 y M21. Aporta a H7 (mes 16) y H12 (mes 21). \\
Esfuerzo estimado & 160 HH (SRE 64, OPS 64 y TEC 32). PERT media. \\
Referencias & Subdocumento 4, 4.1.3.4 y 4.1.9; Subdocumento 5, Tabla \refext{S5}{tab:objetivos-recuperacion-mecanismo-datos-sostiene}; Bases Técnicas, RT-03.10 y Capítulo B. T-12: RNF-DES-012; RNF-DIS-001, 003 y 006. \\
\end{tabla}
\begin{tabla}{I{0.269}I{0.731}}{Diccionario del entregable 1.5.6, Acta de certificación de la Etapa 2.}{tab:diccionario-entregable-1-5-6-acta}
  {\cab{1.5.6} & \cab{Acta de certificación de la Etapa 2}}
  {Fuente: elaboración propia a partir de la Guía del PMBOK (PMI, 2017, p. 162).}
Descripción del trabajo & Certificación de la solución de la Etapa 2. El desarrollo queda cerrado en el mes 18, inclusive. \\
Responsable & Kavon Kermani (LCP). \\
Hitos asociados & M17–M18. Aporta a H10 (mes 18). \\
Esfuerzo estimado & 264 HH, 72 externas (LCP 64, LI 16, SRE 80, QA 32 y PRF ext. 72). \\
1.5.6.1 & Informe de pruebas de carga de la Etapa 2 \\
Descripción del trabajo & Ajustar el perfil de carga a la Etapa 2, ejecutar la prueba y analizar resultados, con cambio de turno. \\
Entregable & Informe de pruebas de carga, estrés y resistencia con ambas etapas. \\
Criterio de aceptación & Con ambas etapas, se sostienen 300 transacciones por segundo. Sin degradación sobre 10 \% en 72 horas, con punto de quiebre identificado. Consulta de habilitación en 1,5 segundos o menos en el 95 \% de las respuestas. \\
Responsable & Kavon Kermani (LCP). \\
Hitos asociados & M18. Aporta a H10 (mes 18). \\
Esfuerzo estimado & 120 HH (LCP 48 y PRF externo 72). PERT media. \\
Referencias & RT-09.06, 09.07 y Subdocumento 9, 9.2.4.5. T-12: RT-09.02, 09.06 y 09.07. \\
1.5.6.2 & Informe de pruebas de resiliencia de la Etapa 2 \\
Descripción del trabajo & Repetir en Preproducción las cinco fallas de 1.5.3.2 en conectores y módulos de la Etapa 2, y registrar la recuperación. \\
Entregable & Informe de fallas controladas con la Etapa 2. \\
Criterio de aceptación & Sin pérdida de transacciones y recuperación sin intervención. Incluye los conectores de la Etapa 2. \\
Responsable & Lucas Zamora (SRE). \\
Hitos asociados & M18. Aporta a H10 (mes 18). \\
Esfuerzo estimado & 80 HH del SRE. PERT media. \\
Referencias & RT-10.07. \\
1.5.6.3 & Informe de regresión de ambas etapas \\
Descripción del trabajo & Ejecutar la regresión completa de ambas etapas, corregir defectos y repetirla hasta cerrar críticos y altos. \\
Entregable & Informe de regresión de ambas etapas y aceptación de la Etapa 2. \\
Criterio de aceptación & Regresión sin fallas, casos de la Etapa 2 firmados y sin defectos críticos ni altos abiertos. \\
Responsable & Kavon Kermani (LCP). \\
Hitos asociados & M17–M18. Aporta a H9 (mes 17) y H10 (mes 18). \\
Esfuerzo estimado & 64 HH (LCP 16, LI 16 y QA 32). PERT media. \\
Referencias & Art. 17.2.3. \\
\end{tabla}

\apartado{Rama 1.6 Solución implantada}
La rama 1.6 tiene 6 entregables y 22 paquetes de trabajo, que suman 21.920 HH. Las tablas \ref{tab:diccionario-entregable-1-6-1-infraestructura} a \ref{tab:diccionario-entregable-1-6-6-plan} presentan la ficha de cada entregable y, bajo ella, la de cada uno de sus paquetes.

\begin{tabla}{I{0.269}I{0.731}}{Diccionario del entregable 1.6.1, Infraestructura y observabilidad habilitadas.}{tab:diccionario-entregable-1-6-1-infraestructura}
  {\cab{1.6.1} & \cab{Infraestructura y observabilidad habilitadas}}
  {Fuente: elaboración propia a partir de la Guía del PMBOK (PMI, 2017, p. 162).}
Descripción del trabajo & Infraestructura híbrida en la nube y en faena, con seguimiento unificado de los sistemas. Supone que la obra civil de la sala la hace el cliente. \\
Responsable & Benjamín Rodríguez (AS). \\
Hitos asociados & M1–M16. Aporta a H3 (mes 6). \\
Esfuerzo estimado & 616 HH, 192 externas (AS 176, ESI 64, SRE 8, DEV 152, TEC 24 y MON ext. 192). \\
1.6.1.1 & Nube con cinco ambientes habilitados \\
Descripción del trabajo & Crear cinco ambientes en la nube con caché, malla de servicios, identidad y observabilidad. No incluye la puerta operacional. \\
Entregable & Cinco ambientes: Desarrollo, QA, Preproducción, Producción y Recuperación. \\
Criterio de aceptación & Cinco ambientes operativos, con Preproducción equivalente a Producción, y ningún ambiente no productivo conectado a Producción. La consola única diagnostica una falla de integración en 30 minutos o menos. Personal, contratistas y compradores entran con doble factor. \\
Responsable & Benjamín Rodríguez (AS). \\
Hitos asociados & M2–M6. Aporta a H3 (mes 6). \\
Esfuerzo estimado & 248 HH (AS 80, ESI 16 y DEV 152). PERT media. \\
Referencias & RT-04.01, RT-03.03 y Subdocumento 4, 4.1; OpenTelemetry Authors (2024). T-12: RT-02.05, 02.06 y 02.08 a 02.10; RT-03.01 a 03.04, 03.06 y 03.16; RT-04.01, 04.02, 04.08 y 04.13; RT-05.18 y 05.24; RT-09.04 y 09.08; RT-11.07, 11.08, 11.11, 11.12 y 11.14; RT-14.01, 14.07 y 14.08; RT-15.02, RT-16.05, RT-22.09, RF-ADM-005, RNF-DIS-007, RNF-SEG-008, 010, 011 y 015; RNF-OPE-001 y 006; RNF-TRV-02.08, RNF-TRV-11.07 y 11.11. \\
1.6.1.2 & Sala técnica, Nodo de Faena y gabinetes instalados \\
Descripción del trabajo & Recibir la sala, montar el Nodo de Faena y los gabinetes, probar en sitio y entregar el plano de zonas. \\
Entregable & Sala técnica, Nodo de Faena de tres servidores y cuatro gabinetes instalados. \\
Criterio de aceptación & Pruebas de climatización, energía y extinción de incendios aprobadas en sitio y nodo operativo. El respaldo de energía da 30 minutos o más a plena carga y el generador 24 horas continuas. La instalación cumple la NCh Elec. 2777. \\
Responsable & Benjamín Rodríguez (AS). \\
Hitos asociados & M1–M6, M10 y M16. Aporta a H3 (mes 6). \\
Esfuerzo estimado & 224 HH (AS 48, SRE 8, TEC 24 y MON ext. 144). PERT media. \\
Referencias & Subdocumento 4, 4.3.1, Bases Técnicas, RT-06.01 (p. 25), y las Bases técnicas transversales (CMAA, 2026c, RT-06.10. T-12: RT-03.14 y 03.15); RT-06.01 a 06.17, 06.19, 06.22, 06.25 a 06.31 y 06.33; RT-08.02 y 08.04; RNF-INF-001. \\
1.6.1.3 & Segunda fibra, enlace satelital y diodo operativos \\
Descripción del trabajo & Coordinar la segunda fibra, instalar el terminal satelital y habilitar el diodo de datos y el equipo SD-WAN. \\
Entregable & Segunda fibra, enlace satelital, diodo de datos y servidor espejo del historiador. \\
Criterio de aceptación & Acta firmada por la Contraparte Técnica si la conmutación entre fibras ocurre en 5 segundos o menos. Respaldo satelital probado en M6. El diodo se recibe en M5 y pasa datos en un solo sentido. \\
Responsable & Benjamín Rodríguez (AS). \\
Hitos asociados & M2–M9. Aporta a H3 (mes 6). \\
Esfuerzo estimado & 144 HH (AS 48, ESI 48 y MON externo 48). PERT alta. \\
Referencias & Subdocumento 3, 3.3.4 y RT-03.17. T-12: RT-06.32 y 06.33; RNF-SEG-004. \\
\end{tabla}
\begin{tabla}{I{0.269}I{0.731}}{Diccionario del entregable 1.6.2, Dispositivos instalados.}{tab:diccionario-entregable-1-6-2-dispositivos}
  {\cab{1.6.2} & \cab{Dispositivos instalados}}
  {Fuente: elaboración propia a partir de la Guía del PMBOK (PMI, 2017, p. 162).}
Descripción del trabajo & Dispositivos de terreno de ambas etapas y estaciones de administración instalados, configurados y en gestión. Supone que el cliente compra los equipos de terreno según la especificación de Overfore. \\
Responsable & Lucas Zamora (SRE). \\
Hitos asociados & M6–M19. Aporta a H6 (mes 13) y H11 (mes 19). \\
Esfuerzo estimado & 656 HH, todas propias (AS 40, SRE 128 y TEC 488). \\
1.6.2.1 & Botoneras y tabletas de la Etapa 1 instaladas \\
Descripción del trabajo & Especificar, recibir, instalar y probar las botoneras y tabletas de la Etapa 1, con una prueba previa con guantes. \\
Entregable & Seis botoneras, 18 tabletas y dos equipos de portería instalados. \\
Criterio de aceptación & Botoneras y tabletas operativas antes del mes 13, con grado de protección IP65 o IP67 y batería probada en invierno y verano. Bloqueo o borrado remoto en 15 minutos o menos. \\
Responsable & Lucas Zamora (SRE). \\
Hitos asociados & M6 y M8–M12. Aporta a H6 (mes 13). \\
Esfuerzo estimado & 336 HH (AS 16, SRE 64 y TEC 256). PERT media. \\
Supuestos & Si la validación de prototipos exige otro hardware, se entrega una adenda antes de la compra. \\
Referencias & Subdocumento 3, 3.3.3, Subdocumento 4, 4.2.1.3 y Formulario T-11. T-12: RT-03.18; RT-08.10 a 08.12, 08.14 y 08.15; RT-13.08, RF-TER-002, RNF-SEG-016, RNF-USA-003. \\
1.6.2.2 & Lectores, caudalímetros y tabletas de la Etapa 2 \\
Descripción del trabajo & Especificar, recibir, instalar y probar los lectores faciales, caudalímetros y tabletas de la Etapa 2. \\
Entregable & Siete lectores faciales, cuatro caudalímetros y 15 tabletas, instalados. \\
Criterio de aceptación & La Contraparte Técnica recibe siete lectores faciales operativos antes del mes 19. Tres caudalímetros más uno de repuesto y 15 tabletas, con inventario por número de serie. \\
Responsable & Lucas Zamora (SRE). \\
Hitos asociados & M12–M19. Aporta a H11 (mes 19). \\
Esfuerzo estimado & 256 HH (AS 16, SRE 48 y TEC 192). PERT media. \\
Referencias & Subdocumento 3, 3.3.3. T-12: RT-08.10 a 08.12. \\
1.6.2.3 & Estaciones de administración instaladas \\
Descripción del trabajo & Especificar, recibir, instalar, proteger y registrar en la gestión central las estaciones de administración. \\
Entregable & Siete estaciones de administración con dos monitores, instaladas. \\
Criterio de aceptación & Siete estaciones gestionadas en Microsoft Intune antes del mes 13, con informe de cumplimiento aprobado. Inventario con número de serie, usuario y ubicación. \\
Responsable & Lucas Zamora (SRE). \\
Hitos asociados & M11–M12. Aporta a H6 (mes 13). \\
Esfuerzo estimado & 64 HH (AS 8, SRE 16 y TEC 40). PERT media. \\
Referencias & RT-08.07 a 08.09 y Formulario T-11. \\
\end{tabla}
\begin{tabla}{I{0.269}I{0.731}}{Diccionario del entregable 1.6.3, Actas de marcha blanca.}{tab:diccionario-entregable-1-6-3-actas}
  {\cab{1.6.3} & \cab{Actas de marcha blanca}}
  {Fuente: elaboración propia a partir de la Guía del PMBOK (PMI, 2017, p. 162).}
Descripción del trabajo & Cierre de cada marcha blanca, con medición diaria y conciliación contra el registro vigente. Supone que en la Etapa 2 las planillas rigen sólo para los procesos que entran con ella. \\
Responsable & Benjamín Velásquez (LIG). \\
Hitos asociados & M13–M20. Aporta a H6 (mes 13), H7 (mes 16), H11 (mes 19) y H12 (mes 21). \\
Esfuerzo estimado & 12.000 HH, todas propias (LF 212, LI 72, LIG 388, ADT 540, OPS 5.448 y TEC 5.340). \\
1.6.3.1 & Acta de marcha blanca de la Etapa 1 \\
Descripción del trabajo & Probar la vuelta a las planillas, medir a diario, conciliar y cerrar la marcha blanca de la Etapa 1. \\
Entregable & Acta firmada de la marcha blanca de la Etapa 1. \\
Criterio de aceptación & Acta firmada del mes 16 que acredita, en las cuatro últimas semanas, las seis condiciones del Art. 17.3. Incluye 95 \% de los registros en SITCO, 99,9 \% de disponibilidad y como máximo 4 defectos medios y 13 bajos abiertos. \\
Responsable & Benjamín Velásquez (LIG). \\
Hitos asociados & M13–M15. Aporta a H6 (mes 13) y H7 (mes 16). \\
Esfuerzo estimado & 520 HH (LIG 156 y ADT 364). Sin PERT. \\
Supuestos & Las áreas entran escalonadas, cada una tras cuatro semanas de capacitación; la marcha blanca dura los tres meses del Art. 17.1. \\
Referencias & Art. 17.3, 78.2 y 79.4, y Bases Técnicas, RT-09.01. T-12: RT-20.03 y 20.04; RT-21.04, RT-22.05, RNF-DES-001, RNF-DIS-002, RNF-USA-005, TRA-001, 004 y 008. \\
1.6.3.2 & Acta de marcha blanca de la Etapa 2 \\
Descripción del trabajo & Probar la vuelta atrás, conciliar a diario ambas etapas y cerrar la marcha blanca de la Etapa 2. \\
Entregable & Acta firmada de la marcha blanca de la Etapa 2. \\
Criterio de aceptación & Acta firmada del mes 21 que acredita, en las cuatro últimas semanas, las seis condiciones del Art. 17.3. Incluye 95 \% de los registros en SITCO, 99,9 \% de disponibilidad, como máximo 1 defecto medio y 5 bajos abiertos. \\
Responsable & Benjamín Velásquez (LIG). \\
Hitos asociados & M19–M20. Aporta a H11 (mes 19) y H12 (mes 21). \\
Esfuerzo estimado & 360 HH (LI 72, LIG 144 y ADT 144). Sin PERT. \\
Supuestos & Las áreas entran escalonadas: seguridad, mantenimiento y cumplimiento ambiental, cada una tras cuatro semanas de capacitación. \\
Referencias & Art. 17.3 y Subdocumento 3, 3.2.2.5; Subdocumento 9, 9.2.6, Formulario T-17. T-12: RT-20.03 y 20.04. \\
1.6.3.3 & Acompañamiento en faena de la Etapa 1 \\
Descripción del trabajo & Programar y ejecutar el acompañamiento en puesto por área, de día y de noche, y registrar las consultas. \\
Entregable & Acompañamiento presencial en faena durante la marcha blanca de la Etapa 1. \\
Criterio de aceptación & Una persona por turno en cada área en marcha blanca. Una en las semanas 1 y 2, dos en las semanas 3 a 5 y cuatro en las semanas 6 a 13. Registro diario de turnos cubiertos. \\
Responsable & Benjamín Velásquez (LIG). \\
Hitos asociados & M13–M15. Aporta a H6 (mes 13) y H7 (mes 16). \\
Esfuerzo estimado & 6.824 HH (LF 52, LIG 52, OPS 3.360 y TEC 3.360). Sin PERT. \\
Referencias & RT-20.05, Art. 90.2, Bases Técnicas, 13.3 (p. 22), y Formulario T-18. T-12: RT-22.02. \\
1.6.3.4 & Acompañamiento en faena de la Etapa 2 \\
Descripción del trabajo & Programar el acompañamiento por área y acompañar en su puesto a seguridad, mantenimiento y cumplimiento ambiental. Registrar consultas. \\
Entregable & Acompañamiento presencial en faena durante la marcha blanca de la Etapa 2. \\
Criterio de aceptación & Una persona por turno, de día y de noche, en cada área en marcha blanca: una la primera semana, dos la segunda y tres desde la tercera. Registro diario de turnos y consultas. \\
Responsable & Benjamín Velásquez (LIG). \\
Hitos asociados & M19–M20. Aporta a H11 (mes 19) y H12 (mes 21). \\
Esfuerzo estimado & 4.248 HH (LF 144, LIG 36, OPS 2.088 y TEC 1.980). Sin PERT. \\
Referencias & RT-20.05, Art. 90.2 y Bases Técnicas, 13.3 (p. 22). \\
1.6.3.5 & Piloto de la Botonera de Pala \\
Descripción del trabajo & Usar la Botonera en una pala durante dos semanas y medir el tiempo de cada declaración. \\
Entregable & Informe del piloto y acta de aprobación de la gerencia mina. \\
Criterio de aceptación & Catorce días continuos en una pala, con ambos grupos del 7x7 y ambos turnos. El tiempo por declaración no supera el que justifica el diseño. La gerencia de operaciones mina firma el acta antes de extender la ola. \\
Responsable & Patricio Figueroa (LF). \\
Hitos asociados & M13. Aporta a H7 (mes 16). \\
Esfuerzo estimado & 48 HH (LF 16 y ADT 32). PERT media. \\
Supuestos & El acompañamiento en faena del piloto se cuenta en 1.6.3.3, no en este paquete. \\
Referencias & Subdocumento 3, 3.2.2.5 y Tabla \refext{S3}{tab:criterios-aceptacion-etapa-1}; Bases Técnicas, cap. 18, criterio 14, y cap. 10, restricción 8. T-12: RNF-DES-016; TRA-004. \\
\end{tabla}
\begin{tabla}{I{0.269}I{0.731}}{Diccionario del entregable 1.6.4, Actas de paso a producción.}{tab:diccionario-entregable-1-6-4-actas}
  {\cab{1.6.4} & \cab{Actas de paso a producción}}
  {Fuente: elaboración propia a partir de la Guía del PMBOK (PMI, 2017, p. 162).}
Descripción del trabajo & Actas del paso a producción de cada etapa, con despliegue progresivo y reversión probada. Supone que cada paso se hace de madrugada, fuera de embarques y detenciones. \\
Responsable & Lucas Zamora (SRE). \\
Hitos asociados & M16–M23. Aporta a H7 (mes 16) y H12 (mes 21). \\
Esfuerzo estimado & 6.096 HH, todas propias (SRE 216, OPS 2.940 y TEC 2.940). \\
1.6.4.1 & Acta de paso a producción y estabilización de la Etapa 1 \\
Descripción del trabajo & Obtener la aprobación de los comités y desplegar la Etapa 1 en tres olas: portería, mina y planta. \\
Entregable & Acta firmada y informe de estabilización de la Etapa 1. \\
Criterio de aceptación & Acta firmada por la Contraparte Técnica. Cada ola pasa tras 72 horas sin incidentes críticos ni altos, con 95 \% de registros en SITCO. Un incidente crítico sin resolver en 4 horas obliga a revertir en 30 minutos o menos. \\
Responsable & Lucas Zamora (SRE). \\
Hitos asociados & M16–M17. Aporta a H7 (mes 16). \\
Esfuerzo estimado & 2.064 HH (SRE 72, OPS 996 y TEC 996). Sin PERT. \\
Supuestos & Una persona por turno en cada área durante sus dos primeras semanas de producción, y una para toda la faena después, hasta el 31 de mayo de 2028. \\
Referencias & Art. 17.1, RT-20.06, Bases Técnicas, 13.3 (p. 22), y Formulario T-18. T-12: RT-20.01, 20.02, 20.05 y 20.06; RNF-DIS-004, TRA-006. \\
1.6.4.2 & Acta de paso a producción y estabilización de la Etapa 2 \\
Descripción del trabajo & Obtener la aprobación de los comités y desplegar la Etapa 2 en tres olas: seguridad, mantenimiento y cumplimiento ambiental. \\
Entregable & Acta de paso a producción y aceptación final de la Etapa 2. \\
Criterio de aceptación & Acta firmada por la Contraparte Técnica, sin degradar la Etapa 1. Cada ola pasa tras 72 horas sin incidentes críticos ni altos, con conciliación diaria sin diferencias no explicadas. Aceptación final exige la garantía constituida (Art. 37.1). \\
Responsable & Lucas Zamora (SRE). \\
Hitos asociados & M21–M23. Aporta a H12 (mes 21). \\
Esfuerzo estimado & 1.752 HH (SRE 72, OPS 840 y TEC 840). Sin PERT. \\
Supuestos & Una persona por turno en cada área durante sus dos primeras semanas de producción, hasta el 8 de noviembre de 2028. La verificación facial pasa con la primera ola si el acuerdo sindical está firmado. \\
Referencias & Art. 17.1, 17.2, RT-20.06, Bases Técnicas, 13.3 (p. 22), y Formulario T-18. T-12: RT-20.02 y 20.06; RT-22.07, RNF-DIS-004, TRA-006. \\
1.6.4.3 & Soporte de la Etapa 1 entre su estabilización y la operación \\
Descripción del trabajo & Atender incidentes y consultas de la Etapa 1 hasta el 31 de agosto de 2028. No incluye requerimientos nuevos. \\
Entregable & Informe mensual de atención de la Etapa 1, meses 18 a 20. \\
Criterio de aceptación & Una persona por turno, de día y de noche. Incidentes críticos atendidos en 15 minutos y resueltos en 4 horas. Ningún incidente crítico ni alto abierto al inicio de la operación. Una prueba de restauración por mes. \\
Responsable & Lucas Zamora (SRE). \\
Hitos asociados & M18–M20. Aporta a H12 (mes 21). \\
Esfuerzo estimado & 2.280 HH (SRE 72, OPS 1.104 y TEC 1.104). Sin PERT. \\
Supuestos & Turnos de doce horas, de día y de noche, durante los 92 días de junio a agosto de 2028. \\
Referencias & Art. 17.2.2 y 17.2.3, RT-20.06 y RT-07.12. \\
\end{tabla}
\begin{tabla}{I{0.269}I{0.731}}{Diccionario del entregable 1.6.5, Usuarios capacitados.}{tab:diccionario-entregable-1-6-5-usuarios}
  {\cab{1.6.5} & \cab{Usuarios capacitados}}
  {Fuente: elaboración propia a partir de la Guía del PMBOK (PMI, 2017, p. 162).}
Descripción del trabajo & Usuarios capacitados, material de apoyo y traspaso de conocimientos al equipo informático del cliente. Supone que el cliente libera a los usuarios en sus turnos para capacitarse. \\
Responsable & Benjamín Velásquez (LIG). \\
Hitos asociados & M10–M27. Aporta a H6 (mes 13) y H11 (mes 19). \\
Esfuerzo estimado & 1.728 HH, 1.296 externas (SRE 192, LIG 192, OPS 48, REL ext. 1.080 y RTC ext. 216). \\
1.6.5.1 & Usuarios capacitados por perfil y turno \\
Descripción del trabajo & Capacitar por ola y turno a cinco perfiles: usuarios finales, avanzados, administradores, equipo técnico y soporte, con certificado. \\
Entregable & Usuarios capacitados y certificados por perfil y turno. \\
Criterio de aceptación & Usuarios de cada ola capacitados en dos ciclos completos de turnos antes del inicio de su ola, con evaluación aprobada. Administradores y equipo técnico certificados antes del cierre de cada marcha blanca. \\
Responsable & Benjamín Velásquez (LIG). \\
Hitos asociados & M12–M14 y M18–M19. Aporta a H6 (mes 13), H7 (mes 16) y H11 (mes 19). \\
Esfuerzo estimado & 1.200 HH (LIG 120 y REL ext. 1.080). PERT baja. \\
Supuestos & Se capacitan 4.180 personas, dos horas por etapa, en grupos de unas 16, con hasta cuatro relatores en paralelo. \\
Referencias & RT-22.01, 22.04, 22.05 y Art. 90. T-12: RT-21.11; RT-22.01, 22.02, 22.04 y 22.05; TRA-003, 007 y 009. \\
1.6.5.2 & Manuales, guías y base de conocimiento \\
Descripción del trabajo & Redactar manuales por perfil, guías rápidas, preguntas frecuentes y videos en español, validados con usuarios. \\
Entregable & Manuales, guías rápidas y base de conocimiento editables en español. \\
Criterio de aceptación & Manuales, guías, preguntas frecuentes, videos y base de conocimiento publicados antes de la capacitación de cada perfil. Disponibilidad mensual de 99,5 \%. \\
Responsable & Benjamín Velásquez (LIG). \\
Hitos asociados & M10–M18. \\
Esfuerzo estimado & 288 HH (LIG 72 y RTC ext. 216). PERT baja. \\
Referencias & RT-22.02, 22.03 y Art. 90.3. T-12: RT-21.08 y 21.11; RT-22.01 a 22.03 y 22.08; RNF-TRV-21.11, TRA-009. \\
1.6.5.3 & Acta de transferencia tecnológica al equipo de TI \\
Descripción del trabajo & Asesorar durante seis meses a once personas del equipo informático del cliente, y certificar cada competencia. \\
Entregable & Entrega certificada de documentación, código, artefactos y procedimientos (Art. 77.1). \\
Criterio de aceptación & Acta firmada con las competencias transferidas y certificadas, cada una con un ejercicio práctico. \\
Responsable & Lucas Zamora (SRE). \\
Hitos asociados & M22–M27, los seis meses posteriores al paso a producción de la Etapa 2. \\
Esfuerzo estimado & 240 HH (SRE 192 y OPS 48). PERT media. \\
Referencias & RT-22.05, 22.07 y Art. 77.1. T-12: RT-21.04; RP-006 y 010. \\
\end{tabla}
\begin{tabla}{I{0.269}I{0.731}}{Diccionario del entregable 1.6.6, Plan de gestión del cambio e informes de adopción.}{tab:diccionario-entregable-1-6-6-plan}
  {\cab{1.6.6} & \cab{Plan de gestión del cambio e informes de adopción}}
  {Fuente: elaboración propia a partir de la Guía del PMBOK (PMI, 2017, p. 162).}
Descripción del trabajo & Gestión del cambio de la organización, desde el diagnóstico hasta medir la adopción. Supone que las superintendencias apoyan la designación de agentes de cambio. \\
Responsable & Benjamín Velásquez (LIG). \\
Hitos asociados & M2–M26. Sin hito propio. \\
Esfuerzo estimado & 824 HH, 228 externas (LF 96, LIG 324, ADT 56, OPS 60, TEC 60 y GCA ext. 228). \\
1.6.6.1 & Diagnóstico de impacto organizacional \\
Descripción del trabajo & Entrevistar, analizar el impacto y las resistencias por área y perfil, y priorizar las acciones. \\
Entregable & Diagnóstico del impacto por área y perfil, con las resistencias. \\
Criterio de aceptación & Diagnóstico aprobado con impacto y plan de acción por área. Cubre todos los perfiles de la Etapa 1 y, en mantenimiento, las dos medidas de la Etapa 2. \\
Responsable & Patricio Figueroa (LF). \\
Hitos asociados & M2–M4. \\
Esfuerzo estimado & 64 HH (LF 32 y GCA ext. 32). PERT media. \\
Referencias & Art. 89; Bases Técnicas, cap. 8, p. 14. \\
1.6.6.2 & Red de agentes de cambio habilitada \\
Descripción del trabajo & Designar con las superintendencias un agente de cambio por área y turno, capacitarlo y reunirlos cada mes. \\
Entregable & Red de agentes de cambio por área y turno, preparados. \\
Criterio de aceptación & Al menos un agente por área y por turno. Reunión mensual registrada. \\
Responsable & Patricio Figueroa (LF). \\
Hitos asociados & M5–M20. \\
Esfuerzo estimado & 128 HH (LF 64 y GCA externo 64). PERT baja. \\
Referencias & Art. 89. \\
1.6.6.3 & Procedimientos operativos rediseñados \\
Descripción del trabajo & Identificar los procedimientos que cambia SITCO Aranda, rediseñarlos con cada área y aprobarlos antes de su marcha blanca. \\
Entregable & Procedimientos operativos afectados por SITCO Aranda, rediseñados y aprobados. \\
Criterio de aceptación & Cada procedimiento afectado aprobado por su área antes de su marcha blanca, sin pasos nuevos para el operador, comprobado por el área. \\
Responsable & Benjamín Velásquez (LIG). \\
Hitos asociados & M8–M18. \\
Esfuerzo estimado & 264 HH (LIG 132 y GCA ext. 132). PERT media. \\
Referencias & Art. 89. \\
1.6.6.4 & Informes de adopción y acción correctiva \\
Descripción del trabajo & Medir cada mes la adopción por área y turno, sin sancionar a nadie, y seguir el plan de acción. \\
Entregable & Informes de uso, abandono de lo manual y satisfacción, con acciones correctivas. \\
Criterio de aceptación & Informe mensual con adopción por función, área, turno y persona. Meta: 90 \% de los cambios declarados por el propio operador. Un área bajo la meta recibe en dos semanas un plan de acción. \\
Responsable & Benjamín Velásquez (LIG). \\
Hitos asociados & M13–M26. \\
Esfuerzo estimado & 224 HH (LIG 168 y ADT 56). PERT baja. \\
Referencias & Art. 89. T-12: RNF-USA-005; TRA-005 y 010. \\
1.6.6.5 & Charlas de inicio de turno por ola \\
Descripción del trabajo & Preparar una charla breve por ola, con imágenes reales, y darla al inicio de cada turno durante dos semanas. \\
Entregable & Material de la charla de cada ola y registro de consultas. \\
Criterio de aceptación & Charla dada al inicio de los 28 turnos de las dos semanas de cada ola. Cada consulta registrada y respondida antes de la charla siguiente. \\
Responsable & Benjamín Velásquez (LIG). \\
Hitos asociados & M13, M14 y M19. \\
Esfuerzo estimado & 144 HH (LIG 24, OPS 60 y TEC 60). PERT baja. \\
Referencias & Subdocumento 6, 6.1.5.2 (Tabla \refext{T9}{tab:canales-control-comunicaciones} del Formulario T-9) y 6.1.5.3. \\
\end{tabla}

\apartado{Rama 1.7 Innovaciones}
La rama 1.7 tiene 5 entregables y 13 paquetes de trabajo, que suman 3.000 HH: las 2.920 HH de la cartera de innovaciones del Formulario T-19 y las 80 HH del sistema de gestión de la inteligencia artificial (1.7.2.4), que no ejecuta ninguna innovación. Las tablas \ref{tab:diccionario-entregable-1-7-1-control} a \ref{tab:diccionario-entregable-1-7-5-balance} presentan la ficha de cada entregable y, bajo ella, la de cada uno de sus paquetes.

\begin{tabla}{I{0.269}I{0.731}}{Diccionario del entregable 1.7.1, Control anticipado del arsénico en el lote comercial.}{tab:diccionario-entregable-1-7-1-control}
  {\cab{1.7.1} & \cab{Control anticipado del arsénico en el lote comercial}}
  {Fuente: elaboración propia a partir de la Guía del PMBOK (PMI, 2017, p. 162).}
Descripción del trabajo & Estimación del arsénico del lote comercial antes del embarque, para avisar si el lote supera el límite del comprador. Supone que la línea base de 2025 tiene 4 de 31 embarques penalizados. \\
Responsable & Benjamín Velásquez (LDA). \\
Hitos asociados & M8–M54. Aporta a H5 (mes 12). \\
Esfuerzo estimado & 464 HH, todas propias (LDA 192, DEV 240 y ADT 32). \\
1.7.1.1 & Control del arsénico implementado \\
Descripción del trabajo & Construir la estimación con la Etapa 1, integrarla en la formación de lotes y entregarla a pruebas. \\
Entregable & Componente de control del arsénico dentro de Trazabilidad y Embarque. \\
Criterio de aceptación & Alerta antes del embarque si el lote proyecta superar el límite de arsénico del comprador. Diferencia dentro del margen en todos los lotes de la marcha blanca y en nueve de cada diez lotes cada año. \\
Responsable & Benjamín Velásquez (LDA). \\
Hitos asociados & M8–M12. Aporta a H5 (mes 12). \\
Esfuerzo estimado & 400 HH (LDA 160 y DEV 240). PERT media. \\
Referencias & Subdocumento 13, 13.1.3. T-12: RT-26.02; RF-TRZ-020. \\
1.7.1.2 & Informe de verificación de la innovación 1 \\
Descripción del trabajo & Medir la diferencia por lote en la marcha blanca, emitir la línea base y medir cada embarque. \\
Entregable & Informe de embarques penalizados por arsénico, con línea base y meta. \\
Criterio de aceptación & Evaluaciones en los meses 30, 42 y 54, frente a la línea base de 2025. Meta de 2 embarques penalizados como máximo entre los meses 16 y 27, y de 1 cada doce meses desde el mes 28. \\
Responsable & Benjamín Velásquez (LDA). \\
Hitos asociados & M15, M30, M42 y M54. \\
Esfuerzo estimado & 64 HH (LDA 32 y ADT 32). PERT baja. \\
Referencias & Subdocumento 13, 13.1.4, 13.1.5 y RT-26.05, 26.08. \\
\end{tabla}
\begin{tabla}{I{0.269}I{0.731}}{Diccionario del entregable 1.7.2, Clasificación asistida de detenciones.}{tab:diccionario-entregable-1-7-2-clasificacion}
  {\cab{1.7.2} & \cab{Clasificación asistida de detenciones}}
  {Fuente: elaboración propia a partir de la Guía del PMBOK (PMI, 2017, p. 162).}
Descripción del trabajo & Propuesta automática de la causa de cada detención que las reglas no clasifican, para que el supervisor la confirme o corrija. Se activa cuando hay seis meses de clasificaciones confirmadas en cada área. \\
Responsable & Gabriel Sanzana (LD). \\
Hitos asociados & M3–M56. Aporta a H10 (mes 18). \\
Esfuerzo estimado & 1.464 HH, 552 externas (ESI 48, LD 320, DEV 264, ADT 280 y CDA ext. 552). \\
1.7.2.1 & Clasificador de detenciones implementado \\
Descripción del trabajo & Construir modelos de mina y planta que propongan la causa de cada detención, e integrarlos en la App SITCO Terreno. \\
Entregable & Clasificador de detenciones de mina y planta, con procedimiento de reentrenamiento documentado. \\
Criterio de aceptación & Modelos probados en QA en el mes 18. Precisión de 85 \% o más y falsos positivos de 15 \% o menos. Si la precisión del trimestre baja de 80 \%, se desactiva. \\
Responsable & Gabriel Sanzana (LD). \\
Hitos asociados & M13–M18. Aporta a H10 (mes 18). \\
Esfuerzo estimado & 880 HH (LD 88, DEV 264 y CDA ext. 528). PERT alta. \\
Referencias & Subdocumento 13, 13.2.2, 13.2.3, cap. 18 de las Bases técnicas transversales y RT-26.06; Art. 86; NIST (2023). T-12: RT-05.30; RT-18.01 a 18.09; RT-26.02 y 26.06; RNF-USA-009, RP-002. \\
1.7.2.2 & Informe de verificación de la innovación 2 \\
Descripción del trabajo & Medir la línea base, revisar una muestra de 5 \% y evaluar en los meses 27 y 33. \\
Entregable & Informe de los indicadores de la clasificación asistida. \\
Criterio de aceptación & Propuestas aceptadas sin cambios: 85 \% o más en mina (mes 27) y planta (mes 33). Precisión de 85 \% o más, con piso de 80 \%. Tiempo del supervisor 50 \% menor al mes 27. \\
Responsable & Gabriel Sanzana (LD). \\
Hitos asociados & M16–M33. \\
Esfuerzo estimado & 48 HH (LD 24 y CDA ext. 24). PERT baja. \\
Referencias & Subdocumento 13, 13.2.5, Tabla B.5 del Anexo B y RT-26.05; Bases Técnicas, 7.2. T-12: RT-05.30; RT-18.08, RP-002. \\
1.7.2.3 & Operación del clasificador \\
Descripción del trabajo & Revisar cada mes una muestra de 5 \% (mínimo 60 casos), vigilar que no empeore y reentrenarlo cada trimestre. \\
Entregable & Registro mensual de la muestra y versión reentrenada cada trimestre. \\
Criterio de aceptación & La Contraparte Técnica recibe cada versión con su precisión y tasa de falsos positivos. Precisión de 85 \% o más cada trimestre. Si la tasa de rechazo sube más de diez puntos, se reentrena. \\
Responsable & Gabriel Sanzana (LD). \\
Hitos asociados & M22–M56, con muestra mensual y reentrenamiento trimestral. \\
Esfuerzo estimado & 456 HH (LD 176 y ADT 280). PERT baja. \\
Supuestos & El equipo propio reentrena desde la activación, sin depender del científico de datos externo después del mes 33. \\
Referencias & Subdocumento 13, 13.2.3 y 13.2.4, y RT-18.01 a 18.09. T-12: RT-05.30; RT-26.06, RP-002. \\
1.7.2.4 & Sistema de gestión de la inteligencia artificial con alcance en el clasificador \\
Descripción del trabajo & Diagnosticar brechas de gestión de inteligencia artificial en el clasificador, según ISO/IEC 42001:2023, auditarla y preparar la certificación. \\
Entregable & Informe de brechas, auditoría interna y expediente de certificación. \\
Criterio de aceptación & Auditoría interna en el mes 9 y certificación en el mes 12. Expediente recibido por la Contraparte Técnica en el mes 12. Cada brecha con responsable y fecha de cierre. \\
Responsable & David Martínez (ESI). \\
Hitos asociados & M3, M9 y M12. \\
Esfuerzo estimado & 80 HH (ESI 48 y LD 32). PERT media. \\
Referencias & Subdocumento 1, 1.4 (Tabla \refext{S1}{tab:certificaciones-institucionales}), y cap. 18 de las Bases técnicas transversales. T-12: RT-18.04; RT-26.06. \\
\end{tabla}
\begin{tabla}{I{0.269}I{0.731}}{Diccionario del entregable 1.7.3, Factores de conciliación estimados con incertidumbre.}{tab:diccionario-entregable-1-7-3-factores}
  {\cab{1.7.3} & \cab{Factores de conciliación estimados con incertidumbre}}
  {Fuente: elaboración propia a partir de la Guía del PMBOK (PMI, 2017, p. 162).}
Descripción del trabajo & Factores de conciliación del pesaje de camiones con incertidumbre, estimados por un programa que detecta errores de medición graves y señala el instrumento sospechoso. Supone que metalurgia declara la incertidumbre de cada instrumento. \\
Responsable & Benjamín Rodríguez (AS). \\
Hitos asociados & M3–M27. Aporta a H5 (mes 12). \\
Esfuerzo estimado & 600 HH, todas propias (AS 144, LDA 24, DEV 288 y ADT 144). \\
1.7.3.1 & Modelo de incertidumbre de los instrumentos de pesaje \\
Descripción del trabajo & Declarar con metalurgia la incertidumbre de los tres instrumentos de pesaje y documentarla. \\
Entregable & Modelo de incertidumbre de los tres instrumentos de pesaje. \\
Criterio de aceptación & Metalurgia aprueba la incertidumbre de los tres instrumentos, con su origen documentado. La de la balanza de alimentación queda fija en el modelo. \\
Responsable & Benjamín Rodríguez (AS). \\
Hitos asociados & M3–M6. Aporta a H5 (mes 12). \\
Esfuerzo estimado & 72 HH (AS 24, LDA 24 y ADT 24). PERT media. \\
Referencias & Subdocumento 13, 13.3.3 y RT-26.04. \\
1.7.3.2 & Motor de reconciliación implementado \\
Descripción del trabajo & Construir el motor de reconciliación, integrarlo a Conciliación Metalúrgica y aplicarlo a doce meses históricos. \\
Entregable & Motor de reconciliación en la capa analítica. \\
Criterio de aceptación & Motor implementado en el mes 12 y aplicado a doce meses históricos. Funciona en paralelo con los factores documentados y no modifica la conciliación oficial. \\
Responsable & Benjamín Rodríguez (AS). \\
Hitos asociados & M7–M12. Aporta a H5 (mes 12). \\
Esfuerzo estimado & 480 HH (AS 96, DEV 288 y ADT 96). PERT media. \\
Supuestos & Los datos históricos llegan con campos personales enmascarados. Los datos en línea se usan desde el mes 10. \\
Referencias & Subdocumento 13, 13.3.3 y VDI 2048-1 (Verein Deutscher Ingenieure, 2017). T-12: RT-26.02. \\
1.7.3.3 & Informe de verificación de la innovación 3 \\
Descripción del trabajo & Comparar el motor con los factores documentados durante la marcha blanca y medir sus tres indicadores. \\
Entregable & Informe de los tres indicadores de la innovación 3. \\
Criterio de aceptación & Informe con mediciones al cierre de la marcha blanca y en los meses 21 y 27. Errores gruesos detectados en el turno siguiente. Cero factores sin base estadística. Sesgo por camión con intervalo de confianza de 95 \%. \\
Responsable & Benjamín Rodríguez (AS). \\
Hitos asociados & M15, M21, M27. \\
Esfuerzo estimado & 48 HH (AS 24 y ADT 24). PERT baja. \\
Referencias & Subdocumento 13, 13.3.5 y RT-26.05. T-12: RT-26.05 y 26.08. \\
\end{tabla}
\begin{tabla}{I{0.269}I{0.731}}{Diccionario del entregable 1.7.4, Servicio gestionado con traspaso programado.}{tab:diccionario-entregable-1-7-4-servicio}
  {\cab{1.7.4} & \cab{Servicio gestionado con traspaso programado}}
  {Fuente: elaboración propia a partir de la Guía del PMBOK (PMI, 2017, p. 162).}
Descripción del trabajo & Servicio con traspaso programado de funciones al equipo del cliente en tres bloques, con actas en los meses 27, 33, 39 y 45. Supone que el cliente designa administradores funcionales para cada bloque. \\
Responsable & Claudio Cabello (GS). \\
Hitos asociados & M21–M45, actas en M27, M33, M39 y M45. Sin hito propio. \\
Esfuerzo estimado & 152 HH, todas del GS. \\
1.7.4.1 & Plan de traspaso programado del servicio \\
Descripción del trabajo & Definir los tres bloques, diseñar los roles de administrador funcional y aprobar el plan. \\
Entregable & Plan con las funciones de cada bloque y su mes de traspaso. \\
Criterio de aceptación & La Contraparte Técnica aprueba el plan con los tres bloques y los meses 27, 33 y 39. El equipo del cliente resuelve al menos el 80 \% de las solicitudes funcionales en el mes 45. \\
Responsable & Claudio Cabello (GS). \\
Hitos asociados & M21–M27. \\
Esfuerzo estimado & 56 HH del GS. PERT media. \\
Referencias & Subdocumento 13, 13.4.3; Art. 30; Art. 77 y 77.1. T-12: RT-26.02, 26.05 y 26.07. \\
1.7.4.2 & Actas de traspaso de servicios al cliente \\
Descripción del trabajo & Operar cada bloque con sus administradores funcionales, verificar su autonomía y firmar el acta. \\
Entregable & Acta de traspaso de cada bloque al cliente. \\
Criterio de aceptación & Cuatro actas firmadas por el equipo informático del cliente en los meses 27, 33, 39 y 45. El acta del mes 45 verifica que el cliente resuelve 80 \% de las solicitudes funcionales. \\
Responsable & Claudio Cabello (GS). \\
Hitos asociados & M27, M33, M39 y M45. \\
Esfuerzo estimado & 96 HH del GS. PERT baja. \\
Referencias & Subdocumento 13, 13.4.5. \\
\end{tabla}
\begin{tabla}{I{0.269}I{0.731}}{Diccionario del entregable 1.7.5, Balance de dilución evitada.}{tab:diccionario-entregable-1-7-5-balance}
  {\cab{1.7.5} & \cab{Balance de dilución evitada}}
  {Fuente: elaboración propia a partir de la Guía del PMBOK (PMI, 2017, p. 162).}
Descripción del trabajo & Balance de la dilución evitada y de su efecto en energía y agua. Supone que el caudalímetro de planta de la Etapa 2 está instalado. \\
Responsable & Benjamín Velásquez (LDA). \\
Hitos asociados & M16–M46. Aporta a H10 (mes 18). \\
Esfuerzo estimado & 320 HH, todas propias (LDA 136, DEV 144 y ADT 40). \\
1.7.5.1 & Balance de dilución evitada implementado \\
Descripción del trabajo & Construir el cálculo en el módulo de Cumplimiento Ambiental, integrarlo con los tableros y entregarlo a pruebas. \\
Entregable & Cálculo del balance de dilución dentro de Cumplimiento Ambiental. \\
Criterio de aceptación & Balance implementado en el mes 18 y calculado cada mes desde el mes 22, con el caudalímetro instalado. Cobertura del balance de 90 \% o más. \\
Responsable & Benjamín Velásquez (LDA). \\
Hitos asociados & M16–M18. Aporta a H10 (mes 18). \\
Esfuerzo estimado & 240 HH (LDA 96 y DEV 144). PERT media. \\
Referencias & Subdocumento 13, 13.5.3. T-12: RT-26.02. \\
1.7.5.2 & Informe de verificación de la innovación 5 \\
Descripción del trabajo & Calcular la línea base de los meses 16 a 21, recalcularla en el mes 28 y evaluar cada semestre. \\
Entregable & Informe del indicador de dilución, con línea base y meta. \\
Criterio de aceptación & Informe con la línea base recalculada en el mes 28. Indicador de dilución: medio punto porcentual menos en el mes 34 y un punto menos en el mes 46, como beneficio esperado, no compromiso. \\
Responsable & Benjamín Velásquez (LDA). \\
Hitos asociados & M21, M28, M34, M40 y M46. \\
Esfuerzo estimado & 80 HH (LDA 40 y ADT 40). PERT baja. \\
Referencias & Subdocumento 13, 13.5.5; RT-26.05. \\
\end{tabla}

\apartado{Rama 1.8 Servicio de operación y mantención}
La rama 1.8 tiene 10 entregables y 29 paquetes de trabajo, que suman 122.672 HH. Los turnos 24x7 se cubren con cinco personas por puesto: un puesto suma 8.760 horas al año y una persona trabaja 1.824 (152 HH por 12 meses), ya descontados vacaciones y feriados, lo que da 4,8 personas. La operación comienza el 20 de septiembre de 2028, en el mes 21. Como es posterior a abril de 2028, también en faena rige el promedio de 40 horas semanales: el 7x7 de doce horas se mantiene y cada persona tiene un turno libre adicional cada seis semanas, que cubre la quinta persona del puesto (1.1.7.2). Los analistas de operación suman unas 18.200 HH al año: 17.520 de los dos puestos de la mesa y del NOC, 288 del refuerzo de las detenciones (1.8.1.1) y unas 415 de promedio anual de la mantención correctiva, las actualizaciones, las pruebas periódicas y el traspaso final (1.8.3.1, 1.8.4.1, 1.8.8 y 1.8.10.2). Equivalen a diez personas. La mesa de ayuda, el NOC y el SOC operan desde el centro de operaciones de Overfore en Santiago, y los técnicos trabajan en la faena en turnos 7x7 (Subdocumento 1, Anexo F del Subdocumento 1). Las tablas \ref{tab:diccionario-entregable-1-8-1-servicio} a \ref{tab:diccionario-entregable-1-8-10-acta} presentan la ficha de cada entregable y, bajo ella, la de cada uno de sus paquetes.

\begin{tabla}{I{0.269}I{0.731}}{Diccionario del entregable 1.8.1, Servicio de mesa de ayuda y soporte en faena.}{tab:diccionario-entregable-1-8-1-servicio}
  {\cab{1.8.1} & \cab{Servicio de mesa de ayuda y soporte en faena}}
  {Fuente: elaboración propia a partir de la Guía del PMBOK (PMI, 2017, p. 162).}
Descripción del trabajo & Mesa de ayuda que atiende todos los días las 24 horas, con soporte presencial en la faena. Supone el volumen de contactos con que se dimensiona la mesa. \\
Responsable & Lucas Zamora (SRE). \\
Hitos asociados & M21–M56. Sin hito propio. \\
Esfuerzo estimado & 53.688 HH, todas propias (SRE 120, OPS 27.144 y TEC 26.424). \\
1.8.1.1 & Mesa de ayuda dimensionada y operando \\
Descripción del trabajo & Dimensionar la mesa, definir sus turnos con refuerzo en las detenciones, y atender y registrar cada ticket hasta su cierre. \\
Entregable & Mesa de ayuda dimensionada y operando las 24 horas. \\
Criterio de aceptación & El informe mensual muestra 80 \% de llamadas atendidas en menos de 20 segundos, 70 \% de resolución en el primer contacto y abandono no superior a 5 \%. Cada ticket cerrado tiene encuesta de satisfacción. \\
Responsable & Lucas Zamora (SRE). \\
Hitos asociados & M21–M56, con refuerzo en las seis detenciones de planta. \\
Esfuerzo estimado & 27.144 HH del OPS. Sin PERT. \\
Supuestos & Se suponen 0,3 contactos por usuario y mes, para 4.180 usuarios permanentes, y un analista de refuerzo por turno en las detenciones. \\
Referencias & RT-21.05 a 21.14. T-12: RT-16.33; RT-21.05 a 21.15; RT-22.08, RNF-TRV-21.13 y 21.15; RP-006, 007 y 010. \\
1.8.1.2 & Soporte presencial en faena \\
Descripción del trabajo & Mantener un técnico en faena en cada turno y renovar cada año sus acreditaciones y exámenes. \\
Entregable & Registro de atención presencial en faena, con un técnico por turno. \\
Criterio de aceptación & Un técnico en faena en cada turno de día y de noche, sin acreditación ni examen de altura vencidos. En incidentes críticos con presencia, llega al puesto afectado en 60 minutos desde el escalamiento. \\
Responsable & Lucas Zamora (SRE). \\
Hitos asociados & M21–M56. \\
Esfuerzo estimado & 26.544 HH (SRE 120 y TEC 26.424). Sin PERT. \\
Supuestos & Cinco técnicos cubren el puesto en turnos 7x7. Las atenciones del puerto se dan a distancia. \\
Referencias & Bases Técnicas, RT-21.06 y RT-21.16 (p. 26), y Subdocumento 3, 3.2.2.6. T-12: RT-21.05, 21.06, 21.07 y 21.16; RP-007 y 008. \\
\end{tabla}
\begin{tabla}{I{0.269}I{0.731}}{Diccionario del entregable 1.8.2, Centro de operaciones de red y seguridad 24x7.}{tab:diccionario-entregable-1-8-2-centro}
  {\cab{1.8.2} & \cab{Centro de operaciones de red y seguridad 24x7}}
  {Fuente: elaboración propia a partir de la Guía del PMBOK (PMI, 2017, p. 162).}
Descripción del trabajo & Operación continua de la vigilancia de red y de seguridad, con gestión de la capacidad. Supone que la vigilancia cubre la nube y la faena. \\
Responsable & Lucas Zamora (SRE). \\
Hitos asociados & M21–M56. Sin hito propio. \\
Esfuerzo estimado & 53.160 HH, de las que 80 son externas (ESI 288, SRE 232, OPS 26.280, ASG 26.280 y RFE externo 80). \\
1.8.2.1 & Centro de operaciones de red 24x7 \\
Descripción del trabajo & Vigilar la nube, la faena, los gabinetes y los enlaces, y atender cada alerta según su severidad. \\
Entregable & Monitoreo continuo con alertas por efecto en el negocio. \\
Criterio de aceptación & La Contraparte Técnica verifica cada mes la disponibilidad por severidad: 99,9 \% crítica, 99,5 \% alta, 99,0 \% media y 98,0 \% baja. Incidentes críticos: respuesta en 15 minutos y resolución en 4 horas (Art. 78.2). Las ventanas aprobadas no cuentan. \\
Responsable & Lucas Zamora (SRE). \\
Hitos asociados & M21–M56. \\
Esfuerzo estimado & 26.280 HH del OPS. Sin PERT. \\
Supuestos & En el cambio de turno, un analista del NOC apoya la mesa de ayuda. \\
Referencias & RT-21.01 y cap. 14 de las Bases técnicas transversales. T-12: RT-06.19; RT-09.09, RT-10.01 y 10.02; RT-14.02, 14.04 y 14.09; RT-21.01 y 21.13; RNF-DIS-002, RNF-OPE-001, RP-006. \\
1.8.2.2 & Centro de operaciones de seguridad 24x7 \\
Descripción del trabajo & Vigilar eventos de seguridad 24x7, responder a ellos y continuar el escaneo de vulnerabilidades que 1.5.4.3 mantiene hasta el mes 20. \\
Entregable & Detección y respuesta a incidentes de seguridad 24x7. \\
Criterio de aceptación & Incidentes críticos comunicados al cliente en 2 horas, y brechas de seguridad o de datos personales en 24 horas. Vulnerabilidades críticas corregidas en 7 días, altas en 15 y medias en 30 días corridos (Art. 21.1). \\
Responsable & David Martínez (ESI). \\
Hitos asociados & M21–M56. \\
Esfuerzo estimado & 26.568 HH (ESI 288 y ASG 26.280). Sin PERT. \\
Referencias & RT-11.17, RT-11.18 y RT-11.19. T-12: RT-11.04 y 11.14 a 11.19; RT-21.01, RNF-SEG-011, RNF-TRV-11.04. \\
1.8.2.3 & Informes trimestrales de capacidad \\
Descripción del trabajo & Analizar la capacidad cada trimestre y proponer ajustes, y revisar la cobertura del rajo cada seis meses. \\
Entregable & Informe trimestral de capacidad y revisión semestral de cobertura. \\
Criterio de aceptación & Informe trimestral al Comité de Operación, con proyección a doce meses y tres años. Revisión en los meses 27, 33, 39, 45 y 51 marca zonas de sombra nuevas. \\
Responsable & Lucas Zamora (SRE). \\
Hitos asociados & M21–M56, trimestral. \\
Esfuerzo estimado & 312 HH (SRE 232 y RFE externo 80). PERT baja. \\
Referencias & RT-09.09 y Bases Técnicas, RT-03.24 (p. 25). T-12: RP-004. \\
\end{tabla}
\begin{tabla}{I{0.269}I{0.731}}{Diccionario del entregable 1.8.3, Servicio de mantención correctiva.}{tab:diccionario-entregable-1-8-3-servicio}
  {\cab{1.8.3} & \cab{Servicio de mantención correctiva}}
  {Fuente: elaboración propia a partir de la Guía del PMBOK (PMI, 2017, p. 162).}
Descripción del trabajo & Corrección de incidentes y análisis de su causa raíz, con plazos según la severidad del contrato. \\
Responsable & Lucas Zamora (SRE). \\
Hitos asociados & M21–M56. Sin hito propio. \\
Esfuerzo estimado & 3.072 HH, todas propias (SRE 768, DEV 1.728 y OPS 576). \\
1.8.3.1 & Registro de incidentes resueltos \\
Descripción del trabajo & Registrar cada incidente, diagnosticar su causa, corregirlo por las puertas de calidad y cerrarlo a tiempo. \\
Entregable & Registro de los incidentes y su resolución. \\
Criterio de aceptación & Incidentes resueltos dentro del plazo de su severidad, verificados por la Contraparte Técnica. Corrección de críticos dentro de 4 horas, según el Art. 78.2; los cambios menores definitivos se aplican en la ventana de 03:00 a 05:00 con aprobación previa. Ticket cerrado con confirmación del usuario o sin ella a los siete días. \\
Responsable & Lucas Zamora (SRE). \\
Hitos asociados & M21–M56. \\
Esfuerzo estimado & 2.880 HH (SRE 576, DEV 1.728 y OPS 576). Sin PERT. \\
Referencias & RT-21.15, 21.17. T-12: RT-21.03, 21.09, 21.15, 21.17 y 21.22. \\
1.8.3.2 & Informes de análisis de causa raíz \\
Descripción del trabajo & Reconstruir la secuencia de cada incidente crítico, identificar su causa, definir acciones correctivas y seguirlas hasta cerrarlas. \\
Entregable & Informe de causa raíz por incidente crítico y alto, con acciones correctivas. \\
Criterio de aceptación & Informe en cinco días hábiles desde un incidente crítico, y en diez desde uno de severidad alta. Overfore sigue las acciones hasta su cierre. \\
Responsable & Lucas Zamora (SRE). \\
Hitos asociados & M21–M56. \\
Esfuerzo estimado & 192 HH del SRE. PERT baja. \\
Referencias & RT-14.06. T-12: RT-11.19. \\
\end{tabla}
\begin{tabla}{I{0.269}I{0.731}}{Diccionario del entregable 1.8.4, Servicio de mantención preventiva y normativa.}{tab:diccionario-entregable-1-8-4-servicio}
  {\cab{1.8.4} & \cab{Servicio de mantención preventiva y normativa}}
  {Fuente: elaboración propia a partir de la Guía del PMBOK (PMI, 2017, p. 162).}
Descripción del trabajo & Mantención preventiva de los componentes y del equipamiento, y adecuación de la solución a las normas nuevas. Supone que los cambios se aplican en la ventana de 03:00 a 05:00. \\
Responsable & Lucas Zamora (SRE). \\
Hitos asociados & M21–M56. Sin hito propio. \\
Esfuerzo estimado & 696 HH, de las que 48 son externas (LD 48, SRE 192, DEV 144, OPS 120, TEC 144 y ABG externo 48). \\
1.8.4.1 & Plan anual de actualizaciones de base \\
Descripción del trabajo & Fijar con el cliente las detenciones de marzo y septiembre, planificar las versiones, probarlas y aplicarlas en ventana. \\
Entregable & Plan anual de actualización de componentes, con reversión y detenciones fechadas. \\
Criterio de aceptación & Plan aprobado cada año y aplicado sin incidentes críticos, en la ventana de 03:00 a 05:00. Aviso al cliente con diez días hábiles y plan de reversión. Prueba de carga antes de cada actualización mayor. \\
Responsable & Lucas Zamora (SRE). \\
Hitos asociados & M21–M56, anual. \\
Esfuerzo estimado & 240 HH (SRE 120 y OPS 120). PERT baja. \\
Referencias & RT-21.22. T-12: RT-10.05. \\
1.8.4.2 & Plan de ciclo de vida del equipamiento \\
Descripción del trabajo & Mantener, reponer y dar de baja los equipos de faena, y medir cada semestre la sala técnica. \\
Entregable & Plan de mantención, reposición y disposición del equipamiento de faena. \\
Criterio de aceptación & Plan vigente sin equipos fuera de soporte. Informe semestral de medición eléctrica y prueba contra incendios. Medios de salida borrados con certificado. Disposición conforme a la Ley N° 20.920. \\
Responsable & Lucas Zamora (SRE). \\
Hitos asociados & M21–M56. \\
Esfuerzo estimado & 216 HH (SRE 72 y TEC 144). PERT baja. \\
Referencias & RT-06.10, RT-08.16 a 08.18 y Subdocumento 4, 4.3.1.3 y 4.3.1.4. T-12: RT-06.10 y 06.18; RT-08.13 y 08.16 a 08.18. \\
1.8.4.3 & Adecuaciones normativas implementadas \\
Descripción del trabajo & Analizar cada norma nueva, construir el cambio en SITCO Aranda y desplegarlo antes de su fecha de exigencia. \\
Entregable & Cambios en la solución exigidos por nuevas regulaciones. \\
Criterio de aceptación & Cada cambio queda vigente antes de la fecha de exigencia de la norma, con el mismo estándar de calidad del desarrollo original. \\
Responsable & Gabriel Sanzana (LD). \\
Hitos asociados & M21–M56. \\
Esfuerzo estimado & 240 HH (LD 48, DEV 144 y ABG ext. 48). PERT alta. \\
Referencias & Subdocumento 3, 3.2.2.6. \\
\end{tabla}
\begin{tabla}{I{0.269}I{0.731}}{Diccionario del entregable 1.8.5, Servicio de mantención evolutiva.}{tab:diccionario-entregable-1-8-5-servicio}
  {\cab{1.8.5} & \cab{Servicio de mantención evolutiva}}
  {Fuente: elaboración propia a partir de la Guía del PMBOK (PMI, 2017, p. 162).}
Descripción del trabajo & Bolsa anual de horas para mejoras que prioriza el cliente, y control de la deuda técnica. Supone que el cliente prioriza la bolsa en el comité anual. \\
Responsable & Lucas Zamora (SRE). \\
Hitos asociados & M21–M56. Sin hito propio. \\
Esfuerzo estimado & 3.208 HH, todas propias (LD 728, SRE 144 y DEV 2.336). \\
1.8.5.1 & Plan anual de mantención evolutiva \\
Descripción del trabajo & Recibir solicitudes, estimar, priorizar con el cliente y construir; probar la replicación en el mes 30. \\
Entregable & Plan anual de mejoras priorizado por el cliente. \\
Criterio de aceptación & Plan aprobado cada año y ejecutado según la prioridad acordada. Bolsa de 960 HH al año. Informe de replicación: la nueva unidad opera sin cambios de código. \\
Responsable & Lucas Zamora (SRE). \\
Hitos asociados & M21–M56, con plan anual y avance mensual. \\
Esfuerzo estimado & 2.920 HH (LD 440, SRE 144 y DEV 2.336). Sin PERT. \\
Referencias & RT-21.19, 21.20, RT-02.12 y 13.2. \\
1.8.5.2 & Registro de deuda técnica \\
Descripción del trabajo & Medir la deuda técnica en cada versión, priorizarla con el Comité de Arquitectura y reducirla en el código nuevo. \\
Entregable & Registro de la deuda técnica y de su reducción. \\
Criterio de aceptación & Deuda técnica del código nuevo de 5 \% o menos en cada versión. Se destinan 8 HH al mes a reducirla. \\
Responsable & Gabriel Sanzana (LD). \\
Hitos asociados & M21–M56. \\
Esfuerzo estimado & 288 HH del LD. PERT baja. \\
Referencias & RT-21.21. \\
\end{tabla}
\begin{tabla}{I{0.269}I{0.731}}{Diccionario del entregable 1.8.6, Actualizaciones del plan de continuidad del negocio.}{tab:diccionario-entregable-1-8-6-actualizaciones}
  {\cab{1.8.6} & \cab{Actualizaciones del plan de continuidad del negocio}}
  {Fuente: elaboración propia a partir de la Guía del PMBOK (PMI, 2017, p. 162).}
Descripción del trabajo & Actualización anual del plan de continuidad del negocio, en los meses 27, 39 y 51. Supone que el cliente informa los cambios de sus procesos críticos. \\
Responsable & Lucas Zamora (SRE). \\
Hitos asociados & M21–M56, actualizaciones en M27, M39 y M51. Sin hito propio. \\
Esfuerzo estimado & 192 HH, 96 externas (SRE 96 y BCP ext. 96). \\
1.8.6.1 & Análisis de impacto en el negocio actualizado \\
Descripción del trabajo & Revisar cada año con el cliente los procesos críticos y actualizar sus tiempos de recuperación y de pérdida de datos. \\
Entregable & Análisis de impacto en el negocio revisado cada año. \\
Criterio de aceptación & Una actualización aprobada cada año, conforme con ISO 22301:2019. \\
Responsable & Lucas Zamora (SRE). \\
Hitos asociados & M21–M56, anual. \\
Esfuerzo estimado & 96 HH (SRE 48 y BCP ext. 48). PERT baja. \\
Referencias & RT-10.03. \\
1.8.6.2 & Plan de continuidad TIC actualizado \\
Descripción del trabajo & Actualizar cada año los procedimientos de recuperación y validarlos en el cambio semestral a la región secundaria. \\
Entregable & Plan de continuidad informática revisado cada año. \\
Criterio de aceptación & Una actualización aprobada cada año, conforme con ISO/IEC 27031:2011. \\
Responsable & Lucas Zamora (SRE). \\
Hitos asociados & M21–M56, anual. \\
Esfuerzo estimado & 96 HH (SRE 48 y BCP ext. 48). PERT baja. \\
Referencias & RT-10.04. \\
\end{tabla}
\begin{tabla}{I{0.269}I{0.731}}{Diccionario del entregable 1.8.7, Informes de servicio y cumplimiento.}{tab:diccionario-entregable-1-8-7-informes}
  {\cab{1.8.7} & \cab{Informes de servicio y cumplimiento}}
  {Fuente: elaboración propia a partir de la Guía del PMBOK (PMI, 2017, p. 162).}
Descripción del trabajo & Informes de niveles de servicio, certificaciones y sostenibilidad, incluida la huella de carbono por lote. Supone que ante diferencias de medición prevalece la del cliente, salvo error metodológico. \\
Responsable & Claudio Cabello (GS). \\
Hitos asociados & M21–M56. Sin hito propio. \\
Esfuerzo estimado & 5.976 HH, todas propias (AS 144, ESI 144, SRE 256, GS 5.000 y ADT 432). \\
1.8.7.1 & Informe mensual de niveles de servicio \\
Descripción del trabajo & Medir cada indicador por severidad, analizar sus incumplimientos y presentarlos al Comité de Operación. \\
Entregable & Informe mensual de niveles de servicio y plan de acción. \\
Criterio de aceptación & Informe recibido en los primeros cinco días hábiles de cada mes, con los siete indicadores del Art. 78.3 y su análisis por severidad. Incidentes críticos restaurados en 4 horas o menos; cambios fallidos de 5 \% o menos. \\
Responsable & Claudio Cabello (GS). \\
Hitos asociados & M21–M56, mensual. \\
Esfuerzo estimado & 576 HH del GS. PERT baja. \\
Referencias & RT-21.18 y Art. 78, 79; Bases Técnicas, cap. 18 y 7.2. T-12: RT-04.12; RT-10.01 y 10.09; RT-14.03, RT-21.13 y 21.18; RF-CON-005, RNF-DES-014, RNF-OPE-006, RNF-TRV-21.13. \\
1.8.7.2 & Informe anual de certificaciones vigentes \\
Descripción del trabajo & Reunir los certificados vigentes y preparar el expediente de registros de seguridad y ambiente para la recertificación. \\
Entregable & Informe anual de certificaciones vigentes y expediente de seguridad y ambiente. \\
Criterio de aceptación & Informe anual con copia de cada certificado vigente. Expediente antes del mes 31, con ISO 45001:2018 e ISO 14001:2015 vigentes y personal de faena con inducción vigente. \\
Responsable & Claudio Cabello (GS). \\
Hitos asociados & M21–M56, anual. \\
Esfuerzo estimado & 64 HH del GS. PERT baja. \\
Referencias & Art. 27; Bases Técnicas, 13.2, p. 22; RT-15.02, Bases Técnicas, cap. 15, p. 25; Subdocumento 1, Tabla \refext{S1}{tab:certificaciones-institucionales}. T-12: RT-15.07 y 15.09; RP-001. \\
1.8.7.3 & Huella de carbono por lote e informe anual de consumo y huella de la solución \\
Descripción del trabajo & Revisar la huella por lote cada mes, preparar el expediente trimestral y emitir el informe anual. \\
Entregable & Huellas de carbono por lote, expedientes trimestrales e informe anual. \\
Criterio de aceptación & Primer expediente recibido por la Contraparte Técnica en el mes 22, con los meses 16 a 21. Informe anual con alcances 1, 2 y 3, y factores de emisión verificados por un tercero. \\
Responsable & Lucas Zamora (SRE). \\
Hitos asociados & M21–M56; mensual, trimestral y anual. \\
Esfuerzo estimado & 544 HH (SRE 112 y ADT 432). PERT baja. \\
Supuestos & El cliente contrata la auditoría externa; el paquete prepara la evidencia. \\
Referencias & Bases Técnicas, 9.6, p. 19; 13.2, p. 22; 16.2; RT-05.23, RT-15.06. T-12: RT-15.03 y 15.06; RF-CUM-015 y 017; RNF-DES-013. \\
1.8.7.4 & Actas de los comités de la operación y plan de mejora continua \\
Descripción del trabajo & Convocar los comités Ejecutivo, de Operación y de Arquitectura, y mantener el plan de mejora continua. \\
Entregable & Actas mensuales de los comités y plan de mejora continua del servicio. \\
Criterio de aceptación & Actas entregadas dentro de dos días hábiles con acuerdos, responsables y plazos. Cada acción de mejora tiene dueño, fecha y efecto medido en los indicadores. \\
Responsable & Claudio Cabello (GS). \\
Hitos asociados & M21–M56, mensual. \\
Esfuerzo estimado & 864 HH (AS 144, ESI 144, SRE 144 y GS 432). Sin PERT. \\
Referencias & Art. 71, RT-19.04 y RT-19.09. \\
1.8.7.5 & Gestión del servicio y relación con el cliente \\
Descripción del trabajo & Atender como contraparte permanente al cliente, gestionar sus solicitudes, renovar licencias y controlar el contrato de servicio. \\
Entregable & Registro de solicitudes del cliente con su atención y seguimiento semanal. \\
Criterio de aceptación & Toda solicitud del cliente tiene respuesta en el plazo comprometido y queda registrada. Cada año, desde el mes 24, se acreditan las pólizas de seguro y las garantías del contrato. \\
Responsable & Claudio Cabello (GS). \\
Hitos asociados & M21–M56, desde la primera ola de la Etapa 2. \\
Esfuerzo estimado & 3.928 HH del GS, saldo de su dedicación. Sin PERT. \\
Supuestos & Gerente de Servicio dedicado desde el inicio de la Operación. \\
Referencias & RT-21.02 y Art. 78. T-12: RT-21.02 y 21.12. \\
\end{tabla}
\begin{tabla}{I{0.269}I{0.731}}{Diccionario del entregable 1.8.8, Informes de pruebas periódicas de la operación.}{tab:diccionario-entregable-1-8-8-informes}
  {\cab{1.8.8} & \cab{Informes de pruebas periódicas de la operación}}
  {Fuente: elaboración propia a partir de la Guía del PMBOK (PMI, 2017, p. 162).}
Descripción del trabajo & Pruebas repetidas durante la operación: resiliencia, intrusión, corte de 48 horas, carga, restauración de respaldos y madurez del proceso. Supone que los cambios de sitio se hacen en las detenciones de marzo y septiembre. \\
Responsable & Kavon Kermani (LCP). \\
Hitos asociados & M21–M56. Sin hito propio. \\
Esfuerzo estimado & 1.712 HH, 216 externas (ESI 48, LCP 368, SRE 352, QA 288, OPS 392, TEC 48, PEN ext. 144 y PRF ext. 72). \\
1.8.8.1 & Informes semestrales de resiliencia y recuperación \\
Descripción del trabajo & Programar cada semestre la prueba de resiliencia y el cambio real, y comparar tiempos con lo comprometido. \\
Entregable & Informes semestrales de resiliencia y de recuperación ante desastres. \\
Criterio de aceptación & Seis pruebas de resiliencia en los meses 24, 30, 36, 42, 48 y 54. Cinco cambios reales en los meses 27, 33, 39, 45 y 51, con recuperación en 4 horas y pérdida de 15 minutos o menos. \\
Responsable & Lucas Zamora (SRE). \\
Hitos asociados & M24 a M54, cada seis meses; las conmutaciones, en las detenciones de marzo y septiembre. \\
Esfuerzo estimado & 216 HH (SRE 136 y OPS 80). PERT baja. \\
Referencias & Subdocumento 9, 9.3.3, RT-07.07 y RT-10.07. T-12: RT-07.06 y 07.07; RNF-OPE-005. \\
1.8.8.2 & Informes anuales de intrusión y corte de 48 horas \\
Descripción del trabajo & Ejecutar con un auditor independiente la prueba de intrusión y el corte de 48 horas, y verificar los cierres. \\
Entregable & Informes anuales de intrusión y corte de 48 horas. \\
Criterio de aceptación & Tres intrusiones de tercero independiente y tres cortes de 48 horas, sin críticos ni altos abiertos. \\
Responsable & David Martínez (ESI). \\
Hitos asociados & M21–M56, anual. \\
Esfuerzo estimado & 480 HH (ESI 48, SRE 144, OPS 96, TEC 48 y PEN ext. 144). PERT media. \\
Supuestos & El auditor independiente se contrata por campaña. \\
Referencias & Subdocumento 9, 9.3.3 y RT-11.20. T-12: RT-11.20 y 11.21. \\
1.8.8.3 & Informes de carga por versión mayor \\
Descripción del trabajo & Ejecutar la prueba de carga en Preproducción y comparar con la línea base. \\
Entregable & Informe de prueba de carga por versión mayor y actualización anual. \\
Criterio de aceptación & Al menos tres pruebas en la operación con carga de 1,5 veces el pico declarado, hasta el punto de quiebre. Resistencia de 24 horas. La degradación no supera 10 \% respecto de la versión anterior. \\
Responsable & Kavon Kermani (LCP). \\
Hitos asociados & M21–M56, por versión mayor. \\
Esfuerzo estimado & 120 HH (LCP 48 y PRF ext. 72). PERT baja. \\
Referencias & Subdocumento 9, 9.3.3, RT-09.06 y Formulario T-13. \\
1.8.8.4 & Puertas de calidad de las versiones de la operación \\
Descripción del trabajo & Verificar que cada versión pasa las puertas de calidad, revisar la regresión automatizada y registrar el resultado. \\
Entregable & Registro de calidad de cada versión liberada a la operación. \\
Criterio de aceptación & Ninguna versión llega a producción sin las puertas de calidad aprobadas y la regresión de ambas etapas sin fallas. \\
Responsable & Kavon Kermani (LCP). \\
Hitos asociados & M21–M56, por versión. \\
Esfuerzo estimado & 576 HH (LCP 288 y QA 288). PERT baja. \\
Supuestos & Se estima una versión por mes. La regresión es automatizada; las horas revisan su resultado. \\
Referencias & Subdocumento 9, 9.2 y RT-21.20. \\
1.8.8.5 & Pruebas mensuales de restauración de respaldos \\
Descripción del trabajo & Restaurar cada mes una muestra de los respaldos en un ambiente aislado y medir el tiempo que toma. \\
Entregable & Registro mensual de restauración de respaldos, con tiempo medido y fallas. \\
Criterio de aceptación & 36 pruebas mensuales con restauración en 4 horas o menos. Respaldos cifrados y copias inmutables verificadas en cada prueba. \\
Responsable & Lucas Zamora (SRE). \\
Hitos asociados & M21–M56, mensual. \\
Esfuerzo estimado & 288 HH (SRE 72 y OPS 216). PERT baja. \\
Referencias & RT-07.10, RT-07.11 y RT-07.12. T-12: RT-07.10 a 07.12 y 07.14; RNF-OPE-003 y 004. \\
1.8.8.6 & Reevaluaciones anuales OWASP SAMM \\
Descripción del trabajo & Repetir cada año la evaluación de las quince prácticas de madurez del desarrollo, comparar con la anterior e informar. \\
Entregable & Informes de madurez en los meses 33 y 45 con sus brechas. \\
Criterio de aceptación & Dos reevaluaciones en los meses 33 y 45; con la del mes 21, una cada doce meses. Ninguna práctica baja del nivel 3 sin plan de recuperación. \\
Responsable & Kavon Kermani (LCP). \\
Hitos asociados & M33 y M45. \\
Esfuerzo estimado & 32 HH del LCP. PERT baja. \\
Referencias & Subdocumento 9, 9.1.3. T-12: RT-11.28. \\
\end{tabla}
\begin{tabla}{I{0.269}I{0.731}}{Diccionario del entregable 1.8.9, Personal del cliente capacitado en la operación.}{tab:diccionario-entregable-1-8-9-personal}
  {\cab{1.8.9} & \cab{Personal del cliente capacitado en la operación}}
  {Fuente: elaboración propia a partir de la Guía del PMBOK (PMI, 2017, p. 162).}
Descripción del trabajo & Capacitación del personal del cliente durante la operación: jornadas semestrales e inducción de personal nuevo. Supone la rotación normal del personal de la faena. \\
Responsable & Benjamín Velásquez (LIG). \\
Hitos asociados & M21–M56, semestral. Sin hito propio. \\
Esfuerzo estimado & 288 HH, 96 externas (LIG 192 y REL ext. 96). \\
1.8.9.1 & Personal actualizado en jornadas semestrales \\
Descripción del trabajo & Preparar y dictar cada jornada con las versiones del año, en grupos de 16 personas y por perfil. \\
Entregable & Dos jornadas de actualización por año. \\
Criterio de aceptación & Dos jornadas por año con registro de asistencia y evaluación certificada por perfil y por rol (Art. 77.1). Contenido actualizado con las versiones del año. \\
Responsable & Benjamín Velásquez (LIG). \\
Hitos asociados & M23 a M53, cada seis meses, en noviembre y mayo. \\
Esfuerzo estimado & 144 HH (LIG 48 y REL ext. 96). PERT baja. \\
Referencias & RT-22.06. T-12: RT-22.02 y 22.06; RP-010. \\
1.8.9.2 & Personal nuevo del cliente inducido \\
Descripción del trabajo & Inducir a cada persona nueva del cliente antes de que opere SITCO, y registrar la inducción. \\
Entregable & Inducción de cada persona nueva del cliente, registrada. \\
Criterio de aceptación & Personal nuevo inducido antes de operar el sistema. Mismo material de 1.6.5.2 actualizado. \\
Responsable & Benjamín Velásquez (LIG). \\
Hitos asociados & M21–M56, continuo. \\
Esfuerzo estimado & 144 HH del LIG. PERT baja. \\
Referencias & RT-22.06. T-12: RP-010. \\
\end{tabla}
\begin{tabla}{I{0.269}I{0.731}}{Diccionario del entregable 1.8.10, Acta de devolución del servicio.}{tab:diccionario-entregable-1-8-10-acta}
  {\cab{1.8.10} & \cab{Acta de devolución del servicio}}
  {Fuente: elaboración propia a partir de la Guía del PMBOK (PMI, 2017, p. 162).}
Descripción del trabajo & Devolución del servicio al cliente al término del contrato, con acta de cierre de los entregables recibidos y pendientes. Supone que el cliente designa al nuevo operador tres meses antes del término. \\
Responsable & Claudio Cabello (GS). \\
Hitos asociados & M54–M56. Sin hito propio. \\
Esfuerzo estimado & 680 HH, todas propias (SRE 252, GS 208, DEV 64 y OPS 156). \\
1.8.10.1 & Código, datos y credenciales entregados \\
Descripción del trabajo & Ejecutar el inventario del plan, entregar y verificar con el cliente. \\
Entregable & Código, documentación, datos y credenciales, en formatos abiertos (Art. 77.1). \\
Criterio de aceptación & Acta de entrega con el inventario completo verificado. Formatos abiertos, con diccionario de datos y un ejemplo de cada archivo entregado. \\
Responsable & Lucas Zamora (SRE). \\
Hitos asociados & M54–M56. \\
Esfuerzo estimado & 160 HH (SRE 96 y DEV 64). PERT media. \\
Referencias & Art. 77.1, 82 y RT-05.06. \\
1.8.10.2 & Informe de acompañamiento de 90 días \\
Descripción del trabajo & Acompañar durante 90 días al nuevo operador en la mesa de ayuda, el NOC y los despliegues. \\
Entregable & Acompañamiento de 90 días al nuevo operador. \\
Criterio de aceptación & Informe de cierre del acompañamiento firmado por el cliente. Sin interrupción del servicio durante el traspaso. \\
Responsable & Claudio Cabello (GS). \\
Hitos asociados & M54–M56. \\
Esfuerzo estimado & 520 HH (SRE 156, GS 208 y OPS 156). Sin PERT. \\
Referencias & Art. 82. \\
\end{tabla}

\seccionsinnumero{Referencias}
\begin{referencias}
\item AsyncAPI Initiative. (2023). AsyncAPI Specification, versión 3.0.0.
\item Compañía Minera Altos de Aranda. (2026a). \emph{Bases administrativas}. Licitación N° TFEP-01/2026.
\item Compañía Minera Altos de Aranda. (2026b). \emph{Bases técnicas del caso Minería}. Licitación N° TFEP-01/2026.
\item Compañía Minera Altos de Aranda. (2026c). \emph{Bases técnicas transversales}. Licitación N° TFEP-01/2026.
\item Decreto con Fuerza de Ley N° 1. (2003, 16 de enero). Fija el texto refundido, coordinado y sistematizado del Código del Trabajo. Diario Oficial de la República de Chile.
\item Gans, N., Koole, G. y Mandelbaum, A. (2003). Telephone call centers: Tutorial, review, and research prospects. Manufacturing \& Service Operations Management, 5(2), 79–141. https://doi.org/10.1287/msom.5.2.79.16071
\item International Electrotechnical Commission. (2013). Enterprise-control system integration. Part 1: Models and terminology (IEC 62264-1:2013).
\item International Electrotechnical Commission. (2020). Security for industrial automation and control systems. Part 3-2: Security risk assessment for system design (IEC 62443-3-2:2020).
\item International Organization for Standardization e International Electrotechnical Commission. (2008). Software engineering. Software product Quality Requirements and Evaluation (SQuaRE). Data quality model (ISO/IEC 25012:2008).
\item International Organization for Standardization e International Electrotechnical Commission. (2011). Information technology. Security techniques. Guidelines for information and communication technology readiness for business continuity (ISO/IEC 27031:2011).
\item International Organization for Standardization e International Electrotechnical Commission. (2015). Information technology. Security techniques. Code of practice for information security controls based on ISO/IEC 27002 for cloud services (ISO/IEC 27017:2015).
\item International Organization for Standardization e International Electrotechnical Commission. (2019). Information technology. Security techniques. Code of practice for protection of personally identifiable information (PII) in public clouds acting as PII processors (ISO/IEC 27018:2019).
\item International Organization for Standardization e International Electrotechnical Commission. (2022a). Information security, cybersecurity and privacy protection. Information security management systems. Requirements (ISO/IEC 27001:2022).
\item International Organization for Standardization e International Electrotechnical Commission. (2022b). Information security, cybersecurity and privacy protection. Information security controls (ISO/IEC 27002:2022).
\item International Organization for Standardization e International Electrotechnical Commission. (2023a). Systems and software engineering. Systems and software Quality Requirements and Evaluation (SQuaRE). Product quality model (ISO/IEC 25010:2023).
\item International Organization for Standardization e International Electrotechnical Commission. (2023b). Information technology. Artificial intelligence. Management system (ISO/IEC 42001:2023).
\item International Organization for Standardization, International Electrotechnical Commission e Institute of Electrical and Electronics Engineers. (2018). Systems and software engineering. Life cycle processes. Requirements engineering (ISO/IEC/IEEE 29148:2018).
\item International Organization for Standardization, International Electrotechnical Commission e Institute of Electrical and Electronics Engineers. (2022). Software, systems and enterprise. Architecture description (ISO/IEC/IEEE 42010:2022).
\item International Organization for Standardization. (2015). Environmental management systems. Requirements with guidance for use (ISO 14001:2015).
\item International Organization for Standardization. (2018a). Occupational health and safety management systems. Requirements with guidance for use (ISO 45001:2018).
\item International Organization for Standardization. (2018b). Risk management. Guidelines (ISO 31000:2018).
\item International Organization for Standardization. (2019). Security and resilience. Business continuity management systems. Requirements (ISO 22301:2019).
\item International Organization for Standardization. (2020). Earth-moving machinery and mobile road construction machinery. Worksite data exchange. Part 3: Telematics data (ISO 15143-3:2020).
\item Ke, G., Meng, Q., Finley, T., Wang, T., Chen, W., Ma, W., Ye, Q. y Liu, T.-Y. (2017). LightGBM: A highly efficient gradient boosting decision tree. Advances in Neural Information Processing Systems, 30, 3146–3154.
\item Ley N° 19.799. (2002, 12 de abril). Sobre documentos electrónicos, firma electrónica y servicios de certificación de dicha firma. Diario Oficial de la República de Chile.
\item Ley N° 20.920. (2016, 1 de junio). Establece marco para la gestión de residuos, la responsabilidad extendida del productor y fomento al reciclaje. Diario Oficial de la República de Chile.
\item Ley N° 21.561. (2023, 26 de abril). Modifica el Código del Trabajo con el objeto de reducir la jornada laboral. Diario Oficial de la República de Chile.
\item Ley N° 21.719. (2024, 13 de diciembre). Regula la protección y el tratamiento de los datos personales y crea la Agencia de Protección de Datos Personales. Diario Oficial de la República de Chile.
\item Linux Foundation. (2022). The Software Package Data Exchange (SPDX) Specification, versión 2.3.
\item McConnell, S. (2006). Software estimation: Demystifying the black art. Microsoft Press.
\item Ministerio de Minería. (2004). \emph{Decreto Supremo N.º 132, Reglamento de Seguridad Minera}. Diario Oficial de la República de Chile, 7 de febrero de 2004.
\item National Institute of Standards and Technology. (2023). Artificial intelligence risk management framework (AI RMF 1.0) (NIST AI 100-1). https://doi.org/10.6028/NIST.AI.100-1
\item Open Source Security Foundation. (2023). Supply-chain Levels for Software Artifacts (SLSA), versión 1.0.
\item OpenAPI Initiative. (2021). OpenAPI Specification, versión 3.1.0.
\item OpenTelemetry Authors. (2024). OpenTelemetry Specification.
\item Overfore S.A. (2026a). Subdocumento 1. Licitación N° TFEP-01/2026.
\item Overfore S.A. (2026b). Subdocumento 1: Anexos. Licitación N° TFEP-01/2026.
\item OWASP Foundation. (2020). Software Assurance Maturity Model (SAMM), versión 2.0.
\item OWASP Foundation. (2021). Application Security Verification Standard, versión 4.0.3.
\item OWASP Foundation. (2024). CycloneDX Bill of Materials Specification (ECMA-424), versión 1.6.
\item Project Management Institute. (2017). \emph{Guía de los fundamentos para la dirección de proyectos (Guía del PMBOK) }(6.ª ed.).
\item Rose, S., Borchert, O., Mitchell, S. y Connelly, S. (2020). Zero trust architecture (NIST Special Publication 800-207). National Institute of Standards and Technology.
\item Verein Deutscher Ingenieure. (2017). Control and quality improvement of process data and their uncertainties by means of correction calculation for operation and acceptance tests (VDI 2048 Part 1). Beuth Verlag.
\item World Business Council for Sustainable Development y World Resources Institute. (2004). The Greenhouse Gas Protocol: A corporate accounting and reporting standard (ed. rev.).
\item World Wide Web Consortium. (2023). Web Content Accessibility Guidelines (WCAG) 2.2.
\end{referencias}
\end{document}
